\documentclass[11pt]{amsart}
\usepackage[margin=27mm]{geometry}
\usepackage{amsmath,amssymb}
\usepackage{tikz}

% amsart loads amsthm, whose proof environment and \qed clash with pf2's.
% We use pf2's proof environment throughout, as Lamport's format requires.
\let\proof\relax
\let\endproof\relax
\let\qed\relax
\usepackage{pf2}

\usepackage[colorlinks=true,linkcolor=blue,citecolor=blue,urlcolor=blue]{hyperref}

% The author's macros.
%
% some calligraphy
%
\newcommand\cA{{\mathcal A}}
\newcommand\cB{{\mathcal B}}
\newcommand\cC{{\mathcal C}}
\newcommand\cD{{\mathcal D}}
\newcommand\cE{{\mathcal E}}
\newcommand\cF{{\mathcal F}}
\newcommand\cG{{\mathcal G}}
\newcommand\cH{{\mathcal H}}
\newcommand\cI{{\mathcal I}}
\newcommand\cJ{{\mathcal J}}
\newcommand\cK{{\mathcal K}}
\newcommand\cL{{\mathcal L}}
\newcommand\cM{{\mathcal M}}
\newcommand\cN{{\mathcal N}}
\newcommand\cP{{\mathcal P}}
\newcommand\cQ{{\mathcal Q}}
\newcommand\cR{{\mathcal R}}
\newcommand\cS{{\mathcal S}}
\newcommand\cT{{\mathcal T}}
\newcommand\cU{{\mathcal U}}
\newcommand\cV{{\mathcal V}}

%
% fields and rings (and a semigroup)
%
\newcommand\cx{{\mathbb C}}% complexes
\newcommand\fld{{\mathbb F}}
\newcommand\flde{{\mathbb E}}
\newcommand\ints{{\mathbb Z}}
\newcommand\nn{{\mathbb N}}%non-negative integers
\newcommand\re{{\mathbb R}}%reals
\newcommand\rats{{\mathbb Q}}

\newcommand\Zu{{\mathbf u}}
\newcommand\Zv{{\mathbf v}}
\newcommand\Zw{{\mathbf w}}
\newcommand\Zy{{\mathbf y}}
\newcommand\Zz{{\mathbf z}}
\newcommand\Zx{{\mathbf x}}
\newcommand\Zj{{\mathbf j}}


%
% the really useful stuff
%
%\newcommand\opk[1]{\mathop{\hbox{\rm #1}}\nolimits} %% old version
\newcommand\opk[1]{\mathop{\mathrm{#1}}\nolimits}
\newcommand\comp[1]{{\mkern2mu\overline{\mkern-2mu#1}}}
\newcommand\diff{\mathbin{\mkern-1.5mu\setminus\mkern-1.5mu}}% for \setminus
\newcommand\gbin[2]{\genfrac{[}{]}{0pt}{}{#1}{#2}}
\newcommand\inv{^{-1}}
\newcommand\res{\mathbin{\mkern-2.0mu\restriction\mkern-2.0mu}}% ch2
\newcommand\sbs{\subseteq}
\newcommand\sps{\supseteq}
\newcommand\seq[3]{#1_{#2},\ldots,#1_{#3}}
\newcommand\supp{\opk{supp}}
\newcommand\im{\opk{im}}
\newcommand\row{\opk{row}}
\newcommand\pmat[1]{\begin{pmatrix} #1 \end{pmatrix}}
\newcommand\cprod{\mathbin{\square}}
%
% matrix theory
%
\newcommand\ip[2]{\langle#1,#2\rangle}
\newcommand\one{{\bf1}}
\newcommand\rk{\opk{rk}}
\newcommand\tr{\opk{tr}}
\newcommand\col{\opk{col}}
\newcommand\mat[3]{\mathrm{Mat}_{#1\times #2}(#3)}

%
% some group theory
%
\newcommand\aut[1]{{\rm Aut}(#1)}
\newcommand\fx[1]{{\rm fix}(#1)}% ch2
\newcommand\grp[1]{\langle #1\rangle}
\newcommand\nrml{\vartriangleleft}
\newcommand\nrmleq{\trianglelefteq}
\newcommand\alt[1]{\opk{Alt}(#1)}
\DeclareMathOperator{\Sym}{Sym}
% End of the author's macros.
% Macros of the blueprint. The author's own macros are at the start of each document; do not duplicate them here.

% The macros of leanblueprint, as no-ops under pdflatex, so that its dependency graph can be
% generated later: \lean{Decl.name}, \uses{label, label}, \leanok.
\newcommand{\lean}[1]{}
\newcommand{\uses}[1]{}
\newcommand{\leanok}{}

% Results that carry the paper's own number, for instance
%   \begin{cgss}{Lemma}{2.8}\label{lem:2.8} ... \end{cgss}
% and results of ours, numbered within the section, through the usual counters.
% Statements are set upright, since pf2's keywords (Assume, Prove) are small capitals, which have
% no italic shape in Computer Modern.
\theoremstyle{definition}
\makeatletter
\newtheorem*{cgss@env}{\cgss@name}
\newcommand{\cgss@name}{}
% The environment sets an anchor and makes the paper's number the text of \ref, so that
% \ref{lem:2.9} prints 2.9 and links to the lemma.
\newenvironment{cgss}[2]{\renewcommand{\cgss@name}{#1 #2}\begin{cgss@env}%
  \phantomsection\def\@currentlabel{#2}}{\end{cgss@env}}
\makeatother
\newtheorem{lemma}{Lemma}[section]
\newtheorem{proposition}[lemma]{Proposition}
\newtheorem{corollary}[lemma]{Corollary}
\newtheorem{definition}[lemma]{Definition}
\newtheorem*{hypothesis}{Hypothesis}
\theoremstyle{remark}
\newtheorem*{remark}{Remark}

% Citing assumptions as in [L95]: \asm{1} gives <1>, the assumption of the enclosing level-1
% step, and \asm[2]{0} gives <0>:2, the second assumption of the theorem itself. Steps are
% cited with pf2's \stepref{label}.
\newcommand{\asm}[2][]{$\langle#2\rangle$\if\relax\detokenize{#1}\relax\else{:}#1\fi}

% Tags of the structured proofs: a step that departs from the paper, and a step that fills a gap.
\newcommand{\new}{\textbf{[new]}}
\newcommand{\repair}{\textbf{[repair]}}

% The paper's sets of Hypothesis 2.4.
\newcommand{\Ga}{\Gamma_a}
\newcommand{\Gb}{\Gamma_b}
\newcommand{\Gc}{\Gamma_c}

% Vectors that the author's macros do not provide, in its pattern: the vectors a, b, c of Hypothesis 2.4,
% a unit vector e, and roots r, s, t.
\newcommand\Za{{\mathbf a}}
\newcommand\Zb{{\mathbf b}}
\newcommand\Zc{{\mathbf c}}
\newcommand\Ze{{\mathbf e}}
\newcommand\Zr{{\mathbf r}}
\newcommand\Zs{{\mathbf s}}
\newcommand\Zt{{\mathbf t}}

% Additional vector names for the arbitrary adjacent pair in Lemma 2.10.
\newcommand\Zp{{\mathbf p}}
\newcommand\Zq{{\mathbf q}}

% Upright Dynkin type letters, following cgss/notation.md: \typ{A}_\ell, \typ{E}_8.
\newcommand{\typ}[1]{\mathrm{#1}}

% Auxiliary results have descriptive headings, not numbers that could clash with CGSS.
% For example, \begin{namedresult}{Corollary}{of Lemma 2.8} inherits cgss's unique anchor;
% \ref prints the descriptive second argument. Existing result environments are unchanged.
\newenvironment{namedresult}[2]{\begin{cgss}{#1}{#2}}{\end{cgss}}

% The all-positive half root used in the Section 3 spanning and inclusion proofs.
\newcommand\Zh{{\mathbf h}}

% A visible citation with its exact Lean name as metadata and a pdfLaTeX-safe display.
% Unlike \path, the display may render Unicode subscripts and Greek letters as TeX math.
% Example: \leancite{LineSystem.E₈}{\texttt{LineSystem.E\ensuremath{{}_8}}}.
\newcommand{\leancite}[2]{#2}

% Local layout for the merged structured proofs; existing section settings are restored.
\newenvironment{mergedproofsection}{\begingroup\emergencystretch=3em}{\par\endgroup}

% A descriptive library-fact label (the paragraph already supplies its hyperlink anchor).
\makeatletter
\newcommand{\factlabel}[2]{\def\@currentlabel{#1}\label{#2}}
\makeatother



\usepackage{microtype}
\theoremstyle{definition}
\newtheorem{theorem}[lemma]{Theorem}
\DeclareMathOperator{\Spec}{Spec}
\DeclareMathOperator{\mult}{mult}
\DeclareMathOperator{\diag}{diag}
\DeclareMathOperator{\spanop}{span}
\newcommand{\ms}[1]{\left[\!\left[#1\right]\!\right]}
\newcommand{\norm}[1]{\left\lVert#1\right\rVert}
\newcommand{\code}[1]{\ifmmode\text{\small\nolinkurl{#1}}\else{\small\nolinkurl{#1}}\fi}
\makeatletter
\g@addto@macro\UrlBreaks{\do\a\do\b\do\c\do\d\do\e\do\f\do\g\do\h\do\i\do\j\do\k\do\l\do\m\do\n\do\o\do\p\do\q\do\r\do\s\do\t\do\u\do\v\do\w\do\x\do\y\do\z\do\A\do\B\do\C\do\D\do\E\do\F\do\G\do\H\do\I\do\J\do\K\do\L\do\M\do\N\do\O\do\P\do\Q\do\R\do\S\do\T\do\U\do\V\do\W\do\X\do\Y\do\Z}
\makeatother
\setlength{\emergencystretch}{3em}
\raggedbottom
\begin{document}
\title{A spectral toolkit for CGSS: structured proofs}
\date{28 September 2026}
\begin{abstract}
We give structured proofs of the spectral toolkit for Cameron, Goethals,
Seidel and Shult. Every proof follows the argument of the ordinary version,
with explicit hypotheses, eigenvalue multiplicities and finite-dimensional
coordinate maps. We retain its proposed Lean names and definitions; these
proof trees guided the formalization.
\end{abstract}
\maketitle
\tableofcontents
\section{Conventions and checked facts}
We follow the conventions of the blueprint (\path{blueprint/conventions.tex})
and the proof format of Lamport~\cite{L95}. The statements use numbered assumptions and a conclusion; the
level-one steps follow the ordinary proofs in \path{spectraltools.tex}.
A proof leaf names its inputs. An assumption at level $k$ is cited as
$\langle k\rangle$, and its numbered part $i$ as $\langle k\rangle{:}i$.
The tags \verb|\lean| and \verb|\uses| are metadata and print nothing.
The ordinary source-status lines distinguish existing project declarations
from proposed interfaces; names alone do not assert a Lean proof. The
ordinary version already contains two corrections, the hypothesis $A\succeq0$ for
principal minors in Theorem~\ref{thm:psd} and the symmetry guard for
graph certificates in Theorem~\ref{thm:checker}, so they are not repairs here. Where the ordinary
version attributes a result to Mathlib, the corresponding leaf applies
that declaration. No step departs from the argument of the ordinary
version, so no step is tagged \new\ or \repair. The proof of interlacing
uses the ordinary proof's dimension witnesses directly; it does not need a
maximum over a space of subspaces.

All sets indexing a matrix are finite, and we will choose decidable
equality and decidable adjacency classically. Entries use functional
notation $A(u,v)$ and $\Zv(u)$. Vectors are columns, $I$ is the identity,
$J$ is the all-ones matrix and $\Zj$ is the all-ones vector, on the
indicated index set. In $\re^U$, the inner product is
$\ip{\Zx}{\Zy}=\sum_u\Zx(u)\Zy(u)$. Matrix multiplication, trace,
transpose and polynomial evaluation have their usual coordinate definitions.
Rank means dimension of the image. Empty sums are zero and empty products
are one, including the determinant of a zero-order matrix. Connectedness
includes nonemptiness. Graphs are simple; we repeat finiteness in every
result. All spectral multisets include repetitions. Dimension differences
in formulas for multiplicities are integer differences proved nonnegative,
not unexplained truncated subtraction.

The following register specifies the exact library facts used in leaves.
Each path is relative to \path{Mathlib/}, at commit \path{082e2d37}
(4 September 2026), the version pinned in \path{lake-manifest.json}. Elementary coordinate expansions, induction, finite-set bijections
and field arithmetic in the proofs are deductions written out below.
\begingroup\makeatletter\def\@currentlabel{1}\subsection*{Facts 1: spec}\phantomsection\label{facts:spec}\makeatother\endgroup
For a finite real symmetric matrix, Hermitian symmetry gives an orthonormal eigenbasis, a decreasing eigenvalue list, the product formula for its characteristic polynomial, its root multiset, and its set spectrum.
\begin{itemize}\raggedright
\item Mathlib: \path{Analysis/Matrix/Spectrum.lean}. \code{Matrix.IsHermitian.eigenvalues}\textsubscript{0}, \code{Matrix.IsHermitian.eigenvalues}\textsubscript{0}\code{_antitone}, \code{Matrix.IsHermitian.eigenvalues}, \code{Matrix.IsHermitian.eigenvectorBasis}, \code{Matrix.IsHermitian.mulVec_eigenvectorBasis}, \code{Matrix.IsHermitian.charpoly_eq}, \code{Matrix.IsHermitian.roots_charpoly_eq_eigenvalues}, \code{Matrix.IsHermitian.roots_charpoly_eq_eigenvalues}\textsubscript{0}, \code{Matrix.IsHermitian.spectrum_real_eq_range_eigenvalues}, \code{Matrix.IsHermitian.trace_eq_sum_eigenvalues}, \code{Matrix.IsHermitian.det_eq_prod_eigenvalues}, \code{Matrix.IsHermitian.rank_eq_card_non_zero_eigs}.
\item Mathlib: \path{Data/Fintype/EquivFin.lean}. \code{Fintype.equivOfCardEq}.
\end{itemize}
\begingroup\makeatletter\def\@currentlabel{2}\subsection*{Facts 2: coord}\phantomsection\label{facts:coord}\makeatother\endgroup
For finite-dimensional real inner product spaces, orthonormal bases give coordinates, and symmetric endomorphisms have orthonormal eigenbases. The function and Euclidean coordinate models have inverse linear maps; both matrix actions are the same coordinate sum. The eigenspace predicate is the eigenvector equation, allowing the zero vector.
\begin{itemize}\raggedright
\item Mathlib: \path{Analysis/InnerProductSpace/PiL2.lean}. \code{EuclideanSpace}, \code{Matrix.toEuclideanLin}, \code{OrthonormalBasis.sum_repr}, \code{OrthonormalBasis.reindex}, \code{exists_orthonormalBasis}.
\item Mathlib: \path{Analysis/Normed/Lp/WithLp.lean}. \code{WithLp.toLp}, \code{WithLp.linearEquiv}, \code{WithLp.ofLp_toLp}, \code{WithLp.toLp_ofLp}.
\item Mathlib: \path{Analysis/Normed/Lp/Matrix.lean}. \code{Matrix.toLpLin_apply}.
\item Mathlib: \path{LinearAlgebra/Matrix/ToLin.lean}. \code{Matrix.toLin'}, \code{Matrix.toLin'_apply}.
\item Mathlib: \path{Analysis/Matrix/Hermitian.lean}. \code{Matrix.isSymmetric_toEuclideanLin_iff}.
\item Mathlib: \path{Analysis/InnerProductSpace/Spectrum.lean}. \code{LinearMap.IsSymmetric.apply_eigenvectorBasis}, \code{LinearMap.IsSymmetric.card_filter_eigenvalues_eq}.
\item Mathlib: \path{LinearAlgebra/Eigenspace/Basic.lean}. \code{Module.End.eigenspace}, \code{Module.End.mem_eigenspace_iff}.
\end{itemize}
\begingroup\makeatletter\def\@currentlabel{3}\subsection*{Facts 3: dim}\phantomsection\label{facts:dim}\makeatother\endgroup
Over a field, finite basis cardinalities are dimensions; linear equivalences preserve dimension. Subspace inclusion bounds dimension. Rank-nullity and the dimension formula for a sum and intersection apply in finite dimension. A finite-dimensional endomorphism is injective iff it is surjective.
\begin{itemize}\raggedright
\item Mathlib: \path{LinearAlgebra/Dimension/StrongRankCondition.lean}. \code{Module.finrank_eq_card_basis}.
\item Mathlib: \path{LinearAlgebra/Dimension/Finrank.lean}. \code{LinearEquiv.finrank_eq}.
\item Mathlib: \path{LinearAlgebra/Dimension/Constructions.lean}. \code{Submodule.finrank_le}, \code{Submodule.finrank_mono}.
\item Mathlib: \path{LinearAlgebra/FiniteDimensional/Lemmas.lean}. \code{Submodule.finrank_sup_add_finrank_inf_eq}, \code{LinearMap.finrank_range_add_finrank_ker}.
\item Mathlib: \path{LinearAlgebra/FiniteDimensional/Basic.lean}. \code{LinearMap.injective_iff_surjective}.
\end{itemize}
\begingroup\makeatletter\def\@currentlabel{4}\subsection*{Facts 4: det}\phantomsection\label{facts:det}\makeatother\endgroup
Over a commutative ring, determinants multiply, reindexing both axes by a bijection preserves determinant, and an upper triangular block determinant is the product of its diagonal-block determinants. A square matrix is a unit iff its determinant is a unit; over a field this means a nonzero determinant.
\begin{itemize}\raggedright
\item Mathlib: \path{LinearAlgebra/Matrix/Determinant/Basic.lean}. \code{Matrix.det_mul}, \code{Matrix.det_reindex_self}, \code{Matrix.det_fromBlocks_zero}\textsubscript{21}.
\item Mathlib: \path{LinearAlgebra/Matrix/NonsingularInverse.lean}. \code{Matrix.isUnit_iff_isUnit_det}.
\end{itemize}
\begingroup\makeatletter\def\@currentlabel{5}\subsection*{Facts 5: roots}\phantomsection\label{facts:roots}\makeatother\endgroup
Polynomial roots over a domain count algebraic multiplicity, and a product of linear factors has exactly their roots. For a nonzero product, root multiplicities add; the nonzero qualification is necessary. Two finite multisets are equal iff every count agrees.
\begin{itemize}\raggedright
\item Mathlib: \path{Algebra/Polynomial/Roots.lean}. \code{Polynomial.count_roots}, \code{Polynomial.roots_mul}, \code{Polynomial.roots_multiset_prod_X_sub_C}.
\item Mathlib: \path{Algebra/Polynomial/RingDivision.lean}. \code{Polynomial.rootMultiplicity_mul}.
\item Mathlib: \path{Data/Multiset/Count.lean}. \code{Multiset.ext}.
\end{itemize}
\begingroup\makeatletter\def\@currentlabel{6}\subsection*{Facts 6: positive}\phantomsection\label{facts:positive}\makeatother\endgroup
For finite real square matrices, positive semidefiniteness means symmetry and a nonnegative quadratic form. The zero-form criterion requires positive semidefiniteness. Principal submatrices retain it.
\begin{itemize}\raggedright
\item Mathlib: \path{LinearAlgebra/Matrix/PosDef.lean}. \code{Matrix.posSemidef_iff_dotProduct_mulVec}, \code{Matrix.PosSemidef.submatrix}.
\item Mathlib: \path{Analysis/Matrix/PosDef.lean}. \code{Matrix.IsHermitian.posSemidef_iff_eigenvalues_nonneg}, \code{Matrix.IsHermitian.posDef_iff_eigenvalues_pos}, \code{Matrix.PosSemidef.det_nonneg}.
\item Mathlib: \path{Analysis/InnerProductSpace/Positive.lean}. \code{Matrix.posSemidef_iff_eq_sum_vecMulVec}.
\item Mathlib: \path{Analysis/Matrix/Order.lean}. \code{Matrix.PosSemidef.dotProduct_mulVec_zero_iff}.
\end{itemize}
\begingroup\makeatletter\def\@currentlabel{7}\subsection*{Facts 7: rank}\phantomsection\label{facts:rank}\makeatother\endgroup
For finite real rectangular matrices, transpose preserves rank and both Gram products have the rank of the original matrix.
\begin{itemize}\raggedright
\item Mathlib: \path{LinearAlgebra/Matrix/Rank.lean}. \code{Matrix.rank_transpose}, \code{Matrix.rank_transpose_mul_self}, \code{Matrix.rank_self_mul_transpose}.
\end{itemize}
\begingroup\makeatletter\def\@currentlabel{8}\subsection*{Facts 8: order}\phantomsection\label{facts:order}\makeatother\endgroup
A finite sum of nonnegative real numbers is zero iff every term is zero, and is positive iff one term is positive. In particular this applies to squares. For a nonnegative real number, the square of its nonnegative square root is itself; its square root is positive iff the number is positive. The two finite-sum declarations are generated by the indicated additive attributes in the source.
\begin{itemize}\raggedright
\item Mathlib: \path{Algebra/Order/BigOperators/Group/Finset.lean}. \code{Finset.sum_eq_zero_iff_of_nonneg}, \code{Finset.sum_pos_iff_of_nonneg}.
\item Mathlib: \path{Analysis/Real/Sqrt.lean}. \code{Real.sq_sqrt}, \code{Real.sqrt_pos}.
\end{itemize}
\begingroup\makeatletter\def\@currentlabel{9}\subsection*{Facts 9: charpoly}\phantomsection\label{facts:charpoly}\makeatother\endgroup
Characteristic polynomials of finite square matrices over a commutative ring are monic, commute with coefficient ring maps and are invariant under conjugation by a unit; a diagonal matrix has the product of its linear factors. For rectangular products of sizes $p,q$, the identity is $t^q\chi_{MN}=t^p\chi_{NM}$.
\begin{itemize}\raggedright
\item Mathlib: \path{LinearAlgebra/Matrix/Charpoly/Coeff.lean}. \code{Matrix.charpoly_monic}.
\item Mathlib: \path{LinearAlgebra/Matrix/Charpoly/Basic.lean}. \code{Matrix.charpoly_map}, \code{Matrix.charpoly_mul_comm'}, \code{Matrix.charpoly_units_conj}, \code{Matrix.charpoly_diagonal}.
\end{itemize}
\begingroup\makeatletter\def\@currentlabel{10}\subsection*{Facts 10: graph}\phantomsection\label{facts:graph}\makeatother\endgroup
For finite simple graphs with classical decidable adjacency, the real adjacency matrix is symmetric, its powers count walks, and graph isomorphisms reindex it. The complement adjacency matrix is $J-I-A$. The Mathlib adjacency-matrix structure bundles zero-one entries, symmetry and zero diagonal, and its graph has adjacency matrix the given matrix. The line graph has actual edges as vertices, adjacent when distinct and intersecting. Component and reachability terminology has its usual simple-graph meaning.
\begin{itemize}\raggedright
\item Mathlib: \path{Combinatorics/SimpleGraph/AdjMatrix.lean}. \code{SimpleGraph.adjMatrix}, \code{SimpleGraph.adjMatrix_pow_apply_eq_card_walk}, \code{SimpleGraph.Iso.reindex_adjMatrix}, \code{SimpleGraph.adjMatrix_mulVec_const_apply_of_regular}, \code{SimpleGraph.trace_adjMatrix}, \code{SimpleGraph.adjMatrix_mul_self_apply_self}, \code{SimpleGraph.compl_adjMatrix_eq_adjMatrix_compl}, \code{SimpleGraph.one_add_adjMatrix_add_compl_adjMatrix_eq_of_one}, \code{Matrix.IsAdjMatrix}, \code{Matrix.IsAdjMatrix.toGraph}, \code{Matrix.IsAdjMatrix.adjMatrix_toGraph_eq}.
\item Mathlib: \path{Combinatorics/SimpleGraph/LapMatrix.lean}. \code{SimpleGraph.isHermitian_adjMatrix}.
\item Mathlib: \path{Combinatorics/SimpleGraph/LineGraph.lean}. \code{SimpleGraph.lineGraph}.
\end{itemize}
\begingroup\makeatletter\def\@currentlabel{11}\subsection*{Facts 11: trace}\phantomsection\label{facts:trace}\makeatother\endgroup
An idempotent endomorphism is a projection onto its range. Over $\re$, the finite-dimensional image and kernel are free, so its trace equals its range dimension. Matrix trace and rank agree with the corresponding coordinate linear map.
\begin{itemize}\raggedright
\item Mathlib: \path{LinearAlgebra/Projection.lean}. \code{LinearMap.isProj_range_iff_isIdempotentElem}.
\item Mathlib: \path{LinearAlgebra/Trace.lean}. \code{LinearMap.IsProj.trace}, \code{Matrix.trace_toLin'_eq}.
\item Mathlib: \path{LinearAlgebra/Matrix/Rank.lean}. \code{Matrix.rank_eq_finrank_range_toLin}.
\end{itemize}
\begingroup\makeatletter\def\@currentlabel{12}\subsection*{Facts 12: rayleigh}\phantomsection\label{facts:rayleigh}\makeatother\endgroup
On a nontrivial finite-dimensional real inner product space, a symmetric linear map has eigenvalues equal to the supremum and infimum of its Rayleigh quotient over nonzero vectors.
\begin{itemize}\raggedright
\item Mathlib: \path{Analysis/InnerProductSpace/Rayleigh.lean}. \code{LinearMap.IsSymmetric.hasEigenvalue_iSup_of_finiteDimensional}, \code{LinearMap.IsSymmetric.hasEigenvalue_iInf_of_finiteDimensional}.
\end{itemize}
\begingroup\makeatletter\def\@currentlabel{13}\subsection*{Facts 13: laplacian}\phantomsection\label{facts:laplacian}\makeatother\endgroup
For a finite simple graph the real Laplacian is positive semidefinite; its kernel functions are constant on reachable pairs, and its kernel dimension is the number of components, including isolated vertices.
\begin{itemize}\raggedright
\item Mathlib: \path{Combinatorics/SimpleGraph/LapMatrix.lean}. \code{SimpleGraph.posSemidef_lapMatrix}, \code{SimpleGraph.lapMatrix_mulVec_eq_zero_iff_forall_reachable}, \code{SimpleGraph.card_connectedComponent_eq_finrank_ker_toLin'_lapMatrix}.
\end{itemize}
\begingroup\makeatletter\def\@currentlabel{14}\subsection*{Facts 14: incidence}\phantomsection\label{facts:incidence}\makeatother\endgroup
The incidence matrix on all unordered pairs, with zero columns off the edge set, has product equal to the degree matrix plus adjacency. Restricting away the zero columns is a separate finite-sum argument.
\begin{itemize}\raggedright
\item Mathlib: \path{Combinatorics/SimpleGraph/IncMatrix.lean}. \code{SimpleGraph.incMatrix_mul_transpose}.
\end{itemize}
\begingroup\makeatletter\def\@currentlabel{15}\subsection*{Facts 15: srg}\phantomsection\label{facts:srg}\makeatother\endgroup
The strongly regular matrix equation holds over any semiring; its parameter equation is in natural numbers with truncated subtraction and assumes positive order. Complement parameters also use natural subtraction, so conversion to integer formulas requires bounds.
\begin{itemize}\raggedright
\item Mathlib: \path{Combinatorics/SimpleGraph/StronglyRegular.lean}. \code{SimpleGraph.IsSRGWith.matrix_eq}, \code{SimpleGraph.IsSRGWith.param_eq}, \code{SimpleGraph.IsSRGWith.compl}.
\end{itemize}
\begingroup\makeatletter\def\@currentlabel{16}\subsection*{Facts 16: vandermonde}\phantomsection\label{facts:vandermonde}\makeatother\endgroup
For any number of distinct nodes over a domain, a coefficient vector with all moments of degrees below that number zero is zero. The index size may be zero.
\begin{itemize}\raggedright
\item Mathlib: \path{LinearAlgebra/Vandermonde.lean}. \code{Matrix.eq_zero_of_forall_pow_sum_mul_pow_eq_zero}.
\end{itemize}
\begingroup\makeatletter\def\@currentlabel{17}\subsection*{Facts 17: rank-one}\phantomsection\label{facts:rank-one}\makeatother\endgroup
Over a commutative ring, the determinant of $I+uv^T$ is $1+v^Tu$; in the declaration a singleton auxiliary type indexes the replicated column and row.
\begin{itemize}\raggedright
\item Mathlib: \path{LinearAlgebra/Matrix/SchurComplement.lean}. \code{Matrix.det_one_add_replicateCol_mul_replicateRow}.
\end{itemize}
\begingroup\makeatletter\def\@currentlabel{18}\subsection*{Facts 18: tree}\phantomsection\label{facts:tree}\makeatother\endgroup
A graph is a tree iff it has a nonempty vertex set and a unique simple path between each pair of vertices. Here a tree means connected and acyclic; an edge-preserving subgraph need not be induced.
\begin{itemize}\raggedright
\item Mathlib: \path{Combinatorics/SimpleGraph/Acyclic.lean}. \code{SimpleGraph.isTree_iff_existsUnique_path}.
\end{itemize}
\begingroup\makeatletter\def\@currentlabel{19}\subsection*{Facts 19: rational}\phantomsection\label{facts:rational}\makeatother\endgroup
A rational has a reduced numerator and positive denominator. Coprimality is preserved on taking a positive power, and a common divisor of coprime natural numbers is one. Taking the absolute value of an integer numerator gives these natural-number statements. The monic rational-root theorem also proves the final integrality implication directly.
\begin{itemize}\raggedright
\item Mathlib: \path{Data/Rat/Defs.lean}. \code{Rat.mk'_num_den}, \code{Rat.pos}.
\item Mathlib: \path{Data/Nat/GCD/Basic.lean}. \code{Nat.coprime_pow_left_iff}, \code{Nat.eq_one_of_dvd_coprimes}.
\item Mathlib: \path{RingTheory/Polynomial/RationalRoot.lean}. \code{isInteger_of_is_root_of_monic}.
\end{itemize}
\begingroup\makeatletter\def\@currentlabel{20}\subsection*{Facts 20: counting}\phantomsection\label{facts:counting}\makeatother\endgroup
For a finite type of cardinality $\ell$, the unordered pairs with repetitions have cardinality $\binom{\ell+1}{2}$. These index the independent entries of a symmetric matrix.
\begin{itemize}\raggedright
\item Mathlib: \path{Data/Sym/Card.lean}. \code{Sym2.card}.
\end{itemize}

\section{Definitions and the spectral theorem}

\begin{definition}\label{def:spectrum}
\lean{Matrix.spectrumMultiset, Matrix.eigenvalueMultiplicity}
\uses{facts:coord,facts:spec}
\mbox{}\par\nobreak
\assume{\begin{pfenum}
\item $U$ is finite and $A:U\times U\to\re$.
\end{pfenum}}
For a real square matrix $A$, put
\[
 \chi_A(t)=\det(tI-A),\qquad
 \Spec(A)=\chi_A.\mathrm{roots},\qquad
 \mult_A(\mu)=\operatorname{count}_{\Spec(A)}(\mu).
\]
Here roots form a multiset over $\re$, with algebraic multiplicity.
We will call $\ker(A-\mu I)$ the \emph{eigenspace} $V_A(\mu)$; an
eigenvector is a nonzero member. The definition of $\Spec$ makes sense
without symmetry, but its cardinality can then be smaller than $n$.
\end{definition}

\begin{theorem}[Coordinates, roots, and multiplicity]\label{thm:multiplicity}
\lean{Matrix.IsHermitian.spectrumMultiset_eq_eigenvalues, Matrix.IsHermitian.eigenvalueMultiplicity_eq_finrank, Matrix.IsHermitian.mem_spectrumMultiset_iff, Matrix.IsHermitian.sortedEigenvectorBasis, Matrix.spectrumMultiset_eq_of_orthonormal_eigenbasis}
\uses{def:spectrum,facts:charpoly,facts:coord,facts:dim,facts:roots,facts:spec}
\mbox{}\par\nobreak
\assume{\begin{pfenum}
\item $U$ is finite, $|U|=n$, and $A:U\times U\to\re$ is symmetric.
\end{pfenum}}
\prove{\raggedright There is an orthonormal basis $(\Zv_i)_{i\in U}$ with
$A\Zv_i=\lambda_i\Zv_i$, where
$\lambda_i$ is the value at $i$ of \code{Matrix.IsHermitian.eigenvalues} for $A$, and
\[
 \chi_A(t)=\prod_{i\in U}(t-\lambda_i),\qquad
 \Spec(A)=\ms{\lambda_i:i\in U}.\]
\[
 \mult_A(\mu)=\dim V_A(\mu)\quad(\mu\in\re).
\]
Thus $|\Spec(A)|=n$; membership in $\Spec(A)$, existence of a nonzero
$\mu$-eigenvector, and membership in Mathlib's set-valued
$\operatorname{spectrum}_{\re}(A)$ are equivalent. Multiplicity also
equals \code{A.charpoly.rootMultiplicity mu}.
More generally, if $(\Zu_j)_{j\in I}$ is any orthonormal basis of $\re^U$
with $A\Zu_j=\mu_j\Zu_j$, then $\Spec(A)=\ms{\mu_j:j\in I}$.
There is also a sorted orthonormal eigenbasis $(\Zw_i)_{i\in\operatorname{Fin}(n)}$
with $A\Zw_i=\nu_i\Zw_i$, where
$\nu_i=\code{hA.eigenvalues}\textsubscript{0}(i)$ and $\nu_i\ge\nu_j$ for $i\le j$.
Its eigenvalues list the spectrum in decreasing order:
$\Spec(A)=\ms{\nu_i:i\in\operatorname{Fin}(n)}$.}
\end{theorem}
\begin{proof}
\pfsketch\ We will diagonalize $A$, read its roots with repetitions, and identify the eigenspace coordinates.
\step{mul-basis}{\pick\ an orthonormal basis $(\Zv_i)$ with $A\Zv_i=\lambda_i\Zv_i$; its column matrix $Q$ satisfies $Q^TQ=QQ^T=I$.}
\begin{proof}
\step{mul-real}{$A$ is Hermitian over $\re$.}
\begin{proof}
\pf\ \raggedright Conjugation on $\re$ is the identity; use the definition of conjugate transpose.~\qed
\end{proof}
\step{mul-eigen}{There is the asserted orthonormal eigenbasis.}
\begin{proof}
\pf\ \raggedright Apply \code{Matrix.IsHermitian.eigenvectorBasis} and \code{Matrix.IsHermitian.mulVec_eigenvectorBasis} in Facts~\ref{facts:spec} to \stepref{mul-real}.~\qed
\end{proof}
\step{mul-orth}{$Q^TQ=I$ and $QQ^T=I$.}
\begin{proof}
\step{mul-leftorth}{$Q^TQ=I$.}
\begin{proof}
\pf\ \raggedright Its $(i,j)$ entry is $\ip{\Zv_i}{\Zv_j}=\delta_{ij}$ by \stepref{mul-eigen}.~\qed
\end{proof}
\step{mul-invertible}{$Q$ is invertible.}
\begin{proof}
\pf\ \raggedright Its columns form a basis by \stepref{mul-eigen}; the basis coordinate map is its inverse.~\qed
\end{proof}
\step{mul-rightorth}{$QQ^T=I$.}
\begin{proof}
\pf\ \raggedright The left inverse in \stepref{mul-leftorth} equals the inverse of \stepref{mul-invertible}.~\qed
\end{proof}
\qedstep
\begin{proof}
\pf\ \raggedright Use \stepref{mul-leftorth}, \stepref{mul-rightorth}.~\qed
\end{proof}
\end{proof}
\qedstep
\begin{proof}
\pf\ \raggedright Use \stepref{mul-eigen}, \stepref{mul-orth}.~\qed
\end{proof}
\end{proof}
\step{mul-roots}{$A=Q\diag(\lambda_i)Q^T$, $\chi_A=\prod_i(t-\lambda_i)$ and $\Spec(A)=\ms{\lambda_i:i\in U}$.}
\begin{proof}
\step{mul-columns}{$AQ=Q\diag(\lambda_i)$.}
\begin{proof}
\pf\ \raggedright Its $i$th column is the eigenvector equation in \stepref{mul-basis}.~\qed
\end{proof}
\step{mul-diag}{$A=Q\diag(\lambda_i)Q^T$.}
\begin{proof}
\pf\ \raggedright Multiply \stepref{mul-columns} on the right by $Q^T$ and use \stepref{mul-basis}.~\qed
\end{proof}
\step{mul-factors}{$\chi_A=\prod_i(t-\lambda_i)$.}
\begin{proof}
\pf\ \raggedright Apply \code{Matrix.IsHermitian.charpoly_eq} in Facts~\ref{facts:spec}; it is the determinant of the conjugation in \stepref{mul-diag}.~\qed
\end{proof}
\step{mul-list}{$\Spec(A)=\ms{\lambda_i:i\in U}$.}
\begin{proof}
\pf\ \raggedright Apply \code{Matrix.IsHermitian.roots_charpoly_eq_eigenvalues} in Facts~\ref{facts:spec} and Definition~\ref{def:spectrum}.~\qed
\end{proof}
\qedstep
\begin{proof}
\pf\ \raggedright Use \stepref{mul-diag}, \stepref{mul-factors}, \stepref{mul-list}.~\qed
\end{proof}
\end{proof}
\step{mul-anybasis}{\assume{$I$ is finite and $(\Zu_j)_{j\in I}$ is an orthonormal basis of $\re^U$ with $A\Zu_j=\mu_j\Zu_j$.}\prove{$\Spec(A)=\ms{\mu_j:j\in I}$.}}
\begin{proof}
\pflet{$P:U\times I\to\re$ has columns $\Zu_j$, and $D\triangleq\diag(\mu_j)$ on $I$.}
\step{mul-any-orth}{$P^TP=I$ and $|I|=n$.}
\begin{proof}
\pf\ \raggedright The entries of $P^TP$ are $\ip{\Zu_j}{\Zu_k}=\delta_{jk}$ by \asm{1}; a basis of $\re^U$ has $n$ members by \code{Module.finrank_eq_card_basis} in Facts~\ref{facts:dim}.~\qed
\end{proof}
\step{mul-any-diag}{$A=PDP^T$.}
\begin{proof}
\pf\ \raggedright Both sides act on each $\Zu_k$ by $\mu_k$, by \asm{1} and \stepref{mul-any-orth}; a linear map is determined by its values on a basis.~\qed
\end{proof}
\step{mul-any-charpoly}{$\chi_A=\prod_j(t-\mu_j)$.}
\begin{proof}
\pf\ \raggedright Apply \code{Matrix.charpoly_mul_comm'} in Facts~\ref{facts:charpoly} to $P$ and $DP^T$: with \stepref{mul-any-diag} and \stepref{mul-any-orth}, $t^n\chi_A=t^n\chi_{DP^TP}=t^n\chi_D$. Cancel $t^n$ in the domain $\re[t]$ and apply \code{Matrix.charpoly_diagonal} in Facts~\ref{facts:charpoly}.~\qed
\end{proof}
\step{mul-any-roots}{$\Spec(A)=\ms{\mu_j:j\in I}$.}
\begin{proof}
\pf\ \raggedright Apply \code{Polynomial.roots_multiset_prod_X_sub_C} in Facts~\ref{facts:roots} to \stepref{mul-any-charpoly} and Definition~\ref{def:spectrum}.~\qed
\end{proof}
\qedstep
\begin{proof}
\pf\ \raggedright Use \stepref{mul-any-roots}.~\qed
\end{proof}
\end{proof}
\step{mul-space}{For each $\mu\in\re$, the multiset count equals the Euclidean eigenspace dimension.}
\begin{proof}
\step{mul-count-reindex}{The count of $\mu$ is the number of sorted indices with eigenvalue $\mu$.}
\begin{proof}
\pf\ \raggedright Apply \code{Matrix.IsHermitian.roots_charpoly_eq_eigenvalues}\textsubscript{0} in Facts~\ref{facts:spec} and the definition of a mapped-multiset count.~\qed
\end{proof}
\step{mul-filter}{This number equals the Euclidean eigenspace dimension.}
\begin{proof}
\pf\ \raggedright Apply \code{LinearMap.IsSymmetric.card_filter_eigenvalues_eq} in Facts~\ref{facts:coord} to $A.\mathrm{toEuclideanLin}$, using \code{Matrix.isSymmetric_toEuclideanLin_iff} and $\dim\re^U=n$. Its eigenvalues are the sorted list by definition.~\qed
\end{proof}
\qedstep
\begin{proof}
\pf\ \raggedright Use \stepref{mul-count-reindex}, \stepref{mul-filter}.~\qed
\end{proof}
\end{proof}
\step{mul-bridge}{The two Lean eigenspace coordinate conventions have equal dimensions.}
\begin{proof}
\step{mul-map}{The map $\Zx\mapsto\code{WithLp.toLp}\,2\,\Zx$ has inverse taking the underlying function.}
\begin{proof}
\pf\ \raggedright These are the inverse coordinate maps defining $\code{EuclideanSpace}$ in Facts~\ref{facts:coord}.~\qed
\end{proof}
\step{mul-intertwine}{This equivalence intertwines $\code{A.toLin'}$ and $\code{A.toEuclideanLin}$.}
\begin{proof}
\pf\ \raggedright Both underlying functions have value $\sum_v A(u,v)\Zx(v)$ at $u$, by Facts~\ref{facts:coord}.~\qed
\end{proof}
\step{mul-restrict}{The equivalence restricts to inverse linear maps of their $\mu$-eigenspaces.}
\begin{proof}
\pf\ \raggedright Use \stepref{mul-intertwine} in the defining eigenvector equations in Facts~\ref{facts:coord}.~\qed
\end{proof}
\qedstep
\begin{proof}
\pf\ \raggedright Linear equivalences preserve dimension (Facts~\ref{facts:dim}); apply \stepref{mul-restrict}.~\qed
\end{proof}
\end{proof}
\step{mul-membership}{The cardinality, three membership criteria and root-multiplicity statement hold, also when $U$ is empty.}
\begin{proof}
\step{mul-card}{$|\Spec(A)|=n$.}
\begin{proof}
\pf\ \raggedright Count the $U$-indexed multiset in \stepref{mul-roots}.~\qed
\end{proof}
\step{mul-nonzero}{$\mu\in\Spec(A)$ iff there is a nonzero $\mu$-eigenvector.}
\begin{proof}
\pf\ \raggedright Combine \stepref{mul-space} and \stepref{mul-bridge}; positive dimension is equivalent to a nonzero vector by Facts~\ref{facts:dim}.~\qed
\end{proof}
\step{mul-spectrum}{$\mu\in\Spec(A)$ iff $\mu\in\operatorname{spectrum}_{\re}(A)$.}
\begin{proof}
\pf\ \raggedright Combine \stepref{mul-roots} with \code{Matrix.IsHermitian.spectrum_real_eq_range_eigenvalues} in Facts~\ref{facts:spec}.~\qed
\end{proof}
\step{mul-rootmult}{$\mult_A(\mu)=\chi_A.\mathrm{rootMultiplicity}(\mu)$.}
\begin{proof}
\pf\ \raggedright Apply \code{Polynomial.count_roots} in Facts~\ref{facts:roots} to Definition~\ref{def:spectrum}.~\qed
\end{proof}
\qedstep
\begin{proof}
\pf\ \raggedright Use \stepref{mul-card}, \stepref{mul-nonzero}, \stepref{mul-spectrum}, \stepref{mul-rootmult}.~\qed
\end{proof}
\end{proof}
\step{mul-sorted}{The eigenvalues $\nu_i=\code{hA.eigenvalues}\textsubscript{0}(i)$ have a sorted orthonormal eigenbasis $(\Zw_i)_{i\in\operatorname{Fin}(n)}$, and $\Spec(A)=\ms{\nu_i:i\in\operatorname{Fin}(n)}$.}
\begin{proof}
\pflet{$e\triangleq\code{Fintype.equivOfCardEq (Fintype.card_fin n)}:\operatorname{Fin}(n)\simeq U$, from Facts~\ref{facts:spec}, and $\Zw_i\triangleq\code{hA.eigenvectorBasis}(e(i))$.}
\step{mul-reindex}{$(\Zw_i)$ is an orthonormal basis.}
\begin{proof}
\pf\ \raggedright It is \code{hA.eigenvectorBasis.reindex} $e^{-1}$; apply \code{OrthonormalBasis.reindex} in Facts~\ref{facts:coord}.~\qed
\end{proof}
\step{mul-sorted-equation}{$A\Zw_i=\nu_i\Zw_i$.}
\begin{proof}
\pf\ \raggedright Apply \code{Matrix.IsHermitian.mulVec_eigenvectorBasis} at $e(i)$ and unfold \code{Matrix.IsHermitian.eigenvalues}, which is defined with this $e$; $e^{-1}(e(i))=i$.~\qed
\end{proof}
\step{mul-antitone}{$\nu_i\ge\nu_j$ whenever $i\le j$.}
\begin{proof}
\pf\ \raggedright Apply \code{Matrix.IsHermitian.eigenvalues}\textsubscript{0}\code{_antitone} in Facts~\ref{facts:spec}.~\qed
\end{proof}
\step{mul-sorted-roots}{$\Spec(A)=\ms{\nu_i:i\in\operatorname{Fin}(n)}$.}
\begin{proof}
\pf\ \raggedright Apply \code{Matrix.IsHermitian.roots_charpoly_eq_eigenvalues}\textsubscript{0} in Facts~\ref{facts:spec} and Definition~\ref{def:spectrum}.~\qed
\end{proof}
\qedstep
\begin{proof}
\pf\ \raggedright Use \stepref{mul-reindex}, \stepref{mul-sorted-equation}, \stepref{mul-antitone}, \stepref{mul-sorted-roots}.~\qed
\end{proof}
\end{proof}
\qedstep
\begin{proof}
\pf\ \raggedright Use \stepref{mul-basis}, \stepref{mul-roots}, \stepref{mul-anybasis}, \stepref{mul-space}, \stepref{mul-bridge}, \stepref{mul-membership}, \stepref{mul-sorted}.~\qed
\end{proof}
\end{proof}

\begin{definition}\label{def:extremes}
\lean{Matrix.IsHermitian.largestEigenvalue, Matrix.IsHermitian.leastEigenvalue}
\uses{facts:spec}
\mbox{}\par\nobreak
\assume{\begin{pfenum}
\item $U$ is finite and nonempty, and $A:U\times U\to\re$ is symmetric.
\end{pfenum}}
For a real symmetric matrix $A$ on $n>0$ indices, define
\[
 \lambda_{\max}(A)=\code{hA.eigenvalues}\textsubscript{0}(0),\qquad
 \lambda_{\min}(A)=\code{hA.eigenvalues}\textsubscript{0}(n-1).
\]
The indices carry their proofs of being less than $n$. These are exactly
\code{LineGraphs/Spectral/Basic.lean}'s definitions. They require
\code{Nonempty U} and \code{hA : A.IsHermitian}; there is no extreme
for the empty matrix.
\end{definition}

\begin{lemma}[Extremal entries and multiset bounds]\label{lem:extremes}
\lean{Matrix.IsHermitian.largestEigenvalue_mem_spectrumMultiset, Matrix.IsHermitian.leastEigenvalue_mem_spectrumMultiset, Matrix.IsHermitian.le_largestEigenvalue_of_mem, Matrix.IsHermitian.leastEigenvalue_le_of_mem}
\uses{def:extremes,facts:spec,thm:multiplicity}
\mbox{}\par\nobreak
\assume{\begin{pfenum}
\item $U$ is finite and nonempty, $|U|=n$, and $A:U\times U\to\re$ is symmetric.
\end{pfenum}}
\prove{\raggedright $\lambda_{\max}(A),\lambda_{\min}(A)\in\Spec(A)$, and
\[
 \lambda_{\min}(A)\le\mu\le\lambda_{\max}(A)
 \qquad(\mu\in\Spec(A)).
\]
Thus they are the maximum and minimum of the multiset and equal
the first and last entries of the decreasing list.}
\end{lemma}
\begin{proof}
\pfsketch\ We will read the two endpoints of the decreasing list $\nu_0\ge\cdots\ge\nu_{n-1}$ of Theorem~\ref{thm:multiplicity} and compare every entry with them.
\step{ext-roots}{$\Spec(A)=\ms{\nu_i:i\in\operatorname{Fin}(n)}$.}
\begin{proof}
\pf\ \raggedright Theorem~\ref{thm:multiplicity}.~\qed
\end{proof}
\step{ext-members}{$\lambda_{\max}(A)=\nu_0$ and $\lambda_{\min}(A)=\nu_{n-1}$ belong to $\Spec(A)$.}
\begin{proof}
\pf\ \raggedright The equalities are Definition~\ref{def:extremes}; use indices $0$ and $n-1$ in \stepref{ext-roots}, which are valid because $n>0$.~\qed
\end{proof}
\step{ext-upper}{Every $\mu\in\Spec(A)$ satisfies $\mu\le\nu_0$.}
\begin{proof}
\pf\ \raggedright By \stepref{ext-roots}, $\mu=\nu_i$ for some $i$; apply \code{Matrix.IsHermitian.eigenvalues}\textsubscript{0}\code{_antitone} in Facts~\ref{facts:spec} to $0\le i$.~\qed
\end{proof}
\step{ext-lower}{Every $\mu\in\Spec(A)$ satisfies $\nu_{n-1}\le\mu$.}
\begin{proof}
\pf\ \raggedright By \stepref{ext-roots}, $\mu=\nu_i$ with $i\le n-1$; apply the same declaration.~\qed
\end{proof}
\qedstep
\begin{proof}
\pf\ \raggedright \stepref{ext-members}, \stepref{ext-upper} and \stepref{ext-lower} are the membership and the two bounds; membership with a bound is the defining property of a maximum or minimum.~\qed
\end{proof}
\end{proof}

\begin{definition}\label{def:graph}
\lean{SimpleGraph.spectrumMultiset, SimpleGraph.eigenvalueMultiplicity, SimpleGraph.Cospectral, SimpleGraph.largestEigenvalue, SimpleGraph.leastEigenvalue, SimpleGraph.seidelMatrix}
\uses{def:extremes,def:spectrum,facts:graph}
\mbox{}\par\nobreak
\assume{\begin{pfenum}
\item $X$ and $Y$ are finite simple graphs; their vertex types may differ.
\end{pfenum}}
The spectrum and multiplicity of a graph $X$ are those of its adjacency
matrix $A(X)$. Two graphs, possibly on different vertex types, are
\emph{cospectral} when their real adjacency characteristic polynomials
are equal. The least and largest graph eigenvalues require a nonempty
vertex set. The \emph{Seidel matrix} is
\[
 S(X)=J-I-2A(X).
\]
It is symmetric, has diagonal zero, and off-diagonal entry $-1$ on an
edge and $1$ on a nonedge.
\end{definition}

\begin{proposition}[Equivalent spectra]\label{prop:cospectral}
\lean{Matrix.IsHermitian.spectrumMultiset_eq_iff_charpoly_eq, SimpleGraph.cospectral_iff_spectrumMultiset_eq, SimpleGraph.Iso.cospectral}
\uses{def:graph,def:spectrum,facts:det,facts:graph,facts:roots,thm:multiplicity}
\mbox{}\par\nobreak
\assume{\begin{pfenum}
\item $U,W$ are finite; $A:U\times U\to\re$ and $B:W\times W\to\re$ are symmetric.
\item For the graph assertions, $X$ and $Y$ are finite simple graphs.
\end{pfenum}}
\prove{\raggedright For real symmetric matrices $A,B$, even on different finite index sets,
$\chi_A=\chi_B$ iff $\Spec(A)=\Spec(B)$ iff
$\mult_A(\mu)=\mult_B(\mu)$ for all real $\mu$.
Thus cospectral graphs have the same number of vertices. Graph
isomorphisms preserve the spectrum.}
\end{proposition}
\begin{proof}
\pfsketch\ We will reconstruct the monic characteristic polynomial from its roots, then relabel adjacency matrices.
\step{cos-equiv}{The three equalities are equivalent.}
\begin{proof}
\step{cos-forward}{$\chi_A=\chi_B$ implies $\Spec(A)=\Spec(B)$.}
\begin{proof}
\pf\ \raggedright Apply the roots function in Definition~\ref{def:spectrum}.~\qed
\end{proof}
\step{cos-back}{$\Spec(A)=\Spec(B)$ implies $\chi_A=\chi_B$.}
\begin{proof}
\pf\ \raggedright Substitute the common multiset in the product formula of Theorem~\ref{thm:multiplicity}.~\qed
\end{proof}
\step{cos-counts}{Multiset equality is equivalent to equality of every count.}
\begin{proof}
\pf\ \raggedright Apply \code{Multiset.ext} in Facts~\ref{facts:roots}.~\qed
\end{proof}
\qedstep
\begin{proof}
\pf\ \raggedright Use \stepref{cos-forward}, \stepref{cos-back}, \stepref{cos-counts}.~\qed
\end{proof}
\end{proof}
\step{cos-graphs}{Cospectral graphs have equal order; isomorphic graphs are cospectral.}
\begin{proof}
\step{cos-orders}{Equal spectra have equal orders.}
\begin{proof}
\pf\ \raggedright Take cardinalities and use Theorem~\ref{thm:multiplicity}.~\qed
\end{proof}
\step{cos-reindex}{An isomorphism reindexes the adjacency matrix.}
\begin{proof}
\pf\ \raggedright Apply \code{SimpleGraph.Iso.reindex_adjMatrix} in Facts~\ref{facts:graph}.~\qed
\end{proof}
\step{cos-poly}{The reindexed characteristic polynomial is unchanged.}
\begin{proof}
\pf\ \raggedright Apply \code{Matrix.det_reindex_self} in Facts~\ref{facts:det} to $tI-A$ from \stepref{cos-reindex}.~\qed
\end{proof}
\qedstep
\begin{proof}
\pf\ \raggedright Use \stepref{cos-orders}, \stepref{cos-poly}.~\qed
\end{proof}
\end{proof}
\qedstep
\begin{proof}
\pf\ \raggedright Use \stepref{cos-equiv}, \stepref{cos-graphs}.~\qed
\end{proof}
\end{proof}

\section{Polynomial calculus, traces, and positivity}

\begin{theorem}[Polynomial calculus and spectral idempotents]\label{thm:calculus}
\lean{Matrix.IsHermitian.spectralIdempotent, Matrix.IsHermitian.aeval_eq_sum_spectralIdempotent, Matrix.IsHermitian.spectrumMultiset_aeval, Matrix.IsHermitian.aeval_eq_zero_iff, Matrix.IsHermitian.trace_spectralIdempotent}
\uses{thm:multiplicity}
\mbox{}\par\nobreak
\assume{\begin{pfenum}
\item $U$ is finite, $|U|=n$, and $A:U\times U\to\re$ is symmetric.
\item $(\theta_r)_{r\in R}$ lists the distinct eigenvalues in decreasing order; $R$ is empty if $n=0$.
\end{pfenum}}
\prove{\raggedright There are
unique symmetric matrices $E_r$ that are the identity on $V_A(\theta_r)$
and zero on every other eigenspace. They satisfy
\[
 \sum_rE_r=I,\quad E_rE_s=\begin{cases}E_r&r=s,\\0&r\ne s,\end{cases}
 \qquad A=\sum_r\theta_rE_r.
\]
For every real polynomial $f$,
\[
 f(A)=\sum_r f(\theta_r)E_r,\quad
 \Spec(f(A))=\operatorname{map}(f,\Spec(A)),\quad
 \mult_{f(A)}(b)=\sum_{r:f(\theta_r)=b}\mult_A(\theta_r).
\]
In particular $f(A)=0$ iff $f(\theta_r)=0$ for every $r$, and
\[
 E_r=\prod_{s\ne r}\frac{A-\theta_sI}{\theta_r-\theta_s},\quad
 \tr E_r=\rk E_r=\mult_A(\theta_r).
\]}
\end{theorem}
\begin{proof}
\pfsketch\ We will select the coordinates of each eigenspace and evaluate polynomials on the eigenbasis.
\step{calc-basis}{\pick\ an orthonormal eigenbasis $(\Zv_i)$ of $A$, with eigenvalues $\lambda_i$; let $Q$ have these columns.}
\begin{proof}
\pf\ \raggedright Theorem~\ref{thm:multiplicity}.~\qed
\end{proof}
\step{calc-proj}{Define $E_r=Q D_rQ^T$, where $D_r(i,i)=1$ if $\lambda_i=\theta_r$ and $0$ otherwise. These are the unique asserted symmetric matrices.}
\begin{proof}
\step{calc-select}{$E_r\Zv_i=\Zv_i$ if $\lambda_i=\theta_r$, and $0$ otherwise.}
\begin{proof}
\pf\ \raggedright Orthonormality of the basis of Theorem~\ref{thm:multiplicity} makes $Q^T\Zv_i$ the $i$th coordinate unit vector; apply $QD_r$ to it.~\qed
\end{proof}
\step{calc-sym}{$E_r^T=E_r$.}
\begin{proof}
\pf\ \raggedright Transpose $QD_rQ^T$, using $D_r^T=D_r$.~\qed
\end{proof}
\step{calc-unique}{Any matrix with the asserted eigenspace actions equals $E_r$.}
\begin{proof}
\pf\ \raggedright The actions in \stepref{calc-select} determine every column in a basis.~\qed
\end{proof}
\qedstep
\begin{proof}
\pf\ \raggedright Use \stepref{calc-select}, \stepref{calc-sym}, \stepref{calc-unique}.~\qed
\end{proof}
\end{proof}
\step{calc-algebra}{$\sum_rE_r=I$, $E_rE_s=\delta_{rs}E_r$ and $A=\sum_r\theta_rE_r$.}
\begin{proof}
\step{calc-sum}{Exactly one $E_r$ fixes each $\Zv_i$.}
\begin{proof}
\pf\ \raggedright The list of distinct values contains $\lambda_i$ exactly once; use \stepref{calc-proj}.~\qed
\end{proof}
\step{calc-product}{$E_rE_s\Zv_i=\delta_{rs}E_r\Zv_i$.}
\begin{proof}
\pf\ \raggedright Apply the two indicator actions in \stepref{calc-proj}.~\qed
\end{proof}
\step{calc-weight}{$\sum_r\theta_rE_r\Zv_i=\lambda_i\Zv_i=A\Zv_i$.}
\begin{proof}
\pf\ \raggedright Use the unique surviving term in \stepref{calc-sum} and Theorem~\ref{thm:multiplicity}.~\qed
\end{proof}
\qedstep
\begin{proof}
\pf\ \raggedright Equality on a basis gives the three matrix identities from \stepref{calc-sum}, \stepref{calc-product} and \stepref{calc-weight}.~\qed
\end{proof}
\end{proof}
\step{calc-poly}{$f(A)\Zv_i=f(\lambda_i)\Zv_i$ and $f(A)=\sum_r f(\theta_r)E_r$.}
\begin{proof}
\step{calc-power0}{$A^0\Zv_i=\lambda_i^0\Zv_i$.}
\begin{proof}
\pf\ \raggedright Both sides equal $\Zv_i$, by the definition of zeroth power.~\qed
\end{proof}
\step{calc-powerS}{If $A^j\Zv_i=\lambda_i^j\Zv_i$, then $A^{j+1}\Zv_i=\lambda_i^{j+1}\Zv_i$.}
\begin{proof}
\pf\ \raggedright Multiply the assumed equation by $A$ and use Theorem~\ref{thm:multiplicity}.~\qed
\end{proof}
\step{calc-powers}{$A^j\Zv_i=\lambda_i^j\Zv_i$ for every $j\in\nn$.}
\begin{proof}
\pf\ \raggedright Induction with \stepref{calc-power0} and \stepref{calc-powerS}.~\qed
\end{proof}
\step{calc-eval}{$f(A)\Zv_i=f(\lambda_i)\Zv_i$.}
\begin{proof}
\pf\ \raggedright Take the finite linear combination of \stepref{calc-powers} with the coefficients of $f$, as specified by polynomial evaluation.~\qed
\end{proof}
\step{calc-equality}{$f(A)=\sum_r f(\theta_r)E_r$.}
\begin{proof}
\pf\ \raggedright The two matrices have the same action on each $\Zv_i$ by \stepref{calc-eval} and \stepref{calc-proj}.~\qed
\end{proof}
\qedstep
\begin{proof}
\pf\ \raggedright Use \stepref{calc-eval}, \stepref{calc-equality}.~\qed
\end{proof}
\end{proof}
\step{calc-multi}{$\Spec(f(A))$ is the mapped multiset; the displayed multiplicities add over equal images, and $f(A)=0$ iff all $f(\theta_r)=0$.}
\begin{proof}
\step{calc-conj}{$f(A)=Q\diag(f(\lambda_i))Q^T$.}
\begin{proof}
\pf\ \raggedright Read the basis action in \stepref{calc-poly}.~\qed
\end{proof}
\step{calc-rootlist}{$\Spec(f(A))=\ms{f(\lambda_i):i\in U}$.}
\begin{proof}
\pf\ \raggedright $f(A)=\sum_rf(\theta_r)E_r$ is symmetric by \stepref{calc-poly} and \stepref{calc-proj}, and $(\Zv_i)$ is an orthonormal eigenbasis of it with eigenvalues $f(\lambda_i)$ by \stepref{calc-poly}; apply the arbitrary-basis assertion of Theorem~\ref{thm:multiplicity}.~\qed
\end{proof}
\step{calc-collisions}{$\mult_{f(A)}(b)=\sum_{r:f(\theta_r)=b}\mult_A(\theta_r)$.}
\begin{proof}
\pf\ \raggedright Partition the indices counted in \stepref{calc-rootlist} by their original value $\lambda_i$.~\qed
\end{proof}
\step{calc-zero}{$f(A)=0$ iff $f(\theta_r)=0$ for all $r$.}
\begin{proof}
\pf\ \raggedright The diagonal matrix in \stepref{calc-conj} is zero exactly when each diagonal entry is zero.~\qed
\end{proof}
\qedstep
\begin{proof}
\pf\ \raggedright Use \stepref{calc-rootlist}, \stepref{calc-collisions}, \stepref{calc-zero}.~\qed
\end{proof}
\end{proof}
\step{calc-interp}{The product formula for $E_r$ and $\tr E_r=\rk E_r=\mult_A(\theta_r)$ hold.}
\begin{proof}
\step{calc-denom}{For $s\ne r$, $\theta_r-\theta_s\ne0$.}
\begin{proof}
\pf\ \raggedright This is distinctness in \asm[2]{0}.~\qed
\end{proof}
\step{calc-value1}{$\prod_{s\ne r}(\theta_r-\theta_s)/(\theta_r-\theta_s)=1$.}
\begin{proof}
\pf\ \raggedright Each factor is $1$ by \stepref{calc-denom}.~\qed
\end{proof}
\step{calc-value0}{For $t\ne r$, $\prod_{s\ne r}(\theta_t-\theta_s)/(\theta_r-\theta_s)=0$.}
\begin{proof}
\pf\ \raggedright The factor $s=t$ is zero.~\qed
\end{proof}
\step{calc-productformula}{The displayed product equals $E_r$.}
\begin{proof}
\pf\ \raggedright Its basis action agrees with \stepref{calc-proj} by \stepref{calc-poly}, \stepref{calc-value1} and \stepref{calc-value0}.~\qed
\end{proof}
\step{calc-rank}{$\rk E_r=\#\{i:\lambda_i=\theta_r\}$.}
\begin{proof}
\pf\ \raggedright Its image has precisely that coordinate basis by \stepref{calc-proj}.~\qed
\end{proof}
\step{calc-trace}{$\tr E_r=\sum_{i:\lambda_i=\theta_r}\sum_u\Zv_i(u)^2=\#\{i:\lambda_i=\theta_r\}$.}
\begin{proof}
\pf\ \raggedright Expand $QD_rQ^T$ on the diagonal and use unit norms from Theorem~\ref{thm:multiplicity}.~\qed
\end{proof}
\qedstep
\begin{proof}
\pf\ \raggedright Use \stepref{calc-productformula}, \stepref{calc-rank}, \stepref{calc-trace}.~\qed
\end{proof}
\end{proof}
\qedstep
\begin{proof}
\pf\ \raggedright Use \stepref{calc-proj}, \stepref{calc-algebra}, \stepref{calc-poly}, \stepref{calc-multi}, \stepref{calc-interp}.~\qed
\end{proof}
\end{proof}

\begin{theorem}[Trace, determinant, rank, and moments]\label{thm:trace}
\lean{Matrix.IsHermitian.trace_pow_eq_sum_spectrumMultiset, Matrix.IsHermitian.rank_aeval_eq_sum_multiplicity, Matrix.trace_eq_rank_of_idempotent}
\uses{facts:spec,facts:trace,thm:calculus,thm:multiplicity}
\mbox{}\par\nobreak
\assume{\begin{pfenum}
\item $A:U\times U\to\re$ is symmetric, $U$ is finite with $|U|=n$, and $(\lambda_i)_{i=1}^n$ lists its eigenvalues with multiplicity.
\item $j\in\nn$, $f\in\re[t]$, and $(\theta_r)$ lists the distinct eigenvalues of $A$; in the last assertion $E:W\times W\to\re$ is idempotent and $W$ is finite.
\end{pfenum}}
\prove{\raggedright \[
 \tr A^j=\sum_{i=1}^n\lambda_i^j,\qquad
 \det A=\prod_{i=1}^n\lambda_i,\qquad
 \rk A=n-\mult_A(0).
\]
For any real polynomial $f$,
$\tr f(A)=\sum_i f(\lambda_i)$ and
$\rk f(A)=\sum_{r:f(\theta_r)\ne0}\mult_A(\theta_r)$.
If $E$ is any real idempotent square matrix, even if not symmetric,
$\tr E=\rk E$ and this number is a nonnegative integer.}
\end{theorem}
\begin{proof}
\pfsketch\ We will compute in diagonal coordinates; for an arbitrary idempotent we will use its image and kernel.
\step{trace-values}{The trace, determinant, rank and polynomial formulas hold.}
\begin{proof}
\step{trace-first}{$\tr A=\sum_i\lambda_i$.}
\begin{proof}
\pf\ \raggedright Apply \code{Matrix.IsHermitian.trace_eq_sum_eigenvalues} in Facts~\ref{facts:spec} to $A$.~\qed
\end{proof}
\step{trace-product}{$\det A=\prod_i\lambda_i$.}
\begin{proof}
\pf\ \raggedright Apply \code{Matrix.IsHermitian.det_eq_prod_eigenvalues} in Facts~\ref{facts:spec}.~\qed
\end{proof}
\step{trace-count}{$\rk A=n-\mult_A(0)$.}
\begin{proof}
\pf\ \raggedright Apply \code{Matrix.IsHermitian.rank_eq_card_non_zero_eigs} in Facts~\ref{facts:spec}; Theorem~\ref{thm:multiplicity} identifies the complementary zero count.~\qed
\end{proof}
\step{trace-polysymmetric}{$f(A)$ is real symmetric and its multiset is $\ms{f(\lambda_i)}$.}
\begin{proof}
\pf\ \raggedright The symmetric spectral idempotents and multiset formula in Theorem~\ref{thm:calculus} give both assertions.~\qed
\end{proof}
\step{trace-f}{$\tr f(A)=\sum_i f(\lambda_i)$.}
\begin{proof}
\pf\ \raggedright Apply \code{Matrix.IsHermitian.trace_eq_sum_eigenvalues} to $f(A)$ using \stepref{trace-polysymmetric}; sum the identified multiset.~\qed
\end{proof}
\step{trace-frank}{$\rk f(A)=\sum_{r:f(\theta_r)\ne0}\mult_A(\theta_r)$.}
\begin{proof}
\pf\ \raggedright Apply \code{Matrix.IsHermitian.rank_eq_card_non_zero_eigs} to $f(A)$ and group the nonzero entries of \stepref{trace-polysymmetric} by original eigenvalue.~\qed
\end{proof}
\step{trace-powers}{$\tr A^j=\sum_i\lambda_i^j$.}
\begin{proof}
\pf\ \raggedright Apply the trace declaration directly to the symmetric matrix $A^j$; Theorem~\ref{thm:calculus} with $t^j$ identifies its multiset, including $j=0$.~\qed
\end{proof}
\qedstep
\begin{proof}
\pf\ \raggedright Use \stepref{trace-first}, \stepref{trace-product}, \stepref{trace-count}, \stepref{trace-f}, \stepref{trace-frank}, \stepref{trace-powers}.~\qed
\end{proof}
\end{proof}
\step{trace-idem}{$\tr E=\rk E\in\nn$.}
\begin{proof}
\step{trace-projection}{$E.\mathrm{toLin'}$ projects onto its range.}
\begin{proof}
\pf\ \raggedright Use \code{LinearMap.isProj_range_iff_isIdempotentElem} in Facts~\ref{facts:trace} with the hypothesis $E^2=E$.~\qed
\end{proof}
\step{trace-idemvalue}{The trace equals the dimension of that range.}
\begin{proof}
\pf\ \raggedright Apply \code{LinearMap.IsProj.trace} in Facts~\ref{facts:trace} to \stepref{trace-projection}; its range and kernel are finite and free over $\re$.~\qed
\end{proof}
\step{trace-coordinate}{This is the matrix trace and matrix rank.}
\begin{proof}
\pf\ \raggedright Apply \code{Matrix.trace_toLin'_eq} and \code{Matrix.rank_eq_finrank_range_toLin} in Facts~\ref{facts:trace}, using the standard coordinate bases for the rank equality; dimension is a natural number.~\qed
\end{proof}
\qedstep
\begin{proof}
\pf\ \raggedright Use \stepref{trace-coordinate}.~\qed
\end{proof}
\end{proof}
\qedstep
\begin{proof}
\pf\ \raggedright Use \stepref{trace-values}, \stepref{trace-idem}.~\qed
\end{proof}
\end{proof}

\begin{theorem}[Positivity and shifts]\label{thm:psd}
\lean{Matrix.IsHermitian.posSemidef_iff_spectrumMultiset_nonneg, Matrix.IsHermitian.posSemidef_sub_smul_one_iff, Matrix.IsHermitian.posSemidef_smul_one_sub_iff}
\uses{facts:positive,thm:calculus,thm:multiplicity}
\mbox{}\par\nobreak
\assume{\begin{pfenum}
\item $U$ is finite and $A:U\times U\to\re$ is symmetric.
\item $b\in\re$ and $\Zx\in\re^U$. The principal-submatrix assertions assume $A\succeq0$; the last assertion also assumes $B:U\times U\to\re$ is positive semidefinite.
\end{pfenum}}
\prove{\raggedright It is positive semidefinite iff every member
of $\Spec(A)$ is nonnegative, and positive definite iff every member is
positive. For any $b\in\re$,
\[
 A-bI\succeq0\ \Longleftrightarrow\ (\forall\mu\in\Spec(A))\ b\le\mu,
 \qquad
 bI-A\succeq0\ \Longleftrightarrow\ (\forall\mu\in\Spec(A))\ \mu\le b.
\]
If $A\succeq0$, then $\Zx^TA\Zx=0$ iff $A\Zx=0$.
If $A\succeq0$, every principal submatrix is positive semidefinite and has
nonnegative determinant. For positive semidefinite real matrices $A,B$ of the same
order, $\tr(AB)\ge0$.}
\end{theorem}
\begin{proof}
\pfsketch\ We will expand the quadratic form in eigenvector coordinates and inspect its nonnegative summands.
\step{psd-signs}{The two positivity equivalences hold.}
\begin{proof}
\step{psd-nonneg}{$A\succeq0$ iff every eigenvalue is nonnegative.}
\begin{proof}
\pf\ \raggedright Apply \code{Matrix.IsHermitian.posSemidef_iff_eigenvalues_nonneg} in Facts~\ref{facts:positive}; Theorem~\ref{thm:multiplicity} identifies the multiset.~\qed
\end{proof}
\step{psd-positive}{$A$ is positive definite iff every eigenvalue is positive.}
\begin{proof}
\pf\ \raggedright Apply \code{Matrix.IsHermitian.posDef_iff_eigenvalues_pos} in Facts~\ref{facts:positive} and the same multiset identification.~\qed
\end{proof}
\qedstep
\begin{proof}
\pf\ \raggedright Use \stepref{psd-nonneg}, \stepref{psd-positive}.~\qed
\end{proof}
\end{proof}
\step{psd-shifts}{$A-bI\succeq0$ iff $b\le\mu$ for every $\mu\in\Spec(A)$, and $bI-A\succeq0$ iff $\mu\le b$ for every $\mu\in\Spec(A)$.}
\begin{proof}
\pf\ \raggedright Apply \code{Matrix.IsHermitian.posSemidef_iff_eigenvalues_nonneg} in Facts~\ref{facts:positive} separately to the symmetric matrices $A-bI$ and $bI-A$; Theorem~\ref{thm:calculus} with $t-b$ and $b-t$ gives their multisets $\ms{\mu-b}$ and $\ms{b-\mu}$.~\qed
\end{proof}
\step{psd-zero}{For $A\succeq0$, $\Zx^TA\Zx=0$ iff $A\Zx=0$.}
\begin{proof}
\pf\ \raggedright Apply \code{Matrix.PosSemidef.dotProduct_mulVec_zero_iff} in Facts~\ref{facts:positive}; real conjugation is trivial.~\qed
\end{proof}
\step{psd-minors}{For $A\succeq0$, every principal submatrix is positive semidefinite with nonnegative determinant.}
\begin{proof}
\step{psd-subpositive}{The principal submatrix is positive semidefinite.}
\begin{proof}
\pf\ \raggedright Apply \code{Matrix.PosSemidef.submatrix} in Facts~\ref{facts:positive} to the inclusion of the principal index set.~\qed
\end{proof}
\step{psd-subdet}{Its determinant is nonnegative.}
\begin{proof}
\pf\ \raggedright Apply \code{Matrix.PosSemidef.det_nonneg} in Facts~\ref{facts:positive} to \stepref{psd-subpositive}.~\qed
\end{proof}
\qedstep
\begin{proof}
\pf\ \raggedright Use \stepref{psd-subpositive}, \stepref{psd-subdet}.~\qed
\end{proof}
\end{proof}
\step{psd-trace}{For $A,B\succeq0$, $\tr(AB)\ge0$.}
\begin{proof}
\step{psd-tracebasis}{\pick\ an orthonormal eigenbasis $(\Zv_i)$ of $A$ with eigenvalues $\lambda_i\ge0$.}
\begin{proof}
\pf\ \raggedright Theorem~\ref{thm:multiplicity} supplies the basis; the spectral criterion in \stepref{psd-signs} makes its eigenvalues nonnegative.~\qed
\end{proof}
\step{psd-tracesum}{$\tr(AB)=\sum_i\lambda_i\Zv_i^TB\Zv_i$.}
\begin{proof}
\pf\ \raggedright Expand $A=\sum_i\lambda_i\Zv_i\Zv_i^T$ and the defining trace sum.~\qed
\end{proof}
\step{psd-traceterms}{Every summand is nonnegative.}
\begin{proof}
\pf\ \raggedright Use \stepref{psd-signs} for $A$ and the defining form inequality for $B$.~\qed
\end{proof}
\qedstep
\begin{proof}
\pf\ \raggedright Use \stepref{psd-tracesum}, \stepref{psd-traceterms}.~\qed
\end{proof}
\end{proof}
\qedstep
\begin{proof}
\pf\ \raggedright Use \stepref{psd-signs}, \stepref{psd-shifts}, \stepref{psd-zero}, \stepref{psd-minors}, \stepref{psd-trace}.~\qed
\end{proof}
\end{proof}

\begin{corollary}[The least-eigenvalue bridge]\label{cor:minus-two}
\lean{SimpleGraph.posSemidef_two_add_adjMatrix_iff, SimpleGraph.neg_two_mem_spectrumMultiset_iff_det_eq_zero, SimpleGraph.leastEigenvalue_eq_neg_two_iff}
\uses{def:graph,def:spectrum,facts:graph,lem:extremes,thm:calculus,thm:psd,thm:trace}
\mbox{}\par\nobreak
\assume{\begin{pfenum}
\item $X$ is a finite simple graph on $U$, and $A=A(X)$.
\item The assertions involving largest or least eigenvalues assume $U\ne\varnothing$.
\end{pfenum}}
\prove{\raggedright \[
 (2I+A)\succeq0\ \Longleftrightarrow\
 (\forall\mu\in\Spec(X))\ -2\le\mu,
 \qquad -2\in\Spec(X)\ \Longleftrightarrow\ \det(2I+A)=0.
\]
If $X$ has a vertex, the first right-hand side is
$\lambda_{\min}(X)\ge-2$. Its least eigenvalue equals $-2$ iff both
the positive semidefiniteness and determinant-zero conditions hold.
Also $2I-A\succeq0$ iff $\lambda_{\max}(X)\le2$ for a nonempty graph.}
\end{corollary}
\begin{proof}
\pfsketch\ We will shift every eigenvalue by two and then take the minimum or maximum of the finite list.
\step{minus-shift}{The two assertions that do not require a vertex hold.}
\begin{proof}
\step{minus-psd}{$2I+A\succeq0$ iff every eigenvalue of $A$ is at least $-2$.}
\begin{proof}
\pf\ \raggedright Use adjacency symmetry (Facts~\ref{facts:graph}) and Theorem~\ref{thm:psd} with $b=-2$.~\qed
\end{proof}
\step{minus-det}{$\det(2I+A)=\prod_i(2+\lambda_i)$.}
\begin{proof}
\pf\ \raggedright Apply Theorems~\ref{thm:calculus} and \ref{thm:trace} to $2+t$.~\qed
\end{proof}
\step{minus-root}{The product is zero iff $-2$ occurs in $\Spec(A)$.}
\begin{proof}
\pf\ \raggedright A finite real product vanishes iff one factor vanishes; use \stepref{minus-det}.~\qed
\end{proof}
\qedstep
\begin{proof}
\pf\ \raggedright Use \stepref{minus-psd}, \stepref{minus-root}.~\qed
\end{proof}
\end{proof}
\step{minus-extreme}{The remaining equivalences hold when $U\ne\varnothing$.}
\begin{proof}
\step{minus-min}{All entries are at least $-2$ iff the minimum is at least $-2$.}
\begin{proof}
\pf\ \raggedright Use the least-element property in Lemma~\ref{lem:extremes}.~\qed
\end{proof}
\step{minus-eq}{The minimum is $-2$ iff all entries are at least $-2$ and $-2$ occurs.}
\begin{proof}
\pf\ \raggedright A finite nonempty list contains its minimum; apply Lemma~\ref{lem:extremes}.~\qed
\end{proof}
\step{minus-max}{$2I-A\succeq0$ iff the maximum is at most $2$.}
\begin{proof}
\pf\ \raggedright Use the upper shift of Theorem~\ref{thm:psd} and Lemma~\ref{lem:extremes}.~\qed
\end{proof}
\qedstep
\begin{proof}
\pf\ \raggedright Use \stepref{minus-min}, \stepref{minus-eq}, \stepref{minus-max}.~\qed
\end{proof}
\end{proof}
\step{minus-empty}{When $U=\varnothing$, both sides of the determinant-zero equivalence are false.}
\begin{proof}
\pf\ \raggedright The empty determinant is $1$ and the empty spectrum has no member, by Definition~\ref{def:spectrum}.~\qed
\end{proof}
\qedstep
\begin{proof}
\pf\ \raggedright Use \stepref{minus-shift}, \stepref{minus-extreme}, \stepref{minus-empty}.~\qed
\end{proof}
\end{proof}

\begin{theorem}[Gram matrices with rank control]\label{thm:gram}
\lean{Matrix.IsHermitian.exists_gram_rank, Matrix.rank_gram_eq_finrank_span, SimpleGraph.exists_representation_iff_posSemidef, SimpleGraph.multiplicity_two_ge_of_gram_representation}
\uses{facts:coord,facts:graph,facts:order,facts:positive,facts:rank,lem:extremes,thm:calculus,thm:multiplicity,thm:psd,thm:trace}
\mbox{}\par\nobreak
\assume{\begin{pfenum}
\item $U$ is finite and $G:U\times U\to\re$ is symmetric, with $\rk G=\rho$.
\item $B$ is any real matrix on finite row and column sets. Vector families below are finite and lie in a finite-dimensional real inner product space.
\item For the graph assertions, $X$ is a finite simple graph with $n$ vertices; the lower-shift occurrence assumes a representation in dimension $\ell<n$, while the upper-shift rank bound allows any $\ell$.
\end{pfenum}}
\prove{\raggedright A real symmetric matrix $G$ is positive semidefinite iff it is a Gram
matrix. If $G\succeq0$ and $\rk G=\rho$, it has a Gram representation in $\re^\rho$
whose vectors span that space. For any real matrix $B$,
\[
 \ker(B^TB)=\ker B,\qquad \rk(B^TB)=\rk B=\rk(BB^T).
\]
Thus the rank of the Gram matrix of a finite vector family equals the
dimension of its span. In particular $2I+A(X)\succeq0$ iff $X$ has a
representation by vectors of squared norm $2$ with inner product $1$
on edges and $0$ on distinct nonedges. If such vectors lie in dimension
$\ell<n$, then $-2$ occurs, with multiplicity at least $n-\ell$.
For a Gram representation of $2I-A(X)$ in dimension $\ell$, every
eigenvalue is at most $2$ and $\mult_X(2)\ge n-\ell$.
If $n>\ell$, the largest eigenvalue is $2$.}
\end{theorem}
\begin{proof}
\pfsketch\ We will take square roots of positive eigendirections, then identify the kernel and image of the coordinate matrix.
\step{gram-existence}{Positive semidefiniteness is equivalent to some Gram representation.}
\begin{proof}
\pf\ \raggedright Apply \code{Matrix.posSemidef_iff_eq_sum_vecMulVec} in Facts~\ref{facts:positive}; the vectors in its sum are the rows of a factor matrix, whose columns have the asserted Gram matrix.~\qed
\end{proof}
\step{gram-factor}{If $G\succeq0$, it equals $B^TB$ for a $\rho$-row matrix whose columns span $\re^\rho$.}
\begin{proof}
\step{gram-eigenbasis}{\pick\ an orthonormal eigenbasis $(\Zv_i)$ for $G$, with eigenvalues $\lambda_i$.}
\begin{proof}
\pf\ \raggedright Theorem~\ref{thm:multiplicity}.~\qed
\end{proof}
\step{gram-positive}{Exactly $\rho$ eigenvalues of $G$ are positive.}
\begin{proof}
\pf\ \raggedright Use Theorems~\ref{thm:psd} and \ref{thm:trace}.~\qed
\end{proof}
\step{gram-construct}{For those indices set $B(i,u)=\sqrt{\lambda_i}\Zv_i(u)$; then $B^TB=G$.}
\begin{proof}
\pf\ \raggedright Expand the product; $(\sqrt{\lambda_i})^2=\lambda_i$ by Facts~\ref{facts:order}, and the omitted eigenvalues are zero.~\qed
\end{proof}
\step{gram-rows}{The rows of $B$ are independent.}
\begin{proof}
\pf\ \raggedright Divide a zero linear combination by the nonzero square-root coefficients and use independence of $(\Zv_i)$.~\qed
\end{proof}
\step{gram-span}{The columns span $\re^\rho$.}
\begin{proof}
\pf\ \raggedright Row rank equals column rank by Facts~\ref{facts:rank}; use \stepref{gram-rows}.~\qed
\end{proof}
\qedstep
\begin{proof}
\pf\ \raggedright Use \stepref{gram-construct}, \stepref{gram-span}.~\qed
\end{proof}
\end{proof}
\step{gram-kernel}{Every $B^TB$ is positive semidefinite and $\ker(B^TB)=\ker B$.}
\begin{proof}
\step{gram-square}{$\Zx^TB^TB\Zx=\norm{B\Zx}^2$.}
\begin{proof}
\pf\ \raggedright Expand matrix multiplication and the Euclidean inner product.~\qed
\end{proof}
\step{gram-nonneg}{$B^TB\succeq0$.}
\begin{proof}
\pf\ \raggedright It is symmetric and the form in \stepref{gram-square} is nonnegative; apply Facts~\ref{facts:positive}.~\qed
\end{proof}
\step{gram-kerforward}{$B^TB\Zx=0$ implies $B\Zx=0$.}
\begin{proof}
\pf\ \raggedright In \stepref{gram-square} the left side vanishes, so a sum of coordinate squares is zero (Facts~\ref{facts:order}).~\qed
\end{proof}
\step{gram-kerback}{$B\Zx=0$ implies $B^TB\Zx=0$.}
\begin{proof}
\pf\ \raggedright Apply $B^T$ to the assumed zero vector.~\qed
\end{proof}
\qedstep
\begin{proof}
\pf\ \raggedright Use \stepref{gram-nonneg}, \stepref{gram-kerforward}, \stepref{gram-kerback}.~\qed
\end{proof}
\end{proof}
\step{gram-ranks}{$\rk(B^TB)=\rk B=\rk(BB^T)$; a Gram rank is the dimension of the vector span.}
\begin{proof}
\step{gram-ranknull}{$\rk(B^TB)=\rk B$.}
\begin{proof}
\pf\ \raggedright Apply \code{Matrix.rank_transpose_mul_self} in Facts~\ref{facts:rank} to the real rectangular matrix $B$.~\qed
\end{proof}
\step{gram-other}{$\rk(BB^T)=\rk B$.}
\begin{proof}
\pf\ \raggedright Apply \code{Matrix.rank_self_mul_transpose} in Facts~\ref{facts:rank} to $B$.~\qed
\end{proof}
\step{gram-family}{The coordinate matrix of a family has image equal to its span.}
\begin{proof}
\pf\ \raggedright The matrix image consists exactly of linear combinations of its columns.~\qed
\end{proof}
\step{gram-familyrank}{Its Gram rank equals the dimension of that span.}
\begin{proof}
\pf\ \raggedright Apply \code{Matrix.rank_transpose_mul_self} to the coordinate matrix of the family in an orthonormal basis (Facts~\ref{facts:coord}); its image is the vector span.~\qed
\end{proof}
\qedstep
\begin{proof}
\pf\ \raggedright Use \stepref{gram-ranknull}, \stepref{gram-other}, \stepref{gram-familyrank}.~\qed
\end{proof}
\end{proof}
\step{gram-graph}{The graph representation criterion, the multiplicity bounds for both $2I+A$ and $2I-A$, and the stated occurrences hold.}
\begin{proof}
\step{gram-entries}{$2I+A(X)$ has diagonal $2$, edge entries $1$ and other off-diagonal entries $0$.}
\begin{proof}
\pf\ \raggedright Use the definition of adjacency matrix in Facts~\ref{facts:graph}.~\qed
\end{proof}
\step{gram-represent}{Its positivity is equivalent to the specified Gram representation.}
\begin{proof}
\pf\ \raggedright Apply the equivalence in \code{Matrix.posSemidef_iff_eq_sum_vecMulVec} (Facts~\ref{facts:positive}) directly to $2I+A$; its entries in \stepref{gram-entries} are exactly the required inner products.~\qed
\end{proof}
\step{gram-boundrank}{$\rk(2I+A)\le\ell$.}
\begin{proof}
\pf\ \raggedright Choose orthonormal coordinates for the representing family. Its coordinate matrix has at most $\ell$ rows; \code{Matrix.rank_transpose_mul_self} in Facts~\ref{facts:rank} bounds the Gram rank by $\ell$.~\qed
\end{proof}
\step{gram-minusmult}{$\mult_X(-2)=n-\rk(2I+A)\ge n-\ell$.}
\begin{proof}
\pf\ \raggedright Apply Theorems~\ref{thm:calculus} and \ref{thm:trace} to the shift, then \stepref{gram-boundrank}.~\qed
\end{proof}
\step{gram-upperrank}{For a representation of $2I-A$ in dimension $\ell$, $\mult_X(2)\ge n-\ell$.}
\begin{proof}
\pf\ \raggedright The coordinate Gram matrix has rank at most $\ell$ by \code{Matrix.rank_transpose_mul_self} (Facts~\ref{facts:rank}); Theorems~\ref{thm:calculus} and \ref{thm:trace} identify its nullity with $\mult_X(2)$.~\qed
\end{proof}
\step{gram-uppertop}{Every eigenvalue is at most $2$, with maximum $2$ if $n>\ell$.}
\begin{proof}
\pf\ \raggedright The Gram form is positive semidefinite, so Theorem~\ref{thm:psd} gives the bound. The positive multiplicity in \stepref{gram-upperrank} and Lemma~\ref{lem:extremes} give equality.~\qed
\end{proof}
\qedstep
\begin{proof}
\pf\ \raggedright Use \stepref{gram-represent}, \stepref{gram-minusmult}, \stepref{gram-upperrank}, \stepref{gram-uppertop}.~\qed
\end{proof}
\end{proof}
\qedstep
\begin{proof}
\pf\ \raggedright Use \stepref{gram-existence}, \stepref{gram-factor}, \stepref{gram-kernel}, \stepref{gram-ranks}, \stepref{gram-graph}.~\qed
\end{proof}
\end{proof}

\begin{lemma}[Trace and rank inequality]\label{lem:trace-rank-bound}
\lean{Matrix.IsHermitian.trace_sq_le_rank_mul_trace_sq}
\uses{facts:order,thm:multiplicity,thm:trace}
\mbox{}\par\nobreak
\assume{\begin{pfenum}
\item $G$ is a real positive semidefinite matrix on a finite index set, of rank $\rho$.
\end{pfenum}}
\prove{\raggedright $(\tr G)^2\le\rho\tr(G^2)$, with equality iff all nonzero eigenvalues
are equal; equality includes $G=0$.}
\end{lemma}
\begin{proof}
\pfsketch\ We will write $\rho$ times the sum of squares of the nonzero eigenvalues, minus the square of their sum, as a sum of squared differences.
\pflet{$\eta_1,\ldots,\eta_\rho$ are the nonzero members of $\Spec(G)$, with repetition; there are $\rho$ of them by Theorem~\ref{thm:trace}.}
\step{tracerank-expand}{$\rho\sum_i\eta_i^2-(\sum_i\eta_i)^2=\sum_{i<j}(\eta_i-\eta_j)^2$.}
\begin{proof}
\pf\ \raggedright Expand the square: each $\eta_i^2$ occurs $\rho-1$ times on the right and each pair contributes $-2\eta_i\eta_j$.~\qed
\end{proof}
\step{tracerank-equal}{The right side is nonnegative and vanishes iff all $\eta_i$ are equal.}
\begin{proof}
\pf\ \raggedright Apply the nonnegative zero-sum criterion in Facts~\ref{facts:order} to \stepref{tracerank-expand}.~\qed
\end{proof}
\step{tracerank-traces}{$\sum_i\eta_i^2=\tr(G^2)$ and $\sum_i\eta_i=\tr G$; when $\rho=0$, $G=0$.}
\begin{proof}
\pf\ \raggedright Theorem~\ref{thm:trace}, since zero eigenvalues add nothing; for $\rho=0$ every eigenvalue is zero, and Theorem~\ref{thm:multiplicity} gives $G=Q\,0\,Q^T=0$.~\qed
\end{proof}
\qedstep
\begin{proof}
\pf\ \raggedright Substitute \stepref{tracerank-traces} in \stepref{tracerank-expand} and apply \stepref{tracerank-equal}; for $\rho=0$ both sides are zero.~\qed
\end{proof}
\end{proof}

\begin{theorem}[Trace bounds and tight frames]\label{thm:trace-bound}
\lean{Matrix.spectrumMultiset_gram_of_tightFrame}
\uses{lem:trace-rank-bound,thm:calculus,thm:gram,thm:trace}
\mbox{}\par\nobreak
\assume{\begin{pfenum}
\item $G:U\times U\to\re$ is positive semidefinite on a finite set, with rank $\rho$.
\item For the relative bound, $n$ unit vectors lie in $\re^\ell$, their distinct-pair inner products have absolute value $\alpha$, $0\le\alpha\le1$, and $\ell\alpha^2<1$.
\item For the frame assertion, $B$ is real $\ell\times n$, $\ell>0$, $b>0$, and $BB^T=bI$.
\end{pfenum}}
\prove{\raggedright $(\tr G)^2\le\rho\tr(G^2)$; equality holds iff all its nonzero
eigenvalues are equal, with equality also allowed when $G=0$.
If $n$ unit vectors in $\re^\ell$ have pairwise inner products of
absolute value $\alpha$, where $0\le\alpha\le1$ and
$\ell\alpha^2<1$, then
\[
 n\le\frac{\ell(1-\alpha^2)}{1-\ell\alpha^2}.
\]
If $B$ is a real $\ell\times n$ matrix, $\ell>0$, and
$BB^T=bI$ with $b>0$, then $G=B^TB$ has spectrum
$b^{\ell},0^{n-\ell}$. In particular, with $C=G-2I$,
\[
 (C+2I)(C-(b-2)I)=0.
\]}
\end{theorem}
\begin{proof}
\pfsketch\ We will expand a sum of squared differences, apply it to a Gram matrix, and compute the frame polynomial.
\step{bound-squares}{$(\tr G)^2\le\rho\tr(G^2)$ with the stated equality criterion.}
\begin{proof}
\pf\ \raggedright Apply Lemma~\ref{lem:trace-rank-bound} to $G$.~\qed
\end{proof}
\step{bound-relative}{$n\le\ell(1-\alpha^2)/(1-\ell\alpha^2)$.}
\begin{proof}
\step{bound-gramtrace}{The unit-vector Gram matrix has $\tr G=n$ and $\tr(G^2)=n+\alpha^2n(n-1)$.}
\begin{proof}
\pf\ \raggedright Count its $n$ diagonal squares $1$ and $n(n-1)$ off-diagonal squares $\alpha^2$ in the defining trace sum.~\qed
\end{proof}
\step{bound-rho}{$\rho\le\ell$.}
\begin{proof}
\pf\ \raggedright Apply Theorem~\ref{thm:gram} to the given family.~\qed
\end{proof}
\step{bound-positive-n}{\case{$n>0$.}}
\begin{proof}
\step{bound-inequality}{$n(1-\ell\alpha^2)\le\ell(1-\alpha^2)$.}
\begin{proof}
\pf\ \raggedright Apply Lemma~\ref{lem:trace-rank-bound} to the Gram matrix of the given unit vectors. Substitute \stepref{bound-gramtrace} and \stepref{bound-rho}; divide by $n>0$ and collect terms.~\qed
\end{proof}
\qedstep
\begin{proof}
\pf\ \raggedright Divide \stepref{bound-inequality} by the positive number $1-\ell\alpha^2$.~\qed
\end{proof}
\end{proof}
\step{bound-zero-n}{\case{$n=0$.}}
\begin{proof}
\pf\ \raggedright The numerator $\ell(1-\alpha^2)$ is nonnegative and the denominator is positive.~\qed
\end{proof}
\qedstep
\begin{proof}
\pf\ \raggedright The cases $n>0$ and $n=0$ exhaust $n\in\nn$.~\qed
\end{proof}
\end{proof}
\step{bound-frame}{The frame spectrum and polynomial identity hold.}
\begin{proof}
\step{bound-full}{$\rk B=\ell$ and $\rk(B^TB)=\ell\le n$.}
\begin{proof}
\pf\ \raggedright Apply Theorem~\ref{thm:gram} to $BB^T=bI$ with $b\ne0$.~\qed
\end{proof}
\step{bound-poly}{For $G=B^TB$, $G^2=B^T(BB^T)B=bG$.}
\begin{proof}
\pf\ \raggedright Associativity and the frame hypothesis.~\qed
\end{proof}
\step{bound-roots}{Every eigenvalue of $G$ is $0$ or $b$.}
\begin{proof}
\pf\ \raggedright Apply the zero-polynomial criterion of Theorem~\ref{thm:calculus} to \stepref{bound-poly}.~\qed
\end{proof}
\step{bound-count}{There are $\ell$ copies of $b$ and $n-\ell$ zeros.}
\begin{proof}
\pf\ \raggedright Count the nonzero eigenvalues using \stepref{bound-full}, \stepref{bound-roots} and Theorem~\ref{thm:trace}.~\qed
\end{proof}
\step{bound-C}{$(C+2I)(C-(b-2)I)=0$ for $C=G-2I$.}
\begin{proof}
\pf\ \raggedright Substitute $G=C+2I$ in $G(G-bI)=0$ from \stepref{bound-poly}.~\qed
\end{proof}
\qedstep
\begin{proof}
\pf\ \raggedright Use \stepref{bound-count}, \stepref{bound-C}.~\qed
\end{proof}
\end{proof}
\qedstep
\begin{proof}
\pf\ \raggedright Use \stepref{bound-squares}, \stepref{bound-relative}, \stepref{bound-frame}.~\qed
\end{proof}
\end{proof}

\section{Rayleigh quotients and symmetric Perron--Frobenius}

\begin{theorem}[Rayleigh extrema and equality]\label{thm:rayleigh}
\lean{Matrix.IsHermitian.rayleigh_le_largestEigenvalue, Matrix.IsHermitian.rayleigh_eq_largestEigenvalue_iff, Matrix.IsHermitian.leastEigenvalue_le_rayleigh}
\uses{def:extremes,facts:coord,facts:order,facts:rayleigh,lem:extremes,thm:multiplicity}
\mbox{}\par\nobreak
\assume{\begin{pfenum}
\item $U$ is finite and nonempty, $|U|=n$, $A:U\times U\to\re$ is symmetric, and $\Zx\in\re^U\setminus\{0\}$.
\item For the coordinate formulas, $(\Zv_i)_{i\in U}$ is any orthonormal eigenbasis of $A$, $A\Zv_i=\lambda_i\Zv_i$, and $\Zx=\sum_ib_i\Zv_i$.
\item For the orthogonality assertion, $\Zy,\Zz\in\re^U$ satisfy $A\Zy=s\Zy$, $A\Zz=t\Zz$, where $s,t\in\re$ and $s\ne t$.
\end{pfenum}}
\prove{\raggedright For
$\Zx\ne0$, put $R_A(\Zx)=\Zx^TA\Zx/\norm{\Zx}^2$. Then
\[
 \lambda_{\min}(A)\le R_A(\Zx)\le\lambda_{\max}(A).
\]
In any orthonormal eigenbasis, writing $\Zx=\sum_i b_i\Zv_i$ gives
\[
 R_A(\Zx)=\frac{\sum_i\lambda_i b_i^2}{\sum_i b_i^2},\qquad
 \sum_i b_i^2=\norm{\Zx}^2>0,
\]
\[
 \begin{aligned}
 \lambda_{\max}\norm{\Zx}^2-\Zx^TA\Zx
   &=\sum_i(\lambda_{\max}-\lambda_i)b_i^2,\\
 \Zx^TA\Zx-\lambda_{\min}\norm{\Zx}^2
   &=\sum_i(\lambda_i-\lambda_{\min})b_i^2.
 \end{aligned}
\]
Both bounds are attained, including on the unit sphere. Equality on
the right holds iff $A\Zx=\lambda_{\max}(A)\Zx$, and equality on
the left holds iff $A\Zx=\lambda_{\min}(A)\Zx$. Thus the extremes
are the supremum and infimum of the Rayleigh quotient on nonzero vectors.
Eigenvectors of a symmetric matrix for distinct eigenvalues are orthogonal.}
\end{theorem}
\begin{proof}
\pfsketch\ We will express the two gaps as sums of nonnegative terms and identify where equality can occur. Such a basis exists by Theorem~\ref{thm:multiplicity}, and the bounds for $\Zx$ do not depend on the choice.
\step{ray-form}{$R_A(\Zx)=\sum_i\lambda_i b_i^2/\sum_i b_i^2$, with positive denominator $\sum_ib_i^2=\norm{\Zx}^2$.}
\begin{proof}
\step{ray-numerator}{$\Zx^TA\Zx=\sum_i\lambda_i b_i^2$.}
\begin{proof}
\pf\ \raggedright Expand $A\Zx=\sum_ib_i\lambda_i\Zv_i$ by \asm[2]{0}, and use orthonormality.~\qed
\end{proof}
\step{ray-denominator}{$\norm{\Zx}^2=\sum_i b_i^2>0$.}
\begin{proof}
\pf\ \raggedright Orthonormality gives the sum; it is positive since $\Zx\ne0$ by \asm[1]{0}.~\qed
\end{proof}
\qedstep
\begin{proof}
\pf\ \raggedright Divide \stepref{ray-numerator} by \stepref{ray-denominator}.~\qed
\end{proof}
\end{proof}
\step{ray-members}{Every $\lambda_i$ belongs to $\Spec(A)$, so $\lambda_{\min}\le\lambda_i\le\lambda_{\max}$.}
\begin{proof}
\pf\ \raggedright $\Zv_i$ is a nonzero $\lambda_i$-eigenvector; apply the membership criterion of Theorem~\ref{thm:multiplicity} and the bounds of Lemma~\ref{lem:extremes}.~\qed
\end{proof}
\step{ray-upper}{The upper gap is $\sum_i(\lambda_{\max}-\lambda_i)b_i^2\ge0$; thus $R_A(\Zx)\le\lambda_{\max}$, with equality iff $A\Zx=\lambda_{\max}\Zx$.}
\begin{proof}
\step{ray-uppergap}{$\lambda_{\max}\norm{\Zx}^2-\Zx^TA\Zx=\sum_i(\lambda_{\max}-\lambda_i)b_i^2$.}
\begin{proof}
\pf\ \raggedright Subtract the two expressions in \stepref{ray-form}.~\qed
\end{proof}
\step{ray-uppert}{Every summand is nonnegative.}
\begin{proof}
\pf\ \raggedright Use \stepref{ray-members}.~\qed
\end{proof}
\step{ray-upperzero}{The sum is zero iff $b_i=0$ whenever $\lambda_i<\lambda_{\max}$.}
\begin{proof}
\pf\ \raggedright Apply Facts~\ref{facts:order} to \stepref{ray-uppert} and use that a positive factor times a square vanishes iff the coordinate does.~\qed
\end{proof}
\step{ray-uppereigen}{That coefficient condition is the equation $A\Zx=\lambda_{\max}\Zx$.}
\begin{proof}
\pf\ \raggedright Compare the coefficients of $A\Zx=\sum_i\lambda_ib_i\Zv_i$ and $\lambda_{\max}\Zx$ in the basis of \asm[2]{0}.~\qed
\end{proof}
\qedstep
\begin{proof}
\pf\ \raggedright Divide the nonnegative gap by the positive denominator in \stepref{ray-form}; the zero criterion is \stepref{ray-upperzero}, \stepref{ray-uppereigen}.~\qed
\end{proof}
\end{proof}
\step{ray-lower}{The lower gap is $\sum_i(\lambda_i-\lambda_{\min})b_i^2\ge0$; thus $\lambda_{\min}\le R_A(\Zx)$, with equality iff $A\Zx=\lambda_{\min}\Zx$.}
\begin{proof}
\step{ray-lowergap}{$\Zx^TA\Zx-\lambda_{\min}\norm{\Zx}^2=\sum_i(\lambda_i-\lambda_{\min})b_i^2$.}
\begin{proof}
\pf\ \raggedright Subtract using \stepref{ray-form}.~\qed
\end{proof}
\step{ray-lowert}{Every summand is nonnegative.}
\begin{proof}
\pf\ \raggedright Use \stepref{ray-members}.~\qed
\end{proof}
\step{ray-lowerzero}{The sum is zero iff $b_i=0$ whenever $\lambda_i>\lambda_{\min}$.}
\begin{proof}
\pf\ \raggedright Use the zero-sum criterion in Facts~\ref{facts:order} and \stepref{ray-lowert}.~\qed
\end{proof}
\step{ray-lowereigen}{That condition is $A\Zx=\lambda_{\min}\Zx$.}
\begin{proof}
\pf\ \raggedright Compare coefficients in the basis of \asm[2]{0}.~\qed
\end{proof}
\qedstep
\begin{proof}
\pf\ \raggedright Divide the gap by \stepref{ray-form}; use \stepref{ray-lowerzero} and \stepref{ray-lowereigen} for equality.~\qed
\end{proof}
\end{proof}
\step{ray-attain}{Both bounds are attained on unit vectors and are the supremum and infimum on nonzero vectors.}
\begin{proof}
\step{ray-witness}{Unit eigenvectors attain each extreme.}
\begin{proof}
\pf\ \raggedright The sorted basis vectors $\Zw_0$ and $\Zw_{n-1}$ of Theorem~\ref{thm:multiplicity} have eigenvalues $\nu_0=\lambda_{\max}$ and $\nu_{n-1}=\lambda_{\min}$ by Definition~\ref{def:extremes}; their quotients are these eigenvalues.~\qed
\end{proof}
\step{ray-uniform}{For every nonzero $\Zw$, $\lambda_{\min}\le R_A(\Zw)\le\lambda_{\max}$.}
\begin{proof}
\pf\ \raggedright Expand $\Zw=\sum_i c_i\Zv_i$ in the basis of \asm[2]{0}. Orthonormality and the eigenvector equations make its quotient $\sum_i\lambda_i c_i^2/\sum_i c_i^2$, with positive denominator since $\Zw\ne0$, and \stepref{ray-members} bounds this weighted average.~\qed
\end{proof}
\step{ray-sup-library}{The nonzero-vector supremum is an eigenvalue of $A.\mathrm{toEuclideanLin}$.}
\begin{proof}
\pf\ \raggedright Apply \code{LinearMap.IsSymmetric.hasEigenvalue_iSup_of_finiteDimensional} in Facts~\ref{facts:rayleigh}, using matrix symmetry from Facts~\ref{facts:coord}; $n>0$ makes the space nontrivial.~\qed
\end{proof}
\step{ray-inf-library}{The nonzero-vector infimum is also an eigenvalue.}
\begin{proof}
\pf\ \raggedright Apply \code{LinearMap.IsSymmetric.hasEigenvalue_iInf_of_finiteDimensional} in Facts~\ref{facts:rayleigh} to the same symmetric map.~\qed
\end{proof}
\step{ray-library-identify}{These two eigenvalues are $\lambda_{\max}$ and $\lambda_{\min}$.}
\begin{proof}
\pf\ \raggedright The Euclidean inner product in Facts~\ref{facts:coord} is the coordinate sum, so the quotient in \stepref{ray-sup-library} and \stepref{ray-inf-library} is $R_A$. The bounds in \stepref{ray-uniform} are attained by \stepref{ray-witness}, so the supremum is $\lambda_{\max}$ and the infimum is $\lambda_{\min}$.~\qed
\end{proof}
\qedstep
\begin{proof}
\pf\ \raggedright Use \stepref{ray-witness}, \stepref{ray-library-identify}.~\qed
\end{proof}
\end{proof}
\step{ray-orth}{Eigenvectors at distinct eigenvalues are orthogonal.}
\begin{proof}
\step{ray-symmetric}{$s\ip{\Zy}{\Zz}=t\ip{\Zy}{\Zz}$.}
\begin{proof}
\step{ray-left-eigen}{$s\ip{\Zy}{\Zz}=\ip{A\Zy}{\Zz}$.}
\begin{proof}
\pf\ \raggedright Substitute $A\Zy=s\Zy$ from \asm[3]{0} in the first inner-product argument.~\qed
\end{proof}
\step{ray-adjoint}{$\ip{A\Zy}{\Zz}=\ip{\Zy}{A\Zz}$.}
\begin{proof}
\pf\ \raggedright Expand both coordinate sums and use $A(u,v)=A(v,u)$.~\qed
\end{proof}
\step{ray-right-eigen}{$\ip{\Zy}{A\Zz}=t\ip{\Zy}{\Zz}$.}
\begin{proof}
\pf\ \raggedright Substitute $A\Zz=t\Zz$ in the second argument.~\qed
\end{proof}
\qedstep
\begin{proof}
\pf\ \raggedright Transitivity of the three equalities.~\qed
\end{proof}
\end{proof}
\step{ray-cancel}{$\ip{\Zy}{\Zz}=0$.}
\begin{proof}
\pf\ \raggedright Cancel $s-t\ne0$ in \stepref{ray-symmetric}.~\qed
\end{proof}
\qedstep
\begin{proof}
\pf\ \raggedright Use \stepref{ray-cancel}.~\qed
\end{proof}
\end{proof}
\qedstep
\begin{proof}
\pf\ \raggedright Use \stepref{ray-form}, \stepref{ray-upper}, \stepref{ray-lower}, \stepref{ray-attain}, \stepref{ray-orth}.~\qed
\end{proof}
\end{proof}

\begin{lemma}[Rayleigh subspace witnesses]\label{lem:rayleigh-subspaces}
\lean{Matrix.IsHermitian.rayleigh_head_tail, Matrix.IsHermitian.exists_rayleigh_le_of_finrank, Matrix.IsHermitian.exists_le_rayleigh_of_finrank}
\uses{facts:dim,thm:multiplicity,thm:rayleigh}
\mbox{}\par\nobreak
\assume{\begin{pfenum}
\item $U$ is finite, $|U|=n>0$, and $A:U\times U\to\re$ is symmetric.
\item $(\Zw_j)_{j\in\operatorname{Fin}(n)}$ and $\nu_0\ge\cdots\ge\nu_{n-1}$ are the sorted eigenbasis and eigenvalues of Theorem~\ref{thm:multiplicity}, and $0\le i<n$.
\item The head is $\spanop\{\Zw_0,\ldots,\Zw_i\}$ and the tail is $\spanop\{\Zw_i,\ldots,\Zw_{n-1}\}$.
\end{pfenum}}
\prove{\raggedright For $0\le i<n$, the head span of directions $0,\ldots,i$ has dimension
$i+1$ and quotient at least $\nu_i$ on nonzero vectors. The tail span
of directions $i,\ldots,n-1$ has dimension $n-i$ and quotient at most
$\nu_i$. Every subspace of dimension $i+1$ contains a nonzero vector
with quotient at most $\nu_i$; every subspace of dimension $n-i$
contains one with quotient at least $\nu_i$.}
\end{lemma}
\begin{proof}
\pfsketch\ We will bound the quotient on the head and on the tail by the weighted-average formula of Theorem~\ref{thm:rayleigh}, and we will find both witnesses in intersections, by the dimension formula for a sum of subspaces.
\step{witness-dim}{\assume{$P,Q$ are subspaces of $\re^U$ with $\dim P+\dim Q>n$.}\prove{There is $\Zx\ne0$ in $P\cap Q$.}}
\begin{proof}
\step{witness-dimensioneq}{$\dim(P+Q)+\dim(P\cap Q)=\dim P+\dim Q$.}
\begin{proof}
\pf\ \raggedright Apply \code{Submodule.finrank_sup_add_finrank_inf_eq} in Facts~\ref{facts:dim}.~\qed
\end{proof}
\step{witness-intersectionpos}{$\dim(P\cap Q)>0$.}
\begin{proof}
\pf\ \raggedright Use $\dim(P+Q)\le n$ (Facts~\ref{facts:dim}) and \asm{1} in \stepref{witness-dimensioneq}.~\qed
\end{proof}
\step{witness-witness}{There is $\Zx\ne0$ in $P\cap Q$.}
\begin{proof}
\pf\ \raggedright A positive-dimensional subspace has a nonempty basis, by Facts~\ref{facts:dim}.~\qed
\end{proof}
\qedstep
\begin{proof}
\pf\ \raggedright Use \stepref{witness-witness}.~\qed
\end{proof}
\end{proof}
\step{witness-tail}{The tail has dimension $n-i$, and $R_A(\Zx)\le\nu_i$ for every nonzero $\Zx$ in it.}
\begin{proof}
\pf\ \raggedright The $\Zw_j$ are independent, by Theorem~\ref{thm:multiplicity}. In the coordinate formula of Theorem~\ref{thm:rayleigh} for the basis $(\Zw_j)$, only indices $j\ge i$ have nonzero coefficients, and $\nu_j\le\nu_i$ for them.~\qed
\end{proof}
\step{witness-head}{The head has dimension $i+1$, and $R_A(\Zx)\ge\nu_i$ for every nonzero $\Zx$ in it.}
\begin{proof}
\pf\ \raggedright In the same coordinate formula only indices $j\le i$ have nonzero coefficients, and $\nu_j\ge\nu_i$ for them; the $\Zw_j$ are independent.~\qed
\end{proof}
\step{witness-low-witness}{Every subspace of dimension $i+1$ contains a nonzero $\Zx$ with $R_A(\Zx)\le\nu_i$.}
\begin{proof}
\pf\ \raggedright Apply \stepref{witness-dim} to that subspace and the tail, since $(i+1)+(n-i)=n+1$; use \stepref{witness-tail}.~\qed
\end{proof}
\step{witness-high-witness}{Every subspace of dimension $n-i$ contains a nonzero $\Zx$ with $R_A(\Zx)\ge\nu_i$.}
\begin{proof}
\pf\ \raggedright Apply \stepref{witness-dim} to that subspace and the head, since $(n-i)+(i+1)=n+1$; use \stepref{witness-head}.~\qed
\end{proof}
\qedstep
\begin{proof}
\pf\ \raggedright Use \stepref{witness-tail}, \stepref{witness-head}, \stepref{witness-low-witness}, \stepref{witness-high-witness}.~\qed
\end{proof}
\end{proof}

\begin{theorem}[Cauchy interlacing]\label{thm:interlacing}
\lean{Matrix.IsHermitian.eigenvalues_principal_submatrix_le, Matrix.IsHermitian.le_eigenvalues_principal_submatrix}
\uses{facts:dim,lem:extremes,lem:rayleigh-subspaces,thm:multiplicity}
\mbox{}\par\nobreak
\assume{\begin{pfenum}
\item $U,W$ are finite with $|U|=n$, $|W|=m$ and $1\le m\le n$.
\item $A:U\times U\to\re$ is symmetric, $\phi:W\hookrightarrow U$ is injective, and $B(u,v)=A(\phi(u),\phi(v))$.
\item $\lambda_i=\nu_i(A)$ and $\beta_i=\nu_i(B)$ are the decreasing eigenvalue lists of the sorted eigenbases of $A$ and $B$ in Theorem~\ref{thm:multiplicity}; by that theorem they list the two spectra with multiplicities.
\end{pfenum}}
\prove{\raggedright \[
 \lambda_i\ge\beta_i\ge\lambda_{i+n-m}\qquad(0\le i<m).
\]
In particular the largest eigenvalue cannot increase, and the least
eigenvalue cannot decrease, on taking a nonempty principal submatrix.}
\end{theorem}
\begin{proof}
\pfsketch\ We will use the two subspace comparisons in the ordinary proof. Their nonzero intersection witnesses give the inequalities directly, without choosing extrema over a Grassmannian.
\step{int-symmetric}{$B$ is symmetric.}
\begin{proof}
\pf\ \raggedright For every $u,v$, $B(u,v)=A(\phi(u),\phi(v))=A(\phi(v),\phi(u))=B(v,u)$ by \asm[2]{0}.~\qed
\end{proof}
\step{int-extension}{Extension by zero $\iota:\re^W\to\re^U$ preserves dimension, norm and the quadratic form.}
\begin{proof}
\step{int-inj}{$\iota$ is linear and injective.}
\begin{proof}
\pf\ \raggedright Restricting back to $\phi(W)$ is a left inverse, by the definition of extension by zero.~\qed
\end{proof}
\step{int-norm}{$\norm{\iota\Zx}^2=\norm{\Zx}^2$.}
\begin{proof}
\pf\ \raggedright Only the coordinates $\phi(w)$ contribute to the sum of squares.~\qed
\end{proof}
\step{int-form}{$(\iota\Zx)^TA(\iota\Zx)=\Zx^TB\Zx$.}
\begin{proof}
\pf\ \raggedright Only pairs $(\phi(u),\phi(v))$ contribute to the double sum; use \asm[2]{0}.~\qed
\end{proof}
\qedstep
\begin{proof}
\pf\ \raggedright Use \stepref{int-inj} and dimension invariance (Facts~\ref{facts:dim}); the two sum identities give the remaining assertions.~\qed
\end{proof}
\end{proof}
\step{int-upper}{For $0\le i<m$, $\beta_i\le\lambda_i$.}
\begin{proof}
\step{int-Bhead}{The extended span of the first $i+1$ eigendirections of $B$ has dimension $i+1$ and quotient at least $\beta_i$.}
\begin{proof}
\pf\ \raggedright Apply the head assertion of Lemma~\ref{lem:rayleigh-subspaces} to $B$, which is symmetric by \stepref{int-symmetric}, with its sorted basis from Theorem~\ref{thm:multiplicity} and $\beta_i=\nu_i(B)$; then transfer by \stepref{int-extension}.~\qed
\end{proof}
\step{int-compareupper}{Some nonzero vector there has quotient at most $\lambda_i$.}
\begin{proof}
\pf\ \raggedright Apply the first witness assertion of Lemma~\ref{lem:rayleigh-subspaces} to $A$ and the subspace in \stepref{int-Bhead}; its bound is $\nu_i(A)=\lambda_i$ by \asm[3]{0}.~\qed
\end{proof}
\qedstep
\begin{proof}
\pf\ \raggedright The witness in \stepref{int-compareupper} has quotient between $\beta_i$ and $\lambda_i$ by \stepref{int-Bhead}.~\qed
\end{proof}
\end{proof}
\step{int-lower}{For $0\le i<m$, $\lambda_{i+n-m}\le\beta_i$.}
\begin{proof}
\step{int-Btail}{The extended span of $B$-directions $i,\ldots,m-1$ has dimension $m-i$ and quotient at most $\beta_i$.}
\begin{proof}
\pf\ \raggedright Apply the tail assertion of Lemma~\ref{lem:rayleigh-subspaces} to $B$, symmetric by \stepref{int-symmetric}, with $\beta_i=\nu_i(B)$, and transfer by \stepref{int-extension}.~\qed
\end{proof}
\step{int-comparedown}{Some nonzero vector there has quotient at least $\lambda_{i+n-m}$.}
\begin{proof}
\pf\ \raggedright Apply the second witness assertion of Lemma~\ref{lem:rayleigh-subspaces} to $A$ at $i+n-m$, since $m-i=n-(i+n-m)$; its bound is $\lambda_{i+n-m}$ by \asm[3]{0}.~\qed
\end{proof}
\qedstep
\begin{proof}
\pf\ \raggedright The witness in \stepref{int-comparedown} has quotient at most $\beta_i$ by \stepref{int-Btail}.~\qed
\end{proof}
\end{proof}
\step{int-ends}{The largest cannot increase and the least cannot decrease.}
\begin{proof}
\pf\ \raggedright Take $i=0$ in \stepref{int-upper} and $i=m-1$ in \stepref{int-lower}; Lemma~\ref{lem:extremes} identifies the endpoints.~\qed
\end{proof}
\qedstep
\begin{proof}
\pf\ \raggedright Use \stepref{int-upper}, \stepref{int-lower}, \stepref{int-ends}.~\qed
\end{proof}
\end{proof}

\begin{definition}\label{def:support}
\lean{Matrix.positiveSupportGraph}
\uses{facts:graph}
\mbox{}\par\nobreak
\assume{\begin{pfenum}
\item $U$ is finite and $A:U\times U\to\re$ is symmetric and entrywise nonnegative.
\end{pfenum}}
For a symmetric entrywise nonnegative matrix $A$, its \emph{positive
support graph} has the same indices as vertices and an edge $uv$ when
$u\ne v$ and $A(u,v)>0$. Connectedness means that the vertex set is
nonempty and any two vertices are joined by a path. In particular a
single vertex with no edges is connected.
\end{definition}

\begin{theorem}[Symmetric Perron--Frobenius]\label{thm:perron}
\lean{Matrix.IsHermitian.exists_nonneg_eigenvector_largest, Matrix.IsHermitian.exists_pos_eigenvector_largest, Matrix.IsHermitian.multiplicity_largest_eq_one, Matrix.IsHermitian.abs_eigenvalue_le_largest}
\uses{def:support,facts:order,lem:extremes,thm:multiplicity,thm:rayleigh}
\mbox{}\par\nobreak
\assume{\begin{pfenum}
\item $U$ is finite and nonempty; $A:U\times U\to\re$ is symmetric and entrywise nonnegative.
\item $\rho=\lambda_{\max}(A)$. The strict-positivity and simplicity assertions also assume that the positive support graph of $A$ is connected.
\end{pfenum}}
\prove{\raggedright $\rho\ge0$, there is
a nonzero entrywise nonnegative vector $\Zx$ with $A\Zx=\rho\Zx$,
and $|\mu|\le\rho$ for every eigenvalue $\mu$.
If the positive support graph is connected, this vector is strictly
positive in every coordinate, $\mult_A(\rho)=1$, and every
nonzero nonnegative $\rho$-eigenvector is strictly positive.}
\end{theorem}
\begin{proof}
\pfsketch\ We will replace a top eigenvector by its absolute values, propagate zero coordinates along paths, and subtract a minimum ratio.
\step{pf-nonnegative}{$\rho\ge0$ and a nonnegative unit $\rho$-eigenvector exists.}
\begin{proof}
\step{pf-diagonal}{$0\le A(u,u)\le\rho$ for any $u\in U$.}
\begin{proof}
\pf\ \raggedright Test Theorem~\ref{thm:rayleigh} on the coordinate unit vector; entrywise nonnegativity gives the first inequality.~\qed
\end{proof}
\step{pf-choose}{\pick\ a unit $\rho$-eigenvector $\Zy$, and let $\Zx(u)=|\Zy(u)|$; then $\norm{\Zx}=1$.}
\begin{proof}
\pf\ \raggedright $\rho\in\Spec(A)$ by Lemma~\ref{lem:extremes}, and Theorem~\ref{thm:multiplicity} supplies a nonzero eigenvector, which we normalize; absolute value preserves each coordinate square.~\qed
\end{proof}
\step{pf-absform}{$\rho=\Zy^TA\Zy\le\Zx^TA\Zx$.}
\begin{proof}
\pf\ \raggedright For each $u,v$, multiply $\Zy(u)\Zy(v)\le|\Zy(u)||\Zy(v)|$ by $A(u,v)\ge0$ and sum.~\qed
\end{proof}
\step{pf-equality}{$A\Zx=\rho\Zx$.}
\begin{proof}
\pf\ \raggedright By Theorem~\ref{thm:rayleigh}, $\Zx^TA\Zx\le\rho$; with \stepref{pf-absform} this is equality, and the equality criterion of Theorem~\ref{thm:rayleigh} applies.~\qed
\end{proof}
\qedstep
\begin{proof}
\pf\ \raggedright \stepref{pf-diagonal} gives $\rho\ge0$; \stepref{pf-choose} and \stepref{pf-equality} give the nonnegative unit $\rho$-eigenvector $\Zx$.~\qed
\end{proof}
\end{proof}
\step{pf-absolute}{Every eigenvalue $\mu$ satisfies $|\mu|\le\rho$.}
\begin{proof}
\step{pf-mu}{For a nonzero $\mu$-eigenvector $\Zw$, $|\mu|\norm{\Zw}^2=|\Zw^TA\Zw|$.}
\begin{proof}
\pf\ \raggedright Take absolute values of its quadratic eigenvector equation.~\qed
\end{proof}
\step{pf-triangle}{$|\Zw^TA\Zw|\le |\Zw|^TA|\Zw|\le\rho\norm{\Zw}^2$.}
\begin{proof}
\pf\ \raggedright The finite-sum triangle inequality gives the first bound; Theorem~\ref{thm:rayleigh} gives the second, since absolute values preserve the norm.~\qed
\end{proof}
\step{pf-divide}{$|\mu|\le\rho$.}
\begin{proof}
\pf\ \raggedright Combine \stepref{pf-mu}, \stepref{pf-triangle} and divide by $\norm{\Zw}^2>0$.~\qed
\end{proof}
\qedstep
\begin{proof}
\pf\ \raggedright Use \stepref{pf-divide}.~\qed
\end{proof}
\end{proof}
\step{pf-propagate}{\assume{The positive support graph of $A$ is connected; $\Zz$ is entrywise nonnegative, $A\Zz=b\Zz$ for some $b\in\re$, and $\Zz(u_0)=0$.}\prove{$\Zz=0$.}}
\begin{proof}
\step{pf-zero-sum}{If $\Zz(u)=0$, then $\sum_vA(u,v)\Zz(v)=0$.}
\begin{proof}
\pf\ \raggedright Evaluate the eigenvector equation of \asm{1} at $u$.~\qed
\end{proof}
\step{pf-edgezero}{If $\Zz(u)=0$, then $\Zz(v)=0$ for every support neighbour $v$ of $u$.}
\begin{proof}
\pf\ \raggedright Each term in \stepref{pf-zero-sum} is nonnegative; its $v$th term is zero by Facts~\ref{facts:order}, and $A(u,v)>0$ by Definition~\ref{def:support}.~\qed
\end{proof}
\step{pf-path}{$\Zz$ vanishes at every vertex of a support path starting at $u_0$.}
\begin{proof}
\step{pf-path0}{It holds for a path of length zero.}
\begin{proof}
\pf\ \raggedright Its only vertex is $u_0$, and $\Zz(u_0)=0$ by \asm{1}.~\qed
\end{proof}
\step{pf-pathS}{If it holds for a path, it holds for its extension by one edge.}
\begin{proof}
\pf\ \raggedright Apply \stepref{pf-edgezero} to the last vertex of the path and the new edge.~\qed
\end{proof}
\qedstep
\begin{proof}
\pf\ \raggedright Induction on path length using \stepref{pf-path0} and \stepref{pf-pathS}.~\qed
\end{proof}
\end{proof}
\step{pf-allzero}{Every coordinate of $\Zz$ is zero.}
\begin{proof}
\pf\ \raggedright Connectedness in \asm{1} supplies a support path from $u_0$ to each vertex; apply \stepref{pf-path}.~\qed
\end{proof}
\qedstep
\begin{proof}
\pf\ \raggedright Use \stepref{pf-allzero}.~\qed
\end{proof}
\end{proof}
\step{pf-simple}{\assume{The positive support graph of $A$ is connected.}\prove{Every nonzero nonnegative $\rho$-eigenvector is positive, and $\mult_A(\rho)=1$.}}
\begin{proof}
\step{pf-positive}{Every coordinate of a nonzero nonnegative eigenvector is positive.}
\begin{proof}
\pf\ \raggedright A zero coordinate would make the vector zero, by \stepref{pf-propagate} and \asm{1}.~\qed
\end{proof}
\step{pf-minratio}{For a positive $\rho$-eigenvector $\Zx$ and any $\rho$-eigenvector $\Zy$, \pick\ $u_0$ with $\Zy(u_0)/\Zx(u_0)=b\triangleq\min_u\Zy(u)/\Zx(u)$.}
\begin{proof}
\pf\ \raggedright A positive $\rho$-eigenvector exists by \stepref{pf-nonnegative} and \stepref{pf-positive}; a finite nonempty set of real ratios has an attained minimum.~\qed
\end{proof}
\step{pf-difference}{$\Zy-b\Zx$ is a nonnegative top eigenvector with a zero coordinate.}
\begin{proof}
\step{pf-difference-eigen}{$A(\Zy-b\Zx)=\rho(\Zy-b\Zx)$.}
\begin{proof}
\pf\ \raggedright Linearity and the two top-eigenvector equations.~\qed
\end{proof}
\step{pf-difference-nonneg}{Every coordinate of $\Zy-b\Zx$ is nonnegative.}
\begin{proof}
\pf\ \raggedright By \stepref{pf-minratio}, $b\le\Zy(u)/\Zx(u)$; multiply by $\Zx(u)>0$.~\qed
\end{proof}
\step{pf-difference-zero}{One coordinate is zero.}
\begin{proof}
\pf\ \raggedright At an attaining index in \stepref{pf-minratio}, the ratio equals $b$.~\qed
\end{proof}
\qedstep
\begin{proof}
\pf\ \raggedright Use \stepref{pf-difference-eigen}, \stepref{pf-difference-nonneg}, \stepref{pf-difference-zero}.~\qed
\end{proof}
\end{proof}
\step{pf-line}{$\Zy=b\Zx$.}
\begin{proof}
\pf\ \raggedright Apply \stepref{pf-propagate}, with \asm{1}, to the vector of \stepref{pf-difference}.~\qed
\end{proof}
\step{pf-dimension}{The multiplicity of $\rho$ is one.}
\begin{proof}
\pf\ \raggedright The nonzero vector $\Zx$ spans the eigenspace by \stepref{pf-line}; use Theorem~\ref{thm:multiplicity}.~\qed
\end{proof}
\qedstep
\begin{proof}
\pf\ \raggedright Use \stepref{pf-positive}, \stepref{pf-dimension}.~\qed
\end{proof}
\end{proof}
\qedstep
\begin{proof}
\pf\ \raggedright Use \stepref{pf-nonnegative}, \stepref{pf-absolute}, \stepref{pf-simple}.~\qed
\end{proof}
\end{proof}

\begin{lemma}[A positive supervector bounds each eigenvalue]\label{lem:supervector}
\lean{Matrix.eigenvalue_le_of_pos_supervector}
\uses{}
\mbox{}\par\nobreak
\assume{\begin{pfenum}
\item $A$ is real, entrywise nonnegative on a finite nonempty set, $\Zx>0$, and $A\Zx\le b\Zx$.
\item $A\Zy=\theta\Zy$ and $\Zy\ne0$.
\end{pfenum}}
\prove{\raggedright $\theta\le b$. Symmetry of $A$ is not required.}
\end{lemma}
\begin{proof}
\pfsketch\ We will follow the maximum-ratio proof in \code{LineGraphs/Spectral/PositiveVector.lean}.
\step{ratio-max}{\pick\ $u$ maximizing $|\Zy(u)|/\Zx(u)$, and denote the maximum by $c$.}
\begin{proof}
\pf\ \raggedright A finite nonempty set of real numbers has an attained maximum.~\qed
\end{proof}
\step{ratio-positive}{$c>0$ and $|\Zy(u)|=c\Zx(u)>0$.}
\begin{proof}
\pf\ \raggedright A nonzero coordinate of $\Zy$ gives a positive ratio; use maximality and the positivity of $\Zx$.~\qed
\end{proof}
\step{ratio-bound}{$|\Zy(v)|\le c\Zx(v)$ for every $v$.}
\begin{proof}
\pf\ \raggedright Multiply the maximal-ratio inequality by $\Zx(v)>0$.~\qed
\end{proof}
\step{ratio-eigen}{$\theta|\Zy(u)|\le |(A\Zy)(u)|$.}
\begin{proof}
\pf\ \raggedright Use $\theta\le|\theta|$ and the eigenvector equation.~\qed
\end{proof}
\step{ratio-sum}{$|(A\Zy)(u)|\le\sum_v A(u,v)|\Zy(v)|\le c(A\Zx)(u)$.}
\begin{proof}
\pf\ \raggedright The finite-sum triangle inequality and nonnegative entries give the first inequality; sum \stepref{ratio-bound} with these coefficients for the second.~\qed
\end{proof}
\step{ratio-conclude}{$\theta\le b$.}
\begin{proof}
\pf\ \raggedright Combine \stepref{ratio-eigen}, \stepref{ratio-sum} and $c(A\Zx)(u)\le cb\Zx(u)=b|\Zy(u)|$; divide by the positive number in \stepref{ratio-positive}.~\qed
\end{proof}
\qedstep
\begin{proof}
\pf\ \raggedright Use \stepref{ratio-conclude}.~\qed
\end{proof}
\end{proof}

\begin{corollary}[Positive-vector certificates]\label{cor:positive-certificate}
\lean{Matrix.IsHermitian.largestEigenvalue_eq_of_pos_eigenvector, Matrix.IsHermitian.largestEigenvalue_le_of_pos_supervector}
\uses{def:extremes,lem:extremes,lem:supervector,thm:multiplicity}
\mbox{}\par\nobreak
\assume{\begin{pfenum}
\item $A$ is real symmetric and entrywise nonnegative on a finite nonempty set; $\Zx>0$ and $b\in\re$.
\end{pfenum}}
\prove{\raggedright If
$A\Zx=b\Zx$, then $b=\lambda_{\max}(A)$, whether or not its support
is connected. More generally, $A\Zx\le b\Zx$ coordinatewise implies
$\lambda_{\max}(A)\le b$.}
\end{corollary}
\begin{proof}
\pfsketch\ We will apply Lemma~\ref{lem:supervector} to a top eigenvector for the upper bound, and use the eigenvector $\Zx$ itself for the reverse bound.
\step{pos-upper}{If $A\Zx\le b\Zx$, then $\lambda_{\max}(A)\le b$.}
\begin{proof}
\pf\ \raggedright Apply Lemma~\ref{lem:supervector} to the sorted basis vector $\Zw_0$ of Theorem~\ref{thm:multiplicity}, whose eigenvalue is $\nu_0=\lambda_{\max}(A)$ by Definition~\ref{def:extremes}.~\qed
\end{proof}
\step{pos-occurs}{If $A\Zx=b\Zx$, then $b\in\Spec(A)$.}
\begin{proof}
\pf\ \raggedright $\Zx>0$ on a nonempty set implies $\Zx\ne0$; apply Theorem~\ref{thm:multiplicity}.~\qed
\end{proof}
\step{pos-equal}{In that case $b=\lambda_{\max}(A)$.}
\begin{proof}
\pf\ \raggedright Lemma~\ref{lem:extremes} gives $b\le\lambda_{\max}$ from \stepref{pos-occurs}; \stepref{pos-upper} gives the reverse inequality.~\qed
\end{proof}
\qedstep
\begin{proof}
\pf\ \raggedright Use \stepref{pos-upper}, \stepref{pos-equal}.~\qed
\end{proof}
\end{proof}

\begin{theorem}[Strict monotonicity]\label{thm:strict}
\lean{Matrix.IsHermitian.largestEigenvalue_lt_of_lt_submatrix, SimpleGraph.largestEigenvalue_lt_of_proper_subgraph}
\uses{def:extremes,facts:graph,facts:order,thm:perron,thm:rayleigh}
\mbox{}\par\nobreak
\assume{\begin{pfenum}
\item $U,W$ are finite and nonempty. $A:U\times U\to\re$ is symmetric, nonnegative and has connected positive support.
\item $B:W\times W\to\re$ is symmetric and nonnegative; $\phi:W\hookrightarrow U$ is injective and $B(u,v)\le A(\phi(u),\phi(v))$.
\item For the graph assertion, $A=A(X)$ for a connected finite simple graph $X$ on $U$, and $B=A(X')$ for a nonempty subgraph $X'$ on $W$, given by the injective vertex map $\phi$ preserving edges; $X'$ is proper if it omits a vertex or an edge.
\end{pfenum}}
\prove{\raggedright $\lambda_{\max}(B)\le\lambda_{\max}(A)$.
The inequality is strict if $\phi$ is not surjective, or if one of
these entry inequalities is strict. In particular every nonempty
proper subgraph of a connected finite graph has strictly smaller
largest eigenvalue. A subgraph here may omit edges as well as vertices.}
\end{theorem}
\begin{proof}
\pfsketch\ We will extend a top eigenvector by zero. Equality in both quadratic-form inequalities forces full vertex support and equality of every entry.
\step{strict-choose}{\pick\ a nonnegative unit top eigenvector $\Zy$ of $B$; let $\Zz$ be its extension by zero along $\phi$.}
\begin{proof}
\pf\ \raggedright Theorem~\ref{thm:perron} gives a nonzero nonnegative top eigenvector; divide by its positive norm.~\qed
\end{proof}
\step{strict-weak}{$\Zy^TB\Zy=\lambda_{\max}(B)\le\Zz^TA\Zz\le\lambda_{\max}(A)$, and $\norm{\Zz}=1$.}
\begin{proof}
\step{strict-vector}{The chosen $\Zz$ is nonnegative and has norm one.}
\begin{proof}
\pf\ \raggedright In \stepref{strict-choose}, extension by zero along an injection preserves each coordinate square and its sum.~\qed
\end{proof}
\step{strict-forms}{$\lambda_{\max}(B)=\Zy^TB\Zy\le\Zz^TA\Zz\le\lambda_{\max}(A)$.}
\begin{proof}
\pf\ \raggedright The entry inequalities may be multiplied by $\Zy(u)\Zy(v)\ge0$ and summed; the final bound is Theorem~\ref{thm:rayleigh}, with extremes defined in Definition~\ref{def:extremes}.~\qed
\end{proof}
\qedstep
\begin{proof}
\pf\ \raggedright Use \stepref{strict-vector}, \stepref{strict-forms}.~\qed
\end{proof}
\end{proof}
\step{strict-equal}{If the two largest eigenvalues are equal, $\phi$ is surjective and every entry inequality is equality.}
\begin{proof}
\step{strict-eigen}{$\Zz$ is a top eigenvector of $A$.}
\begin{proof}
\pf\ \raggedright Equality in \stepref{strict-weak} gives the equality case of Theorem~\ref{thm:rayleigh}.~\qed
\end{proof}
\step{strict-positive}{$\Zz(u)>0$ for all $u$, so $\phi$ is surjective.}
\begin{proof}
\pf\ \raggedright Connectedness and Theorem~\ref{thm:perron} make \stepref{strict-eigen} positive; extension by zero cannot omit a coordinate.~\qed
\end{proof}
\step{strict-difference}{$0=\sum_{u,v}(A(\phi(u),\phi(v))-B(u,v))\Zy(u)\Zy(v)$.}
\begin{proof}
\pf\ \raggedright Subtract the equal quadratic forms in \stepref{strict-weak}.~\qed
\end{proof}
\step{strict-entries}{$A(\phi(u),\phi(v))=B(u,v)$ for every pair.}
\begin{proof}
\pf\ \raggedright The summands in \stepref{strict-difference} are nonnegative; each vanishes by Facts~\ref{facts:order}, and \stepref{strict-positive} makes its coordinate product positive.~\qed
\end{proof}
\qedstep
\begin{proof}
\pf\ \raggedright Use \stepref{strict-positive}, \stepref{strict-entries}.~\qed
\end{proof}
\end{proof}
\step{strict-final}{Either strictness hypothesis implies strict inequality; the graph assertion follows.}
\begin{proof}
\step{strict-contrapositive}{An omitted vertex or a strict entry excludes equality.}
\begin{proof}
\pf\ \raggedright Take the contrapositive of \stepref{strict-equal}.~\qed
\end{proof}
\step{strict-graph-hyp}{In the graph assertion, $A=A(X)$ and $B=A(X')$ satisfy \asm[1]{0} and \asm[2]{0}.}
\begin{proof}
\pf\ \raggedright Adjacency matrices are symmetric with entries $0,1$ (Facts~\ref{facts:graph}); the positive support graph of $A(X)$ is $X$, which is connected; an edge-preserving injection gives $B(u,v)\le A(\phi(u),\phi(v))$.~\qed
\end{proof}
\step{strict-graph}{A proper subgraph gives one of the two strictness conditions.}
\begin{proof}
\step{strict-missingvertex}{\case{The subgraph omits a vertex.}}
\begin{proof}
\pf\ \raggedright The omitted vertex is not in the image of the vertex injection, which is thus not surjective.~\qed
\end{proof}
\step{strict-missingedge}{\case{The subgraph contains every vertex and omits an edge.}}
\begin{proof}
\pf\ \raggedright At the endpoints of that edge the smaller adjacency entry is $0$ and the larger one is $1$, so an entry inequality is strict.~\qed
\end{proof}
\qedstep
\begin{proof}
\pf\ \raggedright By \asm[3]{0} a proper subgraph omits a vertex or an edge, so the two cases are exhaustive.~\qed
\end{proof}
\end{proof}
\qedstep
\begin{proof}
\pf\ \raggedright Use \stepref{strict-contrapositive}, \stepref{strict-graph-hyp}, \stepref{strict-graph}.~\qed
\end{proof}
\end{proof}
\qedstep
\begin{proof}
\pf\ \raggedright Use \stepref{strict-weak}, \stepref{strict-final}.~\qed
\end{proof}
\end{proof}

\section{Walks, Laplacians, and regular graphs}

\begin{theorem}[Walks and moments]\label{thm:walks}
\lean{SimpleGraph.trace_pow_adjMatrix_eq_closedWalks, SimpleGraph.sum_sq_spectrumMultiset_eq_twice_card_edges, SimpleGraph.trace_cube_adjMatrix_eq_six_mul_triangles, SimpleGraph.adjMatrix_cube_apply_self_eq_twice_triangles}
\uses{def:graph,facts:graph,thm:multiplicity,thm:trace}
\mbox{}\par\nobreak
\assume{\begin{pfenum}
\item $X$ is a finite simple graph on $U$, $|U|=n$, $|E(X)|=m$, $A=A(X)$, and $j\in\nn$.
\item $t$ denotes the number of three-element vertex sets inducing a triangle.
\end{pfenum}}
\prove{\raggedright For $j\ge0$, $(A^j)(u,v)$ counts walks of
length $j$ from $u$ to $v$, and $\tr A^j$ counts closed walks with
a distinguished starting vertex. In particular
\[
 \sum_{\mu\in\Spec(X)}1=n,\quad
 \sum_{\mu\in\Spec(X)}\mu=0,\quad
 \sum_{\mu\in\Spec(X)}\mu^2=2m=\sum_u\deg(u),\quad
 \sum_{\mu\in\Spec(X)}\mu^3=6t,
\]
where $t$ is the number of triangles and all sums count multiplicity.
At each vertex $u$,
$(A^3)(u,u)=2\,\#\{\text{triangles containing }u\}$.
The total number of walks of length $j$ is $\Zj^TA^j\Zj$; it need
not equal $\tr A^j$.}
\end{theorem}
\begin{proof}
\pfsketch\ We will count walks recursively, then count closed walks of lengths zero through three.
\step{walk-entry}{$(A^j)(u,v)$ counts walks of length $j$ from $u$ to $v$.}
\begin{proof}
\pf\ \raggedright Apply \code{SimpleGraph.adjMatrix_pow_apply_eq_card_walk} in Facts~\ref{facts:graph} over $\re$.~\qed
\end{proof}
\step{walk-traces}{Trace counts closed walks with a starting vertex; $\Zj^TA^j\Zj$ counts all walks.}
\begin{proof}
\step{walk-closed}{$\tr A^j=\sum_u(A^j)(u,u)$.}
\begin{proof}
\pf\ \raggedright This is the definition of trace.~\qed
\end{proof}
\step{walk-total}{$\Zj^TA^j\Zj=\sum_{u,v}(A^j)(u,v)$.}
\begin{proof}
\pf\ \raggedright Expand multiplication by the all-ones vectors.~\qed
\end{proof}
\qedstep
\begin{proof}
\pf\ \raggedright Use \stepref{walk-closed}, \stepref{walk-total}.~\qed
\end{proof}
\end{proof}
\step{walk-small}{The first four closed-walk counts are $n,0,2m,6t$.}
\begin{proof}
\step{walk-length0}{There are $n$ length-zero closed walks.}
\begin{proof}
\pf\ \raggedright Choose their single vertex.~\qed
\end{proof}
\step{walk-length1}{There are no length-one closed walks.}
\begin{proof}
\pf\ \raggedright Apply \code{SimpleGraph.trace_adjMatrix} in Facts~\ref{facts:graph}.~\qed
\end{proof}
\step{walk-length2}{There are $2m=\sum_u\deg(u)$ length-two closed walks.}
\begin{proof}
\pf\ \raggedright Apply \code{SimpleGraph.adjMatrix_mul_self_apply_self} in Facts~\ref{facts:graph}, sum over vertices, and count incident vertex-edge pairs by their two endpoints.~\qed
\end{proof}
\step{walk-length3}{There are $6t$ length-three closed walks.}
\begin{proof}
\step{walk-distinct}{In $u,v,w,u$, the first three vertices are distinct.}
\begin{proof}
\pf\ \raggedright Any equality would make one of its three required edges a loop.~\qed
\end{proof}
\step{walk-triangle}{Each such walk gives a triangle and a cyclic ordering with a distinguished start.}
\begin{proof}
\pf\ \raggedright The edges $uv,vw,wu$ in the walk are precisely the three triangle edges.~\qed
\end{proof}
\step{walk-six}{Each triangle gives exactly $3\cdot2=6$ such walks.}
\begin{proof}
\pf\ \raggedright Choose one of its three starting vertices and then one of the two neighbours as the next vertex.~\qed
\end{proof}
\qedstep
\begin{proof}
\pf\ \raggedright Use \stepref{walk-distinct}, \stepref{walk-triangle}, \stepref{walk-six}.~\qed
\end{proof}
\end{proof}
\qedstep
\begin{proof}
\pf\ \raggedright Use \stepref{walk-length0}, \stepref{walk-length1}, \stepref{walk-length2}, \stepref{walk-length3}.~\qed
\end{proof}
\end{proof}
\step{walk-spectral}{The four moment identities follow, with multiplicities.}
\begin{proof}
\pf\ \raggedright Apply Theorem~\ref{thm:trace} to \stepref{walk-small}; the zeroth moment counts $n$ entries by Theorem~\ref{thm:multiplicity}.~\qed
\end{proof}
\step{walk-local}{$(A^3)(u,u)$ is twice the number of triangles containing $u$.}
\begin{proof}
\pf\ \raggedright Apply \code{SimpleGraph.adjMatrix_pow_apply_eq_card_walk} in Facts~\ref{facts:graph} at power $3$: the entry counts closed walks $u,v,w,u$. Looplessness makes $u,v,w$ distinct; each triangle through $u$ supplies precisely the two possible orders of $v,w$.~\qed
\end{proof}
\qedstep
\begin{proof}
\pf\ \raggedright Use \stepref{walk-entry}, \stepref{walk-traces}, \stepref{walk-local}, \stepref{walk-spectral}.~\qed
\end{proof}
\end{proof}

\begin{theorem}[Laplacian and signless Laplacian kernels]\label{thm:laplacian}
\lean{SimpleGraph.signlessLapMatrix, SimpleGraph.signlessLapMatrix_mulVec_eq_zero_iff, SimpleGraph.finrank_ker_signlessLapMatrix_eq_card_bipartite_components}
\uses{facts:dim,facts:laplacian,facts:order,facts:positive,thm:psd}
\mbox{}\par\nobreak
\assume{\begin{pfenum}
\item $X$ is a finite simple graph on $U$, $A=A(X)$ and $D$ is its diagonal degree matrix.
\item $L=D-A$, $K=D+A$, and $\Zx\in\re^U$. A component is called bipartite if its vertices admit a partition crossed by every edge.
\end{pfenum}}
\prove{\raggedright \[
 \Zx^TL\Zx=\sum_{\{u,v\}\in E(X)}(\Zx(u)-\Zx(v))^2,\qquad
 \Zx^TK\Zx=\sum_{\{u,v\}\in E(X)}(\Zx(u)+\Zx(v))^2.
\]
Both matrices are positive semidefinite. Their kernels consist,
respectively, of functions constant across each edge and functions
whose values sum to zero across each edge. Thus
\[
 \dim\ker L=\#\{\text{components of }X\},\qquad
 \dim\ker K=\#\{\text{bipartite components of }X\}.
\]
An isolated vertex counts as a bipartite component.}
\end{theorem}
\begin{proof}
\pfsketch\ We will expand edge sums of squares and construct bases of the two kernels component by component.
\step{lap-forms}{The two displayed quadratic-form identities hold and both matrices are positive semidefinite.}
\begin{proof}
\step{lap-diagonal}{$\sum_{\{u,v\}\in E}(\Zx(u)^2+\Zx(v)^2)=\sum_u\deg(u)\Zx(u)^2$.}
\begin{proof}
\pf\ \raggedright Count one occurrence of $\Zx(u)^2$ for each edge incident with $u$.~\qed
\end{proof}
\step{lap-cross}{$\Zx^TA\Zx=2\sum_{\{u,v\}\in E}\Zx(u)\Zx(v)$.}
\begin{proof}
\pf\ \raggedright The ordered adjacency sum has the two orientations of each edge.~\qed
\end{proof}
\step{lap-minus}{$\Zx^TL\Zx=\sum_{\{u,v\}\in E}(\Zx(u)-\Zx(v))^2$.}
\begin{proof}
\pf\ \raggedright Subtract \stepref{lap-cross} from \stepref{lap-diagonal} and expand the square.~\qed
\end{proof}
\step{lap-plus}{$\Zx^TK\Zx=\sum_{\{u,v\}\in E}(\Zx(u)+\Zx(v))^2$.}
\begin{proof}
\pf\ \raggedright Add \stepref{lap-cross} to \stepref{lap-diagonal} and expand the square.~\qed
\end{proof}
\step{lap-positive}{$L,K\succeq0$.}
\begin{proof}
\step{lap-Lpositive}{$L\succeq0$.}
\begin{proof}
\pf\ \raggedright Apply \code{SimpleGraph.posSemidef_lapMatrix} in Facts~\ref{facts:laplacian} over $\re$.~\qed
\end{proof}
\step{lap-Kpositive}{$K\succeq0$.}
\begin{proof}
\pf\ \raggedright $K$ is symmetric and its form in \stepref{lap-plus} is nonnegative; apply \code{Matrix.posSemidef_iff_dotProduct_mulVec} in Facts~\ref{facts:positive}.~\qed
\end{proof}
\qedstep
\begin{proof}
\pf\ \raggedright Use \stepref{lap-Lpositive}, \stepref{lap-Kpositive}.~\qed
\end{proof}
\end{proof}
\qedstep
\begin{proof}
\pf\ \raggedright Use \stepref{lap-minus}, \stepref{lap-plus}, \stepref{lap-positive}.~\qed
\end{proof}
\end{proof}
\step{lap-edges}{The kernels are exactly the functions satisfying the stated edge relations.}
\begin{proof}
\step{lap-kernelL}{$L\Zx=0$ iff $\Zx(u)=\Zx(v)$ on every edge.}
\begin{proof}
\step{lap-Lreachable}{$L\Zx=0$ iff $\Zx$ is constant on reachable pairs.}
\begin{proof}
\pf\ \raggedright Apply \code{SimpleGraph.lapMatrix_mulVec_eq_zero_iff_forall_reachable} in Facts~\ref{facts:laplacian}.~\qed
\end{proof}
\step{lap-Ltoedges}{Constancy on reachable pairs implies constancy across edges.}
\begin{proof}
\pf\ \raggedright An edge is a path of length one.~\qed
\end{proof}
\step{lap-Lfromedges}{Constancy across edges implies constancy along paths.}
\begin{proof}
\pf\ \raggedright Induct on path length: the empty path has equal endpoints, and each successive edge preserves the value.~\qed
\end{proof}
\qedstep
\begin{proof}
\pf\ \raggedright Use \stepref{lap-Lreachable}, \stepref{lap-Ltoedges}, \stepref{lap-Lfromedges}.~\qed
\end{proof}
\end{proof}
\step{lap-kernelK}{$K\Zx=0$ iff $\Zx(u)=-\Zx(v)$ on every edge.}
\begin{proof}
\pf\ \raggedright Apply the zero-form criterion of Theorem~\ref{thm:psd} and the zero-sum-of-squares criterion in Facts~\ref{facts:order} to the plus-square identity in \stepref{lap-forms}.~\qed
\end{proof}
\qedstep
\begin{proof}
\pf\ \raggedright Use \stepref{lap-kernelL}, \stepref{lap-kernelK}.~\qed
\end{proof}
\end{proof}
\step{lap-Ldimension}{$\dim\ker L$ equals the number of components.}
\begin{proof}
\pf\ \raggedright Apply \code{SimpleGraph.card_connectedComponent_eq_finrank_ker_toLin'_lapMatrix} in Facts~\ref{facts:laplacian}; its kernel already uses the function-coordinate linear map.~\qed
\end{proof}
\step{lap-Kcomponent}{On a fixed connected component the signless kernel has dimension one if bipartite and zero otherwise.}
\begin{proof}
\step{lap-Kpaths}{Along a path of length $j$ from $u$ to $v$, $\Zx(v)=(-1)^j\Zx(u)$.}
\begin{proof}
\pf\ \raggedright Induction on $j$ using the sign relation in \stepref{lap-edges}.~\qed
\end{proof}
\step{lap-Kinject}{Evaluation at a fixed vertex is injective on the component kernel.}
\begin{proof}
\pf\ \raggedright By \stepref{lap-Kpaths}, a zero value there forces all values to be zero.~\qed
\end{proof}
\step{lap-Kpartition}{A nonzero kernel function gives a bipartition into its positive and negative vertices.}
\begin{proof}
\pf\ \raggedright By \stepref{lap-Kpaths} all absolute values equal a positive number; the edge relation in \stepref{lap-edges} reverses the sign.~\qed
\end{proof}
\step{lap-Ksign}{A bipartition gives a nonzero kernel function equal to $1$ and $-1$ on the two parts.}
\begin{proof}
\pf\ \raggedright Each edge has opposite signs, so \stepref{lap-edges} applies; a component has a vertex.~\qed
\end{proof}
\step{lap-Kdim}{The component dimension is one in the bipartite case and zero otherwise.}
\begin{proof}
\step{lap-Katmost}{The dimension is at most one.}
\begin{proof}
\pf\ \raggedright Evaluation injects it into $\re$ by \stepref{lap-Kinject}; apply dimension monotonicity in Facts~\ref{facts:dim}.~\qed
\end{proof}
\step{lap-Kbipartite}{\case{The component is bipartite.}}
\begin{proof}
\pf\ \raggedright The nonzero vector in \stepref{lap-Ksign} gives dimension at least one; with \stepref{lap-Katmost} the dimension is one.~\qed
\end{proof}
\step{lap-Knonbipartite}{\case{The component is not bipartite.}}
\begin{proof}
\pf\ \raggedright By \stepref{lap-Kpartition} any nonzero kernel function would give a bipartition; thus the kernel is the zero space.~\qed
\end{proof}
\qedstep
\begin{proof}
\pf\ \raggedright The two cases are exhaustive.~\qed
\end{proof}
\end{proof}
\qedstep
\begin{proof}
\pf\ \raggedright Use \stepref{lap-Kdim}.~\qed
\end{proof}
\end{proof}
\step{lap-counts}{The two kernel dimensions are the numbers of components and bipartite components.}
\begin{proof}
\step{lap-Kspan}{Extend one signed function from each bipartite component by zero; these functions span $\ker K$.}
\begin{proof}
\pf\ \raggedright Use the one-dimensional or zero component spaces in \stepref{lap-Kcomponent} and restrict any kernel function to each component.~\qed
\end{proof}
\step{lap-Kind}{The extended functions are independent.}
\begin{proof}
\pf\ \raggedright Evaluate a zero combination at the chosen vertex of each component; supports are disjoint.~\qed
\end{proof}
\step{lap-isolated}{An isolated vertex contributes one dimension to each kernel.}
\begin{proof}
\pf\ \raggedright There is no edge relation at that vertex, and its singleton component has a bipartition.~\qed
\end{proof}
\step{lap-countbasis}{The dimension formulas hold, including an empty vertex set.}
\begin{proof}
\pf\ \raggedright Use \stepref{lap-Ldimension} for $L$. For $K$, count the basis given by \stepref{lap-Kspan} and \stepref{lap-Kind}; an empty basis has dimension zero.~\qed
\end{proof}
\qedstep
\begin{proof}
\pf\ \raggedright Use \stepref{lap-countbasis}, \stepref{lap-isolated}.~\qed
\end{proof}
\end{proof}
\qedstep
\begin{proof}
\pf\ \raggedright Use \stepref{lap-forms}, \stepref{lap-edges}, \stepref{lap-counts}.~\qed
\end{proof}
\end{proof}

\begin{theorem}[SP1: regular graphs]\label{thm:sp1}
\lean{SimpleGraph.IsRegularOfDegree.largestEigenvalue_eq, SimpleGraph.IsRegularOfDegree.multiplicity_eq_card_components}
\uses{def:graph,facts:graph,lem:extremes,thm:laplacian,thm:multiplicity,thm:psd}
\mbox{}\par\nobreak
\assume{\begin{pfenum}
\item $X$ is a finite simple $k$-regular graph on a nonempty set $U$, where $k\in\nn$.
\end{pfenum}}
\prove{\raggedright Its largest eigenvalue
is $k$, every eigenvalue lies in $[-k,k]$, and
\[
 \mult_X(k)=\#\{\text{connected components of }X\}.
\]}
\end{theorem}
\begin{proof}
\pfsketch\ We will use the constant vector for occurrence of $k$ and the two Laplacians for bounds and multiplicity.
\step{sp1-eigen}{$k$ occurs and every eigenvalue lies in $[-k,k]$.}
\begin{proof}
\step{sp1-ones}{$A\Zj=k\Zj$ and $\Zj\ne0$.}
\begin{proof}
\pf\ \raggedright Apply \code{SimpleGraph.adjMatrix_mulVec_const_apply_of_regular} in Facts~\ref{facts:graph} with constant $1$; nonemptiness makes $\Zj\ne0$.~\qed
\end{proof}
\step{sp1-bounds}{$kI-A\succeq0$ and $kI+A\succeq0$.}
\begin{proof}
\pf\ \raggedright These are $L,K$ in Theorem~\ref{thm:laplacian} for constant degree.~\qed
\end{proof}
\step{sp1-extremes}{$-k\le\mu\le k$ for every eigenvalue, and the largest is $k$.}
\begin{proof}
\pf\ \raggedright Apply the shifts in Theorem~\ref{thm:psd} to \stepref{sp1-bounds}. By \stepref{sp1-ones} and Theorem~\ref{thm:multiplicity}, $k\in\Spec(X)$; Lemma~\ref{lem:extremes} then gives $\lambda_{\max}(X)=k$.~\qed
\end{proof}
\qedstep
\begin{proof}
\pf\ \raggedright Use \stepref{sp1-ones}, \stepref{sp1-extremes}.~\qed
\end{proof}
\end{proof}
\step{sp1-components}{$\mult_X(k)$ is the number of components.}
\begin{proof}
\step{sp1-ker}{$V_A(k)=\ker(kI-A)$.}
\begin{proof}
\pf\ \raggedright Rearrange the defining eigenvector equation.~\qed
\end{proof}
\step{sp1-count}{The dimension of this space is the number of components.}
\begin{proof}
\pf\ \raggedright Apply Theorem~\ref{thm:laplacian} to \stepref{sp1-ker}.~\qed
\end{proof}
\step{sp1-mult}{Multiplicity is that dimension.}
\begin{proof}
\pf\ \raggedright Apply Theorem~\ref{thm:multiplicity}.~\qed
\end{proof}
\qedstep
\begin{proof}
\pf\ \raggedright Use \stepref{sp1-count}, \stepref{sp1-mult}.~\qed
\end{proof}
\end{proof}
\qedstep
\begin{proof}
\pf\ \raggedright Use \stepref{sp1-eigen}, \stepref{sp1-components}.~\qed
\end{proof}
\end{proof}

\begin{theorem}[SP2: regularity detected by the spectrum]\label{thm:sp2}
\lean{SimpleGraph.isRegular_iff_sum_sq_eigenvalues}
\uses{thm:rayleigh,thm:sp1,thm:walks}
\mbox{}\par\nobreak
\assume{\begin{pfenum}
\item $X$ is a finite simple graph on $n>0$ vertices.
\end{pfenum}}
\prove{\raggedright \[
 X\text{ is regular}\ \Longleftrightarrow\
 \sum_{\mu\in\Spec(X)}\mu^2=n\lambda_{\max}(X).
\]}
\end{theorem}
\begin{proof}
\pfsketch\ We will recognize equality in the Rayleigh bound at the all-ones vector.
\step{sp2-Rayleigh}{The displayed spectral identity is equivalent to regularity.}
\begin{proof}
\step{sp2-forms}{$\norm{\Zj}^2=n$ and $\Zj^TA\Zj=\sum_u\deg(u)=\sum_{\mu\in\Spec(X)}\mu^2$.}
\begin{proof}
\pf\ \raggedright The all-ones sums and Theorem~\ref{thm:walks} give the equalities.~\qed
\end{proof}
\step{sp2-equality}{The identity is equivalent to $A\Zj=\lambda_{\max}\Zj$.}
\begin{proof}
\pf\ \raggedright Divide by $n>0$ in \stepref{sp2-forms} and apply Theorem~\ref{thm:rayleigh}.~\qed
\end{proof}
\step{sp2-degrees}{That equation forces all degrees equal to a natural number $k$.}
\begin{proof}
\pf\ \raggedright Evaluate \stepref{sp2-equality} at each vertex and choose one vertex to set $k=\deg(u)$.~\qed
\end{proof}
\step{sp2-reverse}{If $X$ is $k$-regular, the spectral identity holds.}
\begin{proof}
\pf\ \raggedright Use $A\Zj=k\Zj$ and $\lambda_{\max}=k$ from Theorem~\ref{thm:sp1} in \stepref{sp2-forms}.~\qed
\end{proof}
\qedstep
\begin{proof}
\pf\ \raggedright Use \stepref{sp2-equality}, \stepref{sp2-degrees}, \stepref{sp2-reverse}.~\qed
\end{proof}
\end{proof}
\qedstep
\begin{proof}
\pf\ \raggedright Use \stepref{sp2-Rayleigh}.~\qed
\end{proof}
\end{proof}

\begin{corollary}[Cospectral mates of regular graphs]\label{cor:regular-mates}
\lean{SimpleGraph.Cospectral.isRegularOfDegree, SimpleGraph.Cospectral.connected_of_regular}
\uses{lem:extremes,prop:cospectral,thm:multiplicity,thm:sp1,thm:sp2}
\mbox{}\par\nobreak
\assume{\begin{pfenum}
\item $X$ is a nonempty finite $k$-regular simple graph and $Y$ is a finite simple graph cospectral with $X$.
\end{pfenum}}
\prove{\raggedright $Y$ is $k$-regular, has the same number of components as $X$, and is connected if $X$ is connected.}
\end{corollary}
\begin{proof}
\pfsketch\ We will transfer the order, the second moment and the largest eigenvalue to $Y$, and apply SP2 and SP1 to $Y$.
\step{mates-invariants}{$X,Y$ have equal order $n>0$, equal second moments $\sum_\mu\mu^2$ and equal largest eigenvalues.}
\begin{proof}
\pf\ \raggedright Proposition~\ref{prop:cospectral} gives $\Spec(X)=\Spec(Y)$; the three quantities are functions of the multiset, by Theorem~\ref{thm:multiplicity} and Lemma~\ref{lem:extremes}.~\qed
\end{proof}
\step{mates-regular}{$Y$ is $k$-regular.}
\begin{proof}
\pf\ \raggedright Theorem~\ref{thm:sp2} for $X$ gives $\sum_\mu\mu^2=n\lambda_{\max}$; by \stepref{mates-invariants} this identity holds for $Y$, and Theorem~\ref{thm:sp2} for $Y$ makes $Y$ regular. Theorem~\ref{thm:sp1} for $Y$ and for $X$ identifies both degrees with the common largest eigenvalue $k$.~\qed
\end{proof}
\step{mates-count}{$X$ and $Y$ have the same number of components.}
\begin{proof}
\pf\ \raggedright By Theorem~\ref{thm:sp1} for $X$ and for $Y$, using \stepref{mates-regular}, both numbers equal $\mult(k)$, which is common by Proposition~\ref{prop:cospectral}.~\qed
\end{proof}
\step{mates-connected}{If $X$ is connected, so is $Y$.}
\begin{proof}
\pf\ \raggedright A nonempty graph is connected exactly when it has one component; use \stepref{mates-count}.~\qed
\end{proof}
\qedstep
\begin{proof}
\pf\ \raggedright Use \stepref{mates-regular}, \stepref{mates-count}, \stepref{mates-connected}.~\qed
\end{proof}
\end{proof}

\begin{theorem}[SP3: bipartiteness]\label{thm:sp3}
\lean{SimpleGraph.IsRegularOfDegree.isBipartite_iff_neg_eigenvalue, SimpleGraph.IsBipartite.multiplicity_neg_eq}
\uses{facts:dim,thm:laplacian,thm:multiplicity}
\mbox{}\par\nobreak
\assume{\begin{pfenum}
\item $X$ is a finite connected simple $k$-regular graph, $k\in\nn$.
\item The second assertion concerns any finite simple bipartite graph $Y$ and any $\mu\in\re$.
\end{pfenum}}
\prove{\raggedright It is bipartite iff
$-k\in\Spec(X)$; when this occurs its multiplicity is one.
More generally, for every finite bipartite graph $Y$,
\[
 \mult_Y(\mu)=\mult_Y(-\mu)\qquad(\mu\in\re).
\]}
\end{theorem}
\begin{proof}
\pfsketch\ We will count the signless kernel, then reverse signs on one part of a bipartition.
\step{sp3-kernel}{$X$ is bipartite iff $-k$ occurs, and then its multiplicity is one.}
\begin{proof}
\step{sp3-identify}{$V_A(-k)=\ker(kI+A)$.}
\begin{proof}
\pf\ \raggedright Rearrange the eigenvector equation.~\qed
\end{proof}
\step{sp3-dim}{That kernel has dimension one if $X$ is bipartite and zero otherwise.}
\begin{proof}
\pf\ \raggedright Connectedness leaves one component in Theorem~\ref{thm:laplacian}.~\qed
\end{proof}
\step{sp3-occurrence}{The claimed occurrence and multiplicity follow, including $k=0$.}
\begin{proof}
\pf\ \raggedright Apply Theorem~\ref{thm:multiplicity} to \stepref{sp3-identify}, \stepref{sp3-dim}; no division by $k$ occurs.~\qed
\end{proof}
\qedstep
\begin{proof}
\pf\ \raggedright Use \stepref{sp3-occurrence}.~\qed
\end{proof}
\end{proof}
\step{sp3-sign}{For a finite bipartite graph $Y$, $\mult_Y(\mu)=\mult_Y(-\mu)$.}
\begin{proof}
\pflet{$A\triangleq A(Y)$; fix a bipartition of $Y$, and let $T$ be the diagonal matrix with entry $1$ on the first part and $-1$ on the second.}
\step{sp3-T}{$T^2=I$ and $TAT=-A$.}
\begin{proof}
\step{sp3-square}{$T^2=I$.}
\begin{proof}
\pf\ \raggedright Each diagonal sign squares to $1$, and the off-diagonal entries vanish.~\qed
\end{proof}
\step{sp3-entries}{\assume{$u,v\in V(Y)$.}\prove{$(TAT)(u,v)=-A(u,v)$.}}
\begin{proof}
\step{sp3-diagonal}{\case{$u=v$.}}
\begin{proof}
\pf\ \raggedright Both entries are zero, since $Y$ is loopless and $T$ is diagonal.~\qed
\end{proof}
\step{sp3-edge}{\case{$u\ne v$ and $uv$ is an edge.}}
\begin{proof}
\pf\ \raggedright $(TAT)(u,v)=T(u,u)A(u,v)T(v,v)$, and the edge crosses the two parts, so the sign product is $-1$.~\qed
\end{proof}
\step{sp3-nonedge}{\case{$u\ne v$ and $uv$ is not an edge.}}
\begin{proof}
\pf\ \raggedright Both entries are zero, since $A(u,v)=0$.~\qed
\end{proof}
\qedstep
\begin{proof}
\pf\ \raggedright The three cases are exhaustive.~\qed
\end{proof}
\end{proof}
\qedstep
\begin{proof}
\pf\ \raggedright Use \stepref{sp3-square}, \stepref{sp3-entries}.~\qed
\end{proof}
\end{proof}
\step{sp3-map}{$\Zx\mapsto T\Zx$ maps $V_A(\mu)$ to $V_A(-\mu)$.}
\begin{proof}
\pf\ \raggedright From \stepref{sp3-T}, $AT=-TA$; apply this to $A\Zx=\mu\Zx$.~\qed
\end{proof}
\step{sp3-inverse}{The reverse map is also multiplication by $T$.}
\begin{proof}
\pf\ \raggedright For $A\Zy=-\mu\Zy$, the identity $AT=-TA$ from \stepref{sp3-T} gives $A(T\Zy)=\mu T\Zy$ directly. Both composites are $T^2=I$.~\qed
\end{proof}
\step{sp3-dimensions}{The multiplicities agree.}
\begin{proof}
\pf\ \raggedright The linear equivalence in \stepref{sp3-inverse} preserves dimensions (Facts~\ref{facts:dim}); use Theorem~\ref{thm:multiplicity}.~\qed
\end{proof}
\qedstep
\begin{proof}
\pf\ \raggedright Use \stepref{sp3-dimensions}.~\qed
\end{proof}
\end{proof}
\qedstep
\begin{proof}
\pf\ \raggedright Use \stepref{sp3-kernel}, \stepref{sp3-sign}.~\qed
\end{proof}
\end{proof}

\begin{theorem}[The all-ones direction]\label{thm:ones}
\lean{Matrix.IsHermitian.spectrumMultiset_smul_add_smul_one_add_smul_allOnes, SimpleGraph.IsRegularOfDegree.spectrumMultiset_compl, SimpleGraph.IsRegularOfDegree.spectrumMultiset_seidelMatrix}
\uses{def:graph,facts:coord,facts:graph,facts:order,thm:multiplicity}
\mbox{}\par\nobreak
\assume{\begin{pfenum}
\item $U$ is finite with $|U|=n>0$, $A:U\times U\to\re$ is symmetric, and $A\Zj=k\Zj$ for $k\in\re$.
\item $b,s,t\in\re$. For the complement identity $X$ is any finite simple graph; for the two graph spectra $X$ is a nonempty finite simple $k$-regular graph on $U$ and $A=A(X)$.
\end{pfenum}}
\prove{\raggedright For every $\Zx$, $J\Zx=(\sum_u\Zx(u))\Zj$; thus $J\Zx=0$ on
$\Zj^\perp$. There is an orthogonal decomposition
$\re^U=\operatorname{span}(\Zj)\oplus\Zj^\perp$, of dimensions $1,n-1$.
The orthogonal complement $\Zj^\perp$ is invariant under $A$, so one
can choose an orthonormal eigenbasis starting with $\Zj/\sqrt n$.
Write the remaining eigenvalues, with multiplicity, as $\mu_1,\ldots,
\mu_{n-1}$. For real $b,s,t$, the eigenvalues of $bI+sA+tJ$ are
\[
 b+sk+tn,\qquad b+s\mu_i\quad(1\le i<n).
\]
For every finite simple graph $X$, $A(\comp X)=J-I-A(X)$.
Thus for a nonempty $k$-regular graph,
\[
 \begin{aligned}
 \Spec(\comp X)&=\ms{n-1-k,-1-\mu_1,\ldots,-1-\mu_{n-1}},\\
 \Spec(S(X))&=\ms{n-1-2k,-1-2\mu_1,\ldots,-1-2\mu_{n-1}}.
 \end{aligned}
\]}
\end{theorem}
\begin{proof}
\pfsketch\ We will separate the constant direction and diagonalize on its orthogonal complement.
\step{ones-decompose}{The orthogonal decomposition into $\operatorname{span}(\Zj)$ and $\Zj^\perp$ has dimensions $1,n-1$; the complement is invariant and admits the asserted eigenbasis.}
\begin{proof}
\step{ones-invariant}{For $\ip{\Zx}{\Zj}=0$, $\ip{A\Zx}{\Zj}=0$.}
\begin{proof}
\step{ones-adjoint}{$\ip{A\Zx}{\Zj}=\ip{\Zx}{A\Zj}$.}
\begin{proof}
\pf\ \raggedright Symmetry of $A$, by exchanging its two coordinate summation indices.~\qed
\end{proof}
\step{ones-eigen}{$\ip{\Zx}{A\Zj}=k\ip{\Zx}{\Zj}$.}
\begin{proof}
\pf\ \raggedright Substitute $A\Zj=k\Zj$ and use linearity of the inner product.~\qed
\end{proof}
\step{ones-zero}{$k\ip{\Zx}{\Zj}=0$.}
\begin{proof}
\pf\ \raggedright The assumed inner product is zero.~\qed
\end{proof}
\qedstep
\begin{proof}
\pf\ \raggedright Transitivity of the three equalities.~\qed
\end{proof}
\end{proof}
\step{ones-split}{Every $\Zx$ splits as $\frac{\ip{\Zx}{\Zj}}n\Zj+(\Zx-\frac{\ip{\Zx}{\Zj}}n\Zj)$ into the constant direction and its orthogonal complement.}
\begin{proof}
\pf\ \raggedright The second term has inner product $\ip{\Zx}{\Zj}-\ip{\Zx}{\Zj}=0$ with $\Zj$, since $\norm{\Zj}^2=n$.~\qed
\end{proof}
\step{ones-restrict}{The restriction to $\Zj^\perp$ is symmetric and has an orthonormal eigenbasis.}
\begin{proof}
\pf\ \raggedright Use \stepref{ones-invariant} and the finite-dimensional linear-map spectral theorem in Facts~\ref{facts:coord}.~\qed
\end{proof}
\step{ones-basis}{Adjoining $\Zj/\sqrt n$ gives the asserted basis.}
\begin{proof}
\pf\ \raggedright It has norm one by Facts~\ref{facts:order}; \stepref{ones-split} is an orthogonal direct sum.~\qed
\end{proof}
\qedstep
\begin{proof}
\pf\ \raggedright Use \stepref{ones-invariant}, \stepref{ones-split}, \stepref{ones-basis}.~\qed
\end{proof}
\end{proof}
\step{ones-action}{$J\Zx=(\sum_u\Zx(u))\Zj$, so $J$ kills $\Zj^\perp$, and the displayed list is the spectrum of $bI+sA+tJ$.}
\begin{proof}
\step{ones-J}{$J\Zx=(\sum_u\Zx(u))\Zj$.}
\begin{proof}
\pf\ \raggedright Expand the all-ones matrix product.~\qed
\end{proof}
\step{ones-constant}{On $\Zj/\sqrt n$ the scalar is $b+sk+tn$.}
\begin{proof}
\pf\ \raggedright Apply the three summands and \stepref{ones-J}.~\qed
\end{proof}
\step{ones-perp}{On a $\mu_i$-eigenvector in $\Zj^\perp$ the scalar is $b+s\mu_i$.}
\begin{proof}
\pf\ \raggedright The coordinate sum is zero, so the $J$ term vanishes by \stepref{ones-J}.~\qed
\end{proof}
\step{ones-list}{The list counts every eigenvalue with multiplicity.}
\begin{proof}
\pf\ \raggedright By \stepref{ones-constant} and \stepref{ones-perp}, the basis in \stepref{ones-decompose} is an orthonormal eigenbasis of the symmetric matrix $bI+sA+tJ$ with the displayed eigenvalues; apply the arbitrary-basis assertion of Theorem~\ref{thm:multiplicity}.~\qed
\end{proof}
\qedstep
\begin{proof}
\pf\ \raggedright Use \stepref{ones-J}, \stepref{ones-list}.~\qed
\end{proof}
\end{proof}
\step{ones-graphs}{The complement identity and the complement and Seidel formulas hold.}
\begin{proof}
\step{ones-complement}{$A(\comp X)=J-I-A(X)$ for every finite simple graph $X$.}
\begin{proof}
\pf\ \raggedright Apply \code{SimpleGraph.one_add_adjMatrix_add_compl_adjMatrix_eq_of_one} in Facts~\ref{facts:graph}, which gives $I+A(X)+A(X).\mathrm{compl}=J$, and \code{SimpleGraph.compl_adjMatrix_eq_adjMatrix_compl}, which identifies $A(X).\mathrm{compl}$ with $A(\comp X)$.~\qed
\end{proof}
\step{ones-Seidel}{$S(X)=J-I-2A(X)$.}
\begin{proof}
\pf\ \raggedright This is Definition~\ref{def:graph}.~\qed
\end{proof}
\step{ones-substitute}{For $A=A(X)$ as in \asm[2]{0}, the two lists follow from \stepref{ones-action}.}
\begin{proof}
\pf\ \raggedright By \stepref{ones-complement} and \stepref{ones-Seidel}, $A(\comp X)$ and $S(X)$ are $bI+sA+tJ$ for $(b,s,t)=(-1,-1,1)$ and $(-1,-2,1)$, respectively.~\qed
\end{proof}
\qedstep
\begin{proof}
\pf\ \raggedright Use \stepref{ones-complement}, \stepref{ones-substitute}.~\qed
\end{proof}
\end{proof}
\qedstep
\begin{proof}
\pf\ \raggedright Use \stepref{ones-decompose}, \stepref{ones-action}, \stepref{ones-graphs}.~\qed
\end{proof}
\end{proof}

\begin{corollary}[Removing the constant direction]\label{cor:remove-ones}
\lean{Matrix.PosSemidef.rank_sub_allOnes}
\uses{facts:positive,thm:ones,thm:psd,thm:trace}
\mbox{}\par\nobreak
\assume{\begin{pfenum}
\item $G$ is real positive semidefinite of order $n>0$, $G\Zj=b\Zj$, and $b>0$.
\end{pfenum}}
\prove{\raggedright $G-(b/n)J\succeq0$ and its rank is $\rk G-1$.}
\end{corollary}
\begin{proof}
\pfsketch\ We will apply Theorem~\ref{thm:ones} to $G$; the correction by $J$ replaces the eigenvalue $b$ on the constant direction by zero and leaves the other eigenvalues unchanged.
\step{remove-list}{The eigenvalues of $G-(b/n)J$ are $0$ on $\Zj/\sqrt n$ and the eigenvalues $\mu_1,\ldots,\mu_{n-1}$ of $G$ on $\Zj^\perp$.}
\begin{proof}
\pf\ \raggedright $G$ is symmetric (Facts~\ref{facts:positive}); apply Theorem~\ref{thm:ones} to $G$ with $k=b$ and coefficients $0,1,-b/n$, so that $0+b-(b/n)n=0$.~\qed
\end{proof}
\step{remove-positive}{$G-(b/n)J\succeq0$.}
\begin{proof}
\pf\ \raggedright By Theorem~\ref{thm:ones} for $G$ with coefficients $0,1,0$, $\Spec(G)=\ms{b,\mu_1,\ldots,\mu_{n-1}}$, so each $\mu_i\ge0$ by Theorem~\ref{thm:psd}; with \stepref{remove-list}, Theorem~\ref{thm:psd} applies to $G-(b/n)J$.~\qed
\end{proof}
\step{remove-rank}{$\rk(G-(b/n)J)=\rk G-1$.}
\begin{proof}
\pf\ \raggedright By Theorem~\ref{thm:trace}, $\rk G$ counts the nonzero entries of $\ms{b,\mu_1,\ldots,\mu_{n-1}}$, and $\rk(G-(b/n)J)$ counts those of \stepref{remove-list}; the two lists differ only in $b>0$.~\qed
\end{proof}
\qedstep
\begin{proof}
\pf\ \raggedright Use \stepref{remove-positive}, \stepref{remove-rank}.~\qed
\end{proof}
\end{proof}

\begin{corollary}[A scalar matrix plus the all-ones matrix]\label{cor:scalar-ones}
\lean{Matrix.spectrumMultiset_smul_one_add_smul_allOnes, Matrix.charpoly_smul_one_add_smul_allOnes, Matrix.det_smul_one_add_smul_allOnes}
\uses{facts:rank-one,thm:multiplicity,thm:ones,thm:trace}
\mbox{}\par\nobreak
\assume{\begin{pfenum}
\item $U$ is finite with $|U|=n>0$, $b,t\in\re$, and $I,J$ are the identity and all-ones matrices on $U$.
\end{pfenum}}
\prove{\raggedright \[
 \begin{aligned}
 \Spec(bI+tJ)&=\ms{(b+tn)^1,b^{n-1}},\\
 \chi_{bI+tJ}(z)&=(z-b-tn)(z-b)^{n-1},\\
 \det(bI+tJ)&=(b+tn)b^{n-1}.
 \end{aligned}
\]
The multiset convention combines the two values when $t=0$.}
\end{corollary}
\begin{proof}
\pfsketch\ We will compute on the constant direction and its orthogonal complement.
\step{scalar-actions}{The two actions have scalars $b+tn$ and $b$, in dimensions one and $n-1$.}
\begin{proof}
\step{scalar-constant}{$(bI+tJ)\Zj=(b+tn)\Zj$.}
\begin{proof}
\pf\ \raggedright $J\Zj=n\Zj$ by the defining row sum.~\qed
\end{proof}
\step{scalar-perp}{For $\Zx\perp\Zj$, $(bI+tJ)\Zx=b\Zx$.}
\begin{proof}
\pf\ \raggedright $J\Zx=0$ by the coordinate-sum formula asserted in Theorem~\ref{thm:ones}, applied to the zero matrix with $k=0$.~\qed
\end{proof}
\step{scalar-dim}{The two orthogonal spaces have dimensions $1,n-1$.}
\begin{proof}
\pf\ \raggedright Use the orthogonal decomposition asserted in Theorem~\ref{thm:ones}, applied to the zero matrix with $k=0$.~\qed
\end{proof}
\qedstep
\begin{proof}
\pf\ \raggedright Use \stepref{scalar-constant}, \stepref{scalar-perp}, \stepref{scalar-dim}.~\qed
\end{proof}
\end{proof}
\step{scalar-polynomials}{All three formulas, including coincident values and $n=1$, hold.}
\begin{proof}
\step{scalar-multiset}{$\Spec(bI+tJ)=\ms{(b+tn)^1,b^{n-1}}$.}
\begin{proof}
\pf\ \raggedright The vector $\Zj/\sqrt n$ and an orthonormal basis of $\Zj^\perp$ form an orthonormal eigenbasis of $bI+tJ$ with eigenvalues $b+tn$ and $b$, by \stepref{scalar-actions}; apply the arbitrary-basis assertion of Theorem~\ref{thm:multiplicity}.~\qed
\end{proof}
\step{scalar-charpoly}{$\chi(z)=(z-b-tn)(z-b)^{n-1}$.}
\begin{proof}
\pf\ \raggedright Apply the product formula in Theorem~\ref{thm:multiplicity} to \stepref{scalar-multiset}.~\qed
\end{proof}
\step{scalar-det}{$\det(bI+tJ)=(b+tn)b^{n-1}$.}
\begin{proof}
\step{scalar-det-nonzero}{\case{$b\ne0$.}}
\begin{proof}
\pf\ \raggedright Apply \code{Matrix.det_one_add_replicateCol_mul_replicateRow} in Facts~\ref{facts:rank-one} to the constant vectors $(t/b)\Zj$ and $\Zj$: $\det(I+(t/b)J)=1+tn/b$. Then $\det(bI+tJ)=b^n(1+tn/b)=(b+tn)b^{n-1}$, since $n>0$.~\qed
\end{proof}
\step{scalar-det-zero}{\case{$b=0$.}}
\begin{proof}
\step{scalar-det-one}{\case{$n=1$.}}
\begin{proof}
\pf\ \raggedright By Theorem~\ref{thm:trace} and \stepref{scalar-multiset}, the determinant is the single eigenvalue $tn$, which is $(b+tn)b^0$.~\qed
\end{proof}
\step{scalar-det-many}{\case{$n>1$.}}
\begin{proof}
\pf\ \raggedright By Theorem~\ref{thm:trace} and \stepref{scalar-multiset}, the determinant is a product with a factor $b=0$; the right side $(b+tn)b^{n-1}$ is also zero.~\qed
\end{proof}
\qedstep
\begin{proof}
\pf\ \raggedright The cases $n=1$ and $n>1$ exhaust $n>0$.~\qed
\end{proof}
\end{proof}
\qedstep
\begin{proof}
\pf\ \raggedright The cases $b\ne0$ and $b=0$ are exhaustive.~\qed
\end{proof}
\end{proof}
\qedstep
\begin{proof}
\pf\ \raggedright Use \stepref{scalar-multiset}, \stepref{scalar-charpoly}, \stepref{scalar-det}.~\qed
\end{proof}
\end{proof}
\qedstep
\begin{proof}
\pf\ \raggedright Use \stepref{scalar-polynomials}.~\qed
\end{proof}
\end{proof}

\section{Rectangular products and line graphs}

\begin{theorem}[Rectangular characteristic polynomials]\label{thm:mn}
\lean{Matrix.eigenvalueMultiplicity_mul_comm_of_ne_zero, Matrix.spectrumMultiset_gram_products}
\uses{def:spectrum,facts:charpoly,facts:roots,thm:gram,thm:psd,thm:trace}
\mbox{}\par\nobreak
\assume{\begin{pfenum}
\item $P,Q$ are finite with cardinalities $p,q$; $M:P\times Q\to R$ and $N:Q\times P\to R$, where $R$ is a commutative ring.
\item For the multiplicity and eigenspace assertions $R=\re$ and $\mu\in\re\setminus\{0\}$. For the final assertion $B:P\times Q\to\re$ has rank $\rho$.
\end{pfenum}}
\prove{\raggedright \[
 t^q\chi_{MN}(t)=t^p\chi_{NM}(t).
\]
Over $\re$, the nonzero roots have equal algebraic multiplicities and
\[
 q+\mult_{MN}(0)=p+\mult_{NM}(0).
\]
For $\mu\ne0$, multiplication by $N$ is a linear equivalence
$V_{MN}(\mu)\simeq V_{NM}(\mu)$, with inverse $\mu^{-1}M$.
If $B$ is a real $p\times q$ matrix of rank $\rho$, both $BB^T$ and
$B^TB$ are positive semidefinite; their positive spectra agree with
multiplicity, and their zero multiplicities are $p-\rho$ and $q-\rho$.}
\end{theorem}
\begin{proof}
\pfsketch\ We will compare two block determinants, transfer eigenspaces, and count the zero modes of Gram products.
\step{mn-polynomial}{$t^q\chi_{MN}=t^p\chi_{NM}$.}
\begin{proof}
\pf\ \raggedright Apply \code{Matrix.charpoly_mul_comm'} in Facts~\ref{facts:charpoly} over the commutative coefficient ring.~\qed
\end{proof}
\step{mn-counts}{Nonzero algebraic multiplicities agree and $q+\mult_{MN}(0)=p+\mult_{NM}(0)$.}
\begin{proof}
\step{mn-nonzero-polys}{Both characteristic polynomials and the powers of $t$ are nonzero.}
\begin{proof}
\pf\ \raggedright Characteristic polynomials are monic (Facts~\ref{facts:charpoly}); powers of $t$ over $\re$ are nonzero.~\qed
\end{proof}
\step{mn-addmult}{Root multiplicities in \stepref{mn-polynomial} add over the two factors on each side.}
\begin{proof}
\pf\ \raggedright Apply the product-root multiplicity fact in Facts~\ref{facts:roots}, using \stepref{mn-nonzero-polys}.~\qed
\end{proof}
\step{mn-powerroots}{$t^h$ has multiplicity $h$ at zero and zero at every nonzero number.}
\begin{proof}
\pf\ \raggedright Its roots are $h$ copies of the unique root zero of $t$.~\qed
\end{proof}
\step{mn-bothcounts}{The two claimed equations hold.}
\begin{proof}
\pf\ \raggedright Substitute \stepref{mn-powerroots} in \stepref{mn-addmult} and use Definition~\ref{def:spectrum}.~\qed
\end{proof}
\qedstep
\begin{proof}
\pf\ \raggedright Use \stepref{mn-bothcounts}.~\qed
\end{proof}
\end{proof}
\step{mn-equiv}{$N$ and $\mu^{-1}M$ restrict to inverse linear maps of the two $\mu$-eigenspaces.}
\begin{proof}
\step{mn-forward}{$NM(N\Zx)=\mu N\Zx$ when $MN\Zx=\mu\Zx$.}
\begin{proof}
\pf\ \raggedright Associativity gives $NM(N\Zx)=N(MN\Zx)$; substitute the hypothesis.~\qed
\end{proof}
\step{mn-backward}{$MN(\mu^{-1}M\Zy)=\mu(\mu^{-1}M\Zy)$ when $NM\Zy=\mu\Zy$.}
\begin{proof}
\pf\ \raggedright Commute the scalar with $M$, use associativity, and substitute $NM\Zy=\mu\Zy$.~\qed
\end{proof}
\step{mn-leftinv}{$\mu^{-1}MN\Zx=\Zx$ on $V_{MN}(\mu)$.}
\begin{proof}
\pf\ \raggedright Cancel $\mu\ne0$ in its eigenvector equation.~\qed
\end{proof}
\step{mn-rightinv}{$N(\mu^{-1}M\Zy)=\Zy$ on $V_{NM}(\mu)$.}
\begin{proof}
\pf\ \raggedright Cancel $\mu\ne0$ in the $NM$ eigenvector equation.~\qed
\end{proof}
\qedstep
\begin{proof}
\pf\ \raggedright Use \stepref{mn-forward}, \stepref{mn-backward}, \stepref{mn-leftinv}, \stepref{mn-rightinv}.~\qed
\end{proof}
\end{proof}
\step{mn-gram}{The Gram products have common positive spectrum and zero multiplicities $p-\rho,q-\rho$.}
\begin{proof}
\step{mn-grampos}{$BB^T,B^TB\succeq0$ and both ranks are $\rho$.}
\begin{proof}
\pf\ \raggedright Theorem~\ref{thm:gram}.~\qed
\end{proof}
\step{mn-zeros}{Their zero multiplicities are their orders minus $\rho$.}
\begin{proof}
\pf\ \raggedright Theorem~\ref{thm:trace} applied to \stepref{mn-grampos}.~\qed
\end{proof}
\step{mn-positives}{All remaining eigenvalues are positive, with equal counts.}
\begin{proof}
\pf\ \raggedright Use Theorem~\ref{thm:psd} for positivity. Apply \code{Matrix.charpoly_mul_comm'} directly to $B,B^T$ (Facts~\ref{facts:charpoly}); at nonzero roots the factors $t^p,t^q$ have zero multiplicity, by the product rule in Facts~\ref{facts:roots}.~\qed
\end{proof}
\qedstep
\begin{proof}
\pf\ \raggedright Use \stepref{mn-zeros}, \stepref{mn-positives}.~\qed
\end{proof}
\end{proof}
\qedstep
\begin{proof}
\pf\ \raggedright Use \stepref{mn-polynomial}, \stepref{mn-counts}, \stepref{mn-equiv}, \stepref{mn-gram}.~\qed
\end{proof}
\end{proof}

\begin{theorem}[Incidence identities]\label{thm:incidence}
\lean{SimpleGraph.edgeIncMatrix, SimpleGraph.edgeIncMatrix_mul_transpose, SimpleGraph.transpose_edgeIncMatrix_mul, SimpleGraph.multiplicity_neg_two_lineGraph}
\uses{cor:minus-two,facts:dim,facts:graph,facts:incidence,thm:calculus,thm:gram,thm:laplacian,thm:mn}
\mbox{}\par\nobreak
\assume{\begin{pfenum}
\item $X$ is a finite simple graph with vertex set $U$, $|U|=n$, and $|E(X)|=m$.
\item $M:U\times E(X)\to\re$ is the endpoint incidence matrix, with actual edges as columns; $D$ is the diagonal degree matrix.
\item $b$ is the number of bipartite connected components of $X$.
\end{pfenum}}
\prove{\raggedright \[
 MM^T=D+A(X),\qquad M^TM=2I+A(L(X)).
\]
If $b$ is the number of bipartite connected components of $X$, then
\[
 \rk M=n-b,\qquad
 \mult_{L(X)}(-2)=m-n+b.
\]
The right side is nonnegative. Every line graph has least eigenvalue
at least $-2$ whenever it is nonempty.}
\end{theorem}
\begin{proof}
\pfsketch\ We will count common incidences in both products and apply the signless kernel formula.
\step{inc-first}{$MM^T=D+A(X)$.}
\begin{proof}
\step{inc-library}{Mathlib's full incidence product has degree diagonal and adjacency off the diagonal.}
\begin{proof}
\pf\ \raggedright Apply \code{SimpleGraph.incMatrix_mul_transpose} in Facts~\ref{facts:incidence} over $\re$.~\qed
\end{proof}
\step{inc-restrict}{Restricting its columns to actual edges leaves the product unchanged.}
\begin{proof}
\pf\ \raggedright For every nonedge column all incidence entries are zero; remove precisely these zero summands from each product entry.~\qed
\end{proof}
\step{inc-identify}{The resulting product is $D+A(X)$.}
\begin{proof}
\pf\ \raggedright The diagonal in \stepref{inc-library} is $D$, the off-diagonal is $A$, and \stepref{inc-restrict} identifies the restricted product.~\qed
\end{proof}
\qedstep
\begin{proof}
\pf\ \raggedright Use \stepref{inc-identify}.~\qed
\end{proof}
\end{proof}
\step{inc-second}{$M^TM=2I+A(L(X))$.}
\begin{proof}
\step{inc-edgediag}{$(M^TM)(e,e)=2$.}
\begin{proof}
\pf\ \raggedright Each simple edge has exactly two distinct endpoints.~\qed
\end{proof}
\step{inc-edgeoff}{For $e\ne f$, $(M^TM)(e,f)$ is $1$ if they intersect and $0$ if they do not.}
\begin{proof}
\pf\ \raggedright Distinct unordered two-element sets have intersection of size at most one.~\qed
\end{proof}
\step{inc-linedef}{These off-diagonal entries equal line-graph adjacency.}
\begin{proof}
\pf\ \raggedright This is \code{SimpleGraph.lineGraph} in Facts~\ref{facts:graph}.~\qed
\end{proof}
\qedstep
\begin{proof}
\pf\ \raggedright Use \stepref{inc-edgediag}, \stepref{inc-edgeoff}, \stepref{inc-linedef}.~\qed
\end{proof}
\end{proof}
\step{inc-rank}{$\rk M=n-b$ and $\mult_{L(X)}(-2)=m-n+b\ge0$.}
\begin{proof}
\step{inc-kernel}{$\dim\ker(MM^T)=b$.}
\begin{proof}
\pf\ \raggedright Use \stepref{inc-first} and Theorem~\ref{thm:laplacian}.~\qed
\end{proof}
\step{inc-rankM}{$\rk M=\rk(MM^T)=n-b\le m$.}
\begin{proof}
\pf\ \raggedright Rank-nullity (Facts~\ref{facts:dim}) and Theorem~\ref{thm:gram}.~\qed
\end{proof}
\step{inc-zeros}{$M^TM$ has zero multiplicity $m-(n-b)$.}
\begin{proof}
\pf\ \raggedright Theorem~\ref{thm:mn} and \stepref{inc-rankM}.~\qed
\end{proof}
\step{inc-shift}{That multiplicity equals $\mult_{L(X)}(-2)$.}
\begin{proof}
\pf\ \raggedright Use \stepref{inc-second} and Theorem~\ref{thm:calculus} with the shift by two.~\qed
\end{proof}
\qedstep
\begin{proof}
\pf\ \raggedright Use \stepref{inc-rankM}, \stepref{inc-zeros}, \stepref{inc-shift}.~\qed
\end{proof}
\end{proof}
\step{inc-lower}{Every nonempty line graph has least eigenvalue at least $-2$.}
\begin{proof}
\pf\ \raggedright $M^TM\succeq0$ by Theorem~\ref{thm:gram}; apply Corollary~\ref{cor:minus-two} to \stepref{inc-second}.~\qed
\end{proof}
\qedstep
\begin{proof}
\pf\ \raggedright Use \stepref{inc-first}, \stepref{inc-second}, \stepref{inc-rank}, \stepref{inc-lower}.~\qed
\end{proof}
\end{proof}

\begin{lemma}[Connectivity of a line graph]\label{lem:line-connected}
\lean{SimpleGraph.lineGraph_connected_iff_of_no_isolated, SimpleGraph.degree_lineGraph}
\uses{facts:graph}
\mbox{}\par\nobreak
\assume{\begin{pfenum}
\item $X$ is a finite simple graph with at least one edge and no isolated vertices.
\item $\{u,v\}$ is an edge for the degree assertion.
\end{pfenum}}
\prove{\raggedright $X$ is connected iff $L(X)$ is connected. For any edge
$\{u,v\}$, its degree in the line graph is $\deg_X(u)+\deg_X(v)-2$.}
\end{lemma}
\begin{proof}
\pfsketch\ We will turn vertex paths into edge walks and edge walks into reachability, then count the two disjoint sets of neighbouring edges.
\step{line-forward}{Connectedness of $X$ implies connectedness of $L(X)$.}
\begin{proof}
\step{line-sequence}{Between any two edges, a path between chosen endpoints gives an edge sequence starting and ending with those edges.}
\begin{proof}
\pf\ \raggedright Connectedness supplies the endpoint path; prepend and append the chosen edges.~\qed
\end{proof}
\step{line-walk}{Consecutive distinct members of this sequence are adjacent in $L(X)$.}
\begin{proof}
\pf\ \raggedright They share the intervening endpoint; use the line-graph definition (Facts~\ref{facts:graph}).~\qed
\end{proof}
\step{line-delete}{Removing consecutive repetitions gives a walk between the chosen edges.}
\begin{proof}
\pf\ \raggedright Each retained transition is one of the adjacencies in \stepref{line-walk}; if all entries coincide, the zero-length walk suffices.~\qed
\end{proof}
\qedstep
\begin{proof}
\pf\ \raggedright The line graph has a vertex because $X$ has an edge, and \stepref{line-delete} proves every pair reachable.~\qed
\end{proof}
\end{proof}
\step{line-back}{Connectedness of $L(X)$ implies connectedness of $X$.}
\begin{proof}
\step{line-incident}{Given vertices $u,v$, choose incident edges $e,f$.}
\begin{proof}
\pf\ \raggedright No vertex is isolated.~\qed
\end{proof}
\step{line-chain}{There is a sequence $e=e_0,\ldots,e_h=f$ whose consecutive edges intersect.}
\begin{proof}
\pf\ \raggedright Connectedness gives a walk in $L(X)$ and the definition gives the intersections.~\qed
\end{proof}
\step{line-propagate}{Both endpoints of every $e_i$ are reachable from $u$.}
\begin{proof}
\pf\ \raggedright Induct on $i$: the base uses the edge at $u$; an intersection vertex and the next edge give the inductive extension.~\qed
\end{proof}
\step{line-reachable}{$v$ is reachable from $u$.}
\begin{proof}
\pf\ \raggedright It is an endpoint of $e_h$ in \stepref{line-propagate}.~\qed
\end{proof}
\qedstep
\begin{proof}
\pf\ \raggedright Use \stepref{line-reachable}.~\qed
\end{proof}
\end{proof}
\step{line-degree}{$\deg_{L(X)}(\{u,v\})=\deg_X(u)+\deg_X(v)-2$.}
\begin{proof}
\step{line-union}{Its neighbours are the other edges at $u$, together with the other edges at $v$.}
\begin{proof}
\pf\ \raggedright A line-graph neighbour shares one endpoint with $\{u,v\}$.~\qed
\end{proof}
\step{line-disjoint}{Those two sets are disjoint.}
\begin{proof}
\pf\ \raggedright An edge containing both $u,v$ is the original edge, which was removed.~\qed
\end{proof}
\step{line-count}{The two cardinalities are $\deg(u)-1$ and $\deg(v)-1$.}
\begin{proof}
\pf\ \raggedright Remove the single original edge from each incident-edge set; both degrees are positive.~\qed
\end{proof}
\qedstep
\begin{proof}
\pf\ \raggedright Add the disjoint cardinalities in \stepref{line-union}, \stepref{line-disjoint}, \stepref{line-count}.~\qed
\end{proof}
\end{proof}
\qedstep
\begin{proof}
\pf\ \raggedright Use \stepref{line-forward}, \stepref{line-back}, \stepref{line-degree}.~\qed
\end{proof}
\end{proof}

\begin{theorem}[SP4: regular line graphs]\label{thm:sp4}
\lean{SimpleGraph.regular_or_semiregular_of_lineGraph_regular}
\uses{lem:line-connected}
\mbox{}\par\nobreak
\assume{\begin{pfenum}
\item $X$ is a finite simple graph with at least one edge and no isolated vertices.
\item $L(X)$ is connected and $h$-regular, with $h\in\nn$.
\end{pfenum}}
\prove{\raggedright $X$ is either regular or bipartite
with constant degree on each part. In the latter case, writing the
part sizes as $p,q$ and the respective degrees as $k_1,k_2$, we have
$k_1+k_2=h+2$ and $pk_1=qk_2=|E(X)|$.}
\end{theorem}
\begin{proof}
\pfsketch\ We will propagate the sum of endpoint degrees along paths and split according to whether its two alternating values agree.
\step{sp4-alternate}{$X$ is connected and every degree is one of two positive integers $k_1,k_2$ with $k_1+k_2=h+2$.}
\begin{proof}
\step{sp4-connected}{$X$ is connected and $\deg(u)+\deg(v)=h+2$ on every edge.}
\begin{proof}
\pf\ \raggedright Apply Lemma~\ref{lem:line-connected}.~\qed
\end{proof}
\step{sp4-choose}{\pick\ an edge $u_0v_0$, and let $k_1=\deg(u_0)$, $k_2=\deg(v_0)$; they are positive and $k_1+k_2=h+2$.}
\begin{proof}
\pf\ \raggedright The edge hypothesis supplies $u_0v_0$; its endpoints have positive degree, and \stepref{sp4-connected} applies to it.~\qed
\end{proof}
\step{sp4-path}{Along a path from $u_0$, the degrees alternate $k_1,k_2$.}
\begin{proof}
\pf\ \raggedright Induct on its length: the edge equation in \stepref{sp4-connected} determines the next degree as $h+2$ minus the current one.~\qed
\end{proof}
\step{sp4-all}{Every vertex has degree $k_1$ or $k_2$.}
\begin{proof}
\pf\ \raggedright Connectedness in \stepref{sp4-connected} supplies a path from $u_0$; apply \stepref{sp4-path}.~\qed
\end{proof}
\qedstep
\begin{proof}
\pf\ \raggedright Use \stepref{sp4-connected}, \stepref{sp4-choose}, \stepref{sp4-all}.~\qed
\end{proof}
\end{proof}
\step{sp4-cases}{$X$ is regular, or bipartite with degrees $k_1,k_2$ on parts of sizes $p,q$, where $k_1+k_2=h+2$ and $pk_1=qk_2=|E(X)|$.}
\begin{proof}
\step{sp4-equal}{\case{$k_1=k_2$.}}
\begin{proof}
\pf\ \raggedright Every degree has that value by \stepref{sp4-alternate}, so $X$ is regular.~\qed
\end{proof}
\step{sp4-unequal}{\case{$k_1\ne k_2$.}}
\begin{proof}
\step{sp4-partition}{The two degree classes form a bipartition of $X$ into nonempty parts, of sizes $p,q$ say.}
\begin{proof}
\pf\ \raggedright The endpoint equation in \stepref{sp4-alternate} forces opposite degree classes across every edge, since $k_1\ne k_2$; the edge $u_0v_0$ meets both classes.~\qed
\end{proof}
\step{sp4-firstcount}{$pk_1=|E(X)|$.}
\begin{proof}
\pf\ \raggedright By \stepref{sp4-partition}, each edge has exactly one endpoint in the first part; count the edges at the vertices of that part.~\qed
\end{proof}
\step{sp4-secondcount}{$qk_2=|E(X)|$.}
\begin{proof}
\pf\ \raggedright Count the edges at the vertices of the second part in the same way.~\qed
\end{proof}
\qedstep
\begin{proof}
\pf\ \raggedright Use \stepref{sp4-partition}, \stepref{sp4-firstcount}, \stepref{sp4-secondcount}, and $k_1+k_2=h+2$ from \stepref{sp4-alternate}.~\qed
\end{proof}
\end{proof}
\qedstep
\begin{proof}
\pf\ \raggedright The cases $k_1=k_2$ and $k_1\ne k_2$ are exhaustive.~\qed
\end{proof}
\end{proof}
\qedstep
\begin{proof}
\pf\ \raggedright Use \stepref{sp4-cases}.~\qed
\end{proof}
\end{proof}

\begin{theorem}[Recovering a regular graph's spectrum]\label{thm:line-regular}
\lean{SimpleGraph.spectrumMultiset_of_regular_lineGraph, SimpleGraph.cospectral_of_lineGraph_cospectral}
\uses{lem:line-connected,prop:cospectral,thm:calculus,thm:incidence,thm:mn,thm:multiplicity,thm:sp1,thm:walks}
\mbox{}\par\nobreak
\assume{\begin{pfenum}
\item $X$ is a finite simple $k$-regular graph with at least one edge; $n=|V(X)|$, $m=|E(X)|$.
\item $h=\lambda_{\max}(L(X))$ and $\rho=\rk(2I+A(L(X)))$.
\item For the last assertion $X'$ is a second finite simple regular graph with at least one edge, and $L(X')$ is cospectral with $L(X)$.
\end{pfenum}}
\prove{\raggedright the spectrum of $L(X)$ determines $k,n$
and the full spectrum of $X$. In detail, with
$h=\lambda_{\max}(L(X))$ and $\rho=\rk(2I+A(L(X)))$,
\[
 k=(h+2)/2,\qquad n=2m/k,
\]
and $\Spec(X)$ is obtained by taking all nonzero eigenvalues of
$2I+A(L(X))$, adjoining $n-\rho$ zeros, and subtracting $k$ from every
entry. Equivalently, without negative exponents,
\[
 (t+2)^n\chi_{L(X)}(t)=(t+2)^m\chi_X(t+2-k).
\]
Thus two regular graphs with an edge and cospectral line
graphs are cospectral.}
\end{theorem}
\begin{proof}
\pfsketch\ We will recover the degree from line-graph regularity and then add precisely the missing zero modes of the incidence product. We will prove this for an arbitrary regular graph $Y$ with an edge, so that the reconstruction can be applied to both $X$ and $X'$.
\step{recover-general}{\assume{$Y$ is a finite simple $k_Y$-regular graph with at least one edge, $n_Y=|V(Y)|$, $m_Y=|E(Y)|$, $h_Y=\lambda_{\max}(L(Y))$ and $\rho_Y=\rk(2I+A(L(Y)))$.}\prove{$m_Y=|\Spec(L(Y))|$, $k_Y=(h_Y+2)/2$, $n_Y=2m_Y/k_Y$, $\Spec(Y)$ is obtained from $\Spec(L(Y))$ by the stated reconstruction, and $(t+2)^{n_Y}\chi_{L(Y)}(t)=(t+2)^{m_Y}\chi_Y(t+2-k_Y)$.}}
\begin{proof}
\step{recover-parameters}{$m_Y=|\Spec(L(Y))|$, $k_Y=(h_Y+2)/2>0$ and $n_Y=2m_Y/k_Y$.}
\begin{proof}
\step{recover-kpos}{$k_Y>0$ and $Y$ has no isolated vertex.}
\begin{proof}
\pf\ \raggedright An edge of \asm{1} has an endpoint of positive degree, and all degrees equal $k_Y$.~\qed
\end{proof}
\step{recover-h}{$h_Y=2k_Y-2$.}
\begin{proof}
\pf\ \raggedright Lemma~\ref{lem:line-connected} makes $L(Y)$ $(2k_Y-2)$-regular, and $L(Y)$ is nonempty; Theorem~\ref{thm:sp1} gives its largest eigenvalue.~\qed
\end{proof}
\step{recover-m}{$m_Y=|\Spec(L(Y))|$ and $n_Yk_Y=2m_Y$.}
\begin{proof}
\pf\ \raggedright Theorem~\ref{thm:multiplicity} counts the vertices of $L(Y)$; Theorem~\ref{thm:walks} gives the handshake identity.~\qed
\end{proof}
\step{recover-formulas}{$k_Y=(h_Y+2)/2$ and $n_Y=2m_Y/k_Y$.}
\begin{proof}
\pf\ \raggedright Solve \stepref{recover-h}, \stepref{recover-m}, using \stepref{recover-kpos}.~\qed
\end{proof}
\qedstep
\begin{proof}
\pf\ \raggedright Use \stepref{recover-kpos}, \stepref{recover-m}, \stepref{recover-formulas}.~\qed
\end{proof}
\end{proof}
\step{recover-products}{For the incidence matrix $M$ of $Y$, $MM^T=k_YI+A(Y)$ and $M^TM=2I+A(L(Y))$.}
\begin{proof}
\pf\ \raggedright Theorem~\ref{thm:incidence} and regularity.~\qed
\end{proof}
\step{recover-spectrum}{$\Spec(Y)$ is obtained by taking the nonzero eigenvalues of $2I+A(L(Y))$, adjoining $n_Y-\rho_Y$ zeros, and subtracting $k_Y$.}
\begin{proof}
\step{recover-nonzero}{$k_YI+A(Y)$ has the nonzero spectrum of $2I+A(L(Y))$ and $n_Y-\rho_Y$ zeros.}
\begin{proof}
\pf\ \raggedright Apply Theorem~\ref{thm:mn} to $M$; by \stepref{recover-products} the rank of both products is $\rho_Y$.~\qed
\end{proof}
\step{recover-shift}{Subtracting $k_Y$ from every entry yields $\Spec(Y)$.}
\begin{proof}
\pf\ \raggedright Apply Theorem~\ref{thm:calculus} to $k_YI+A(Y)$ and the polynomial $t-k_Y$, using \stepref{recover-nonzero}.~\qed
\end{proof}
\qedstep
\begin{proof}
\pf\ \raggedright Use \stepref{recover-shift}.~\qed
\end{proof}
\end{proof}
\step{recover-substitute}{$(t+2)^{m_Y}\chi_Y(t+2-k_Y)=(t+2)^{n_Y}\chi_{L(Y)}(t)$.}
\begin{proof}
\pf\ \raggedright Apply Theorem~\ref{thm:mn} to $M$ and $M^T$, whose products are given by \stepref{recover-products}, and substitute $t+2$ for the indeterminate.~\qed
\end{proof}
\qedstep
\begin{proof}
\pf\ \raggedright Use \stepref{recover-parameters}, \stepref{recover-spectrum}, \stepref{recover-substitute}.~\qed
\end{proof}
\end{proof}
\step{recover-X}{The displayed formulas and the reconstruction hold for $X$.}
\begin{proof}
\pf\ \raggedright Apply \stepref{recover-general} to $Y=X$ with \asm[1]{0} and \asm[2]{0}.~\qed
\end{proof}
\step{recover-mates}{$X$ and $X'$ are cospectral.}
\begin{proof}
\pf\ \raggedright Apply \stepref{recover-general} to $Y=X$ and to $Y=X'$. By \asm[3]{0} and Proposition~\ref{prop:cospectral}, $L(X)$ and $L(X')$ have the same spectrum, so the quantities $m$, $h$, $\rho$, and thus the degree, the order and the reconstructed multiset, agree; the last is $\Spec(X)=\Spec(X')$. Apply Proposition~\ref{prop:cospectral} once more.~\qed
\end{proof}
\qedstep
\begin{proof}
\pf\ \raggedright Use \stepref{recover-X}, \stepref{recover-mates}.~\qed
\end{proof}
\end{proof}

\begin{theorem}[Bipartite block spectra]\label{thm:block}
\lean{Matrix.spectrumMultiset_fromBlocks_scalar, Matrix.fromBlocks_scalar_quadratic}
\uses{facts:dim,facts:order,facts:rank,thm:gram,thm:mn,thm:multiplicity}
\mbox{}\par\nobreak
\assume{\begin{pfenum}
\item $P,Q$ are finite of sizes $p,q$; $B:P\times Q\to\re$ has rank $\rho$.
\item $\lambda_1,\ldots,\lambda_\rho$ are the positive eigenvalues of $BB^T$, with repetition, and $s,t\in\re$.
\end{pfenum}}
\prove{\raggedright the matrix
\[
 C=\pmat{sI_p&B\\B^T&tI_q}
\]
has the two eigenvalues
\[
 \gamma_j^{\pm}=\frac{s+t\pm\sqrt{(s-t)^2+4\lambda_j}}2
 \quad(1\le j\le \rho),
\]
as well as $s$ repeated $p-\rho$ times and $t$ repeated $q-\rho$ times.
This is a multiset assertion, so coinciding displayed values add their
multiplicities. For $s=t=0$, it gives $\pm\sqrt{\lambda_j}$ and
$p+q-2\rho$ zeros. Furthermore
\[
 (C-sI)(C-tI)=\pmat{BB^T&0\\0&B^TB}.
\]}
\end{theorem}
\begin{proof}
\pfsketch\ We will pair positive Gram eigendirections into orthogonal two-dimensional spaces and keep both kernel spaces.
\step{block-pairs}{\pick\ orthonormal $\Zu_j$ with $BB^T\Zu_j=\lambda_j\Zu_j$, and put $\Zv_j=B^T\Zu_j/\sqrt{\lambda_j}$. Then $(\Zv_j)$ is orthonormal and $B\Zv_j=\sqrt{\lambda_j}\Zu_j$.}
\begin{proof}
\step{block-basis}{\pick\ the $\Zu_j$ as the positive-eigenvalue vectors of an orthonormal eigenbasis of $BB^T$.}
\begin{proof}
\pf\ \raggedright Theorem~\ref{thm:multiplicity} applied to the symmetric Gram matrix.~\qed
\end{proof}
\step{block-positive}{Each $\sqrt{\lambda_j}>0$ and its square is $\lambda_j$.}
\begin{proof}
\pf\ \raggedright Facts~\ref{facts:order} and $\lambda_j>0$.~\qed
\end{proof}
\step{block-inner}{$\ip{\Zv_i}{\Zv_j}=\lambda_j\ip{\Zu_i}{\Zu_j}/(\sqrt{\lambda_i}\sqrt{\lambda_j})$.}
\begin{proof}
\pf\ \raggedright Move $B^T$ across the inner product and substitute the eigenvector equation.~\qed
\end{proof}
\step{block-orthonormal}{This inner product is $1$ for $i=j$ and $0$ for $i\ne j$.}
\begin{proof}
\pf\ \raggedright Use the orthonormality of $(\Zu_j)$ and \stepref{block-positive} in \stepref{block-inner}.~\qed
\end{proof}
\step{block-Baction}{$B\Zv_j=\sqrt{\lambda_j}\Zu_j$ and $B^T\Zu_j=\sqrt{\lambda_j}\Zv_j$.}
\begin{proof}
\pf\ \raggedright The first is the $BB^T$ equation divided by $\sqrt{\lambda_j}$; the second is the definition of $\Zv_j$.~\qed
\end{proof}
\qedstep
\begin{proof}
\pf\ \raggedright Use \stepref{block-basis}, \stepref{block-orthonormal}, \stepref{block-Baction}.~\qed
\end{proof}
\end{proof}
\step{block-decomp}{The ambient space is the orthogonal sum of the $\rho$ paired planes and kernel spaces of dimensions $p-\rho,q-\rho$.}
\begin{proof}
\step{block-leftkernel}{The remaining left eigenbasis vectors span $\ker B^T$ of dimension $p-\rho$.}
\begin{proof}
\pf\ \raggedright Theorem~\ref{thm:gram} identifies $\ker BB^T=\ker B^T$; Theorem~\ref{thm:mn} counts its zero modes.~\qed
\end{proof}
\step{block-image}{The $\rho$ vectors $\Zv_j$ form a basis of $\im B^T$.}
\begin{proof}
\pf\ \raggedright They are independent by \stepref{block-pairs}, belong to the image, and that image has dimension $\rho$ by Facts~\ref{facts:rank}.~\qed
\end{proof}
\step{block-rightkernel}{The orthogonal complement of this image is $\ker B$, of dimension $q-\rho$.}
\begin{proof}
\pf\ \raggedright $\ip{B^T\Zu}{\Zv}=\ip{\Zu}{B\Zv}$ for every $\Zu$ makes the two membership conditions equivalent; apply rank-nullity (Facts~\ref{facts:dim}).~\qed
\end{proof}
\step{block-total}{The dimensions add to $2\rho+(p-\rho)+(q-\rho)=p+q$.}
\begin{proof}
\pf\ \raggedright Finite-dimensional arithmetic; the summands are orthogonal by \stepref{block-pairs}, \stepref{block-leftkernel}, \stepref{block-rightkernel}.~\qed
\end{proof}
\qedstep
\begin{proof}
\pf\ \raggedright Use \stepref{block-leftkernel}, \stepref{block-image}, \stepref{block-rightkernel}, \stepref{block-total}.~\qed
\end{proof}
\end{proof}
\step{block-roots}{The stated multiset gives all eigenvalues of $C$.}
\begin{proof}
\step{block-plane}{On the $j$th paired plane the matrix is $\pmat{s&\sqrt{\lambda_j}\\\sqrt{\lambda_j}&t}$.}
\begin{proof}
\pf\ \raggedright Apply the four blocks of $C$ using \stepref{block-pairs}.~\qed
\end{proof}
\step{block-quadratic}{Its characteristic polynomial is $(z-s)(z-t)-\lambda_j$.}
\begin{proof}
\pf\ \raggedright Use the two-by-two determinant formula on \stepref{block-plane}.~\qed
\end{proof}
\step{block-discriminant}{The two roots are $\gamma_j^\pm$ and are distinct.}
\begin{proof}
\pf\ \raggedright Complete the square: $(2z-s-t)^2=(s-t)^2+4\lambda_j>0$; use Facts~\ref{facts:order}.~\qed
\end{proof}
\step{block-kernels}{The two kernel spaces have scalars $s$ and $t$, respectively.}
\begin{proof}
\step{block-leftaction}{For $B^T\Zu=0$, $C(\Zu,0)=(s\Zu,0)$.}
\begin{proof}
\pf\ \raggedright The two block rows give $s\Zu+B0=s\Zu$ and $B^T\Zu+t0=0$.~\qed
\end{proof}
\step{block-rightaction}{For $B\Zv=0$, $C(0,\Zv)=(0,t\Zv)$.}
\begin{proof}
\pf\ \raggedright The two block rows give $s0+B\Zv=0$ and $B^T0+t\Zv=t\Zv$.~\qed
\end{proof}
\qedstep
\begin{proof}
\pf\ \raggedright Use \stepref{block-leftaction}, \stepref{block-rightaction}.~\qed
\end{proof}
\end{proof}
\step{block-account}{All eigenvalues and their multiplicities are accounted for.}
\begin{proof}
\pf\ \raggedright Take the union of bases of the invariant orthogonal summands in \stepref{block-decomp}; by \stepref{block-discriminant} and \stepref{block-kernels} it is an orthonormal eigenbasis of the symmetric matrix $C$ with the stated eigenvalues. Apply the arbitrary-basis assertion of Theorem~\ref{thm:multiplicity}.~\qed
\end{proof}
\qedstep
\begin{proof}
\pf\ \raggedright Use \stepref{block-account}.~\qed
\end{proof}
\end{proof}
\step{block-product}{$(C-sI)(C-tI)=BB^T\oplus B^TB$; for $s=t=0$ the spectrum specializes as stated.}
\begin{proof}
\step{block-offdiag}{Both off-diagonal blocks of $C^2-(s+t)C+stI$ vanish.}
\begin{proof}
\pf\ \raggedright They are $(s+t)B-(s+t)B$ and its transpose.~\qed
\end{proof}
\step{block-diag}{The diagonal blocks are $BB^T$ and $B^TB$.}
\begin{proof}
\pf\ \raggedright The scalar terms are $s^2-(s+t)s+st=0$ and $t^2-(s+t)t+st=0$.~\qed
\end{proof}
\step{block-zero}{For $s=t=0$ the paired roots are $\pm\sqrt{\lambda_j}$ and the other $p+q-2\rho$ entries are zero.}
\begin{proof}
\pf\ \raggedright Substitute in \stepref{block-roots}, using $\sqrt{4\lambda_j}=2\sqrt{\lambda_j}$.~\qed
\end{proof}
\qedstep
\begin{proof}
\pf\ \raggedright Use \stepref{block-offdiag}, \stepref{block-diag}, \stepref{block-zero}.~\qed
\end{proof}
\end{proof}
\qedstep
\begin{proof}
\pf\ \raggedright Use \stepref{block-roots}, \stepref{block-product}.~\qed
\end{proof}
\end{proof}

\begin{theorem}[Cospectral line graphs of bipartite graphs]\label{thm:line-bipartite}
\lean{SimpleGraph.lineGraph_cospectral_bipartite}
\uses{cor:regular-mates,lem:line-connected,prop:cospectral,thm:block,thm:calculus,thm:gram,thm:incidence,thm:laplacian,thm:mn,thm:multiplicity,thm:sp4,thm:trace}
\mbox{}\par\nobreak
\assume{\begin{pfenum}
\item $X_1$ is a finite connected simple bipartite graph with an edge, and $L(X_1)$ is regular.
\item $X_2$ is a finite simple graph with no isolated vertices and $L(X_2)$ cospectral with $L(X_1)$.
\end{pfenum}}
\prove{\raggedright $X_2$ is
connected, and precisely one of the following alternatives holds:
\begin{enumerate}
\item $X_2$ is bipartite. Order the part sizes $p\le q$ in each
graph. The part sizes and respective degrees $k_1,k_2$ agree, and
the two biadjacency matrices $B_1,B_2$ have cospectral products
$B_1B_1^T,B_2B_2^T$.
\item $X_1$ is not regular; $X_2$ is regular and not bipartite;
$|V(X_1)|=|V(X_2)|+1$.
\end{enumerate}}
\end{theorem}
\begin{proof}
\pfsketch\ We will recover the order from the signless rank, recover the two degrees from their sum and product, and distinguish the bipartite and nonbipartite cases.
\step{lb-common}{Both line graphs are connected and $h$-regular with the same $m>0$ vertices; $X_2$ is connected and is regular or bipartite semiregular.}
\begin{proof}
\step{lb-lineconnected}{$L(X_1)$ is connected.}
\begin{proof}
\pf\ \raggedright $X_1$ is connected with an edge, so has no isolates; apply Lemma~\ref{lem:line-connected}.~\qed
\end{proof}
\step{lb-transfer}{$L(X_2)$ is connected, regular with the same degree $h$, and has the same order $m>0$.}
\begin{proof}
\pf\ \raggedright Apply Corollary~\ref{cor:regular-mates} and Proposition~\ref{prop:cospectral}.~\qed
\end{proof}
\step{lb-root}{$X_2$ is connected and has the required regular-or-semiregular alternative.}
\begin{proof}
\pf\ \raggedright It has $m>0$ edges and no isolates; apply Lemma~\ref{lem:line-connected} and Theorem~\ref{thm:sp4}.~\qed
\end{proof}
\qedstep
\begin{proof}
\pf\ \raggedright Use \stepref{lb-lineconnected}, \stepref{lb-transfer}, \stepref{lb-root}.~\qed
\end{proof}
\end{proof}
\step{lb-rank}{$2I+A(L(X_1))$ and $2I+A(L(X_2))$ have the same rank and the same nonzero spectrum.}
\begin{proof}
\pf\ \raggedright By \stepref{lb-common} and Proposition~\ref{prop:cospectral} the two line graphs have the same spectrum; shift it by two with Theorem~\ref{thm:calculus} and count nonzero entries with Theorem~\ref{thm:trace}.~\qed
\end{proof}
\pflet{$s\triangleq h+2$, where $h$ is the common degree in \stepref{lb-common}; $R$ is the common rank in \stepref{lb-rank}, and $N\triangleq R+1$.}
\step{lb-order}{\assume{$Z$ is $X_1$, or $Z=X_2$ and $X_2$ is bipartite.}\prove{$|V(Z)|=N$, and the signless Laplacian of $Z$ has the common nonzero spectrum of \stepref{lb-rank} with one zero adjoined.}}
\begin{proof}
\step{lb-onezero}{The signless Laplacian of $Z$ has exactly one zero eigenvalue.}
\begin{proof}
\pf\ \raggedright $Z$ is connected, by \asm[1]{0} or \stepref{lb-common}, and bipartite; apply Theorem~\ref{thm:laplacian} and Theorem~\ref{thm:multiplicity}.~\qed
\end{proof}
\step{lb-N}{$|V(Z)|=R+1$, and the nonzero part of the signless spectrum of $Z$ is that of \stepref{lb-rank}.}
\begin{proof}
\pf\ \raggedright By Theorem~\ref{thm:incidence}, the signless Laplacian and $2I+A(L(Z))$ are the two Gram products of the incidence matrix of $Z$; Theorem~\ref{thm:mn} gives equal nonzero spectra, and with \stepref{lb-onezero} the order is the rank plus one.~\qed
\end{proof}
\qedstep
\begin{proof}
\pf\ \raggedright Use \stepref{lb-onezero}, \stepref{lb-N}.~\qed
\end{proof}
\end{proof}
\step{lb-degrees}{\assume{$Z$ is $X_1$, or $Z=X_2$ and $X_2$ is bipartite.}\prove{$Z$ has parts of sizes $p\le q$ with constant degrees $k_1\ge k_2>0$, where $k_1,k_2$ are the ordered roots of $z^2-sz+ms/N$, $p=m/k_1$ and $q=m/k_2$.}}
\begin{proof}
\step{lb-semireg}{The degrees of $Z$ are constant on the two parts; write them $k_1\ge k_2>0$ on parts of sizes $p\le q$.}
\begin{proof}
\pf\ \raggedright $L(Z)$ is connected and $h$-regular by \asm[1]{0}, \stepref{lb-common}; apply Theorem~\ref{thm:sp4}, and in its regular case note that each part has that common degree. An edge gives positive degrees.~\qed
\end{proof}
\step{lb-count}{$p+q=N$, $k_1+k_2=s$ and $pk_1=qk_2=m$.}
\begin{proof}
\pf\ \raggedright Theorem~\ref{thm:sp4} gives the last two, since $Z$ has $m$ edges by \stepref{lb-common}; \stepref{lb-order} gives the first.~\qed
\end{proof}
\step{lb-product}{$k_1k_2=ms/N$.}
\begin{proof}
\pf\ \raggedright From \stepref{lb-count}, $N=m/k_1+m/k_2=ms/(k_1k_2)$; all denominators are positive.~\qed
\end{proof}
\step{lb-unique}{$k_1,k_2$ are the ordered roots of $z^2-sz+ms/N$, and $p=m/k_1$, $q=m/k_2$.}
\begin{proof}
\pf\ \raggedright Use their sum and product in \stepref{lb-count}, \stepref{lb-product}; ordering fixes which root is first.~\qed
\end{proof}
\qedstep
\begin{proof}
\pf\ \raggedright Use \stepref{lb-semireg}, \stepref{lb-unique}.~\qed
\end{proof}
\end{proof}
\step{lb-bipartite}{\case{$X_2$ is bipartite.}}
\begin{proof}
\step{lb-C}{The signless Laplacians of $X_1$ and $X_2$ are $C_i=\pmat{k_1I&B_i\\B_i^T&k_2I}$ with common $p,q,k_1,k_2$ and common spectrum.}
\begin{proof}
\pf\ \raggedright Apply \stepref{lb-degrees} and \stepref{lb-order} to $Z=X_1$ and $Z=X_2$, using \asm{1}; order the vertices by parts.~\qed
\end{proof}
\step{lb-map}{$B_1B_1^T\oplus B_1^TB_1$ and $B_2B_2^T\oplus B_2^TB_2$ have the same spectrum.}
\begin{proof}
\pf\ \raggedright By Theorem~\ref{thm:block}, $(C_i-k_1I)(C_i-k_2I)$ is this direct sum; apply Theorem~\ref{thm:calculus} with $f(z)=(z-k_1)(z-k_2)$ to \stepref{lb-C}.~\qed
\end{proof}
\step{lb-halving}{Each positive multiplicity in $B_iB_i^T$ is half the corresponding multiplicity in the direct sum.}
\begin{proof}
\pf\ \raggedright Theorem~\ref{thm:mn} makes the positive counts of $B_iB_i^T$ and $B_i^TB_i$ equal; direct-sum eigenbases add counts.~\qed
\end{proof}
\step{lb-zeros}{$B_1B_1^T$ and $B_2B_2^T$ have the same zero multiplicity.}
\begin{proof}
\pf\ \raggedright Each has order $p$; subtract the common sum of positive multiplicities in \stepref{lb-map}, \stepref{lb-halving} from $p$.~\qed
\end{proof}
\step{lb-cospectral}{$B_1B_1^T$ and $B_2B_2^T$ are cospectral.}
\begin{proof}
\pf\ \raggedright Apply Proposition~\ref{prop:cospectral} to \stepref{lb-halving}, \stepref{lb-zeros}.~\qed
\end{proof}
\qedstep
\begin{proof}
\pf\ \raggedright Alternative (1) holds by \stepref{lb-C} and \stepref{lb-cospectral}; alternative (2) fails, since it requires $X_2$ not to be bipartite.~\qed
\end{proof}
\end{proof}
\step{lb-nonbipartite}{\case{$X_2$ is not bipartite.}}
\begin{proof}
\step{lb-regular}{$X_2$ is regular of degree $s/2$.}
\begin{proof}
\step{lb-sp4alt}{$X_2$ is regular, or bipartite with constant degree on each part.}
\begin{proof}
\pf\ \raggedright This is \stepref{lb-common}.~\qed
\end{proof}
\step{lb-exclude-semiregular}{\case{$X_2$ is bipartite with constant degree on each part.}}
\begin{proof}
\pf\ \raggedright This case assumption contradicts \asm{1}, so the case does not occur.~\qed
\end{proof}
\step{lb-regular-branch}{\case{$X_2$ is regular, of degree $k$.}}
\begin{proof}
\pf\ \raggedright Lemma~\ref{lem:line-connected} gives line degree $2k-2=h$, and $s=h+2$, so $k=s/2$.~\qed
\end{proof}
\qedstep
\begin{proof}
\pf\ \raggedright The two cases of \stepref{lb-sp4alt} are exhaustive.~\qed
\end{proof}
\end{proof}
\step{lb-nonsingular}{$|V(X_2)|=R=N-1$.}
\begin{proof}
\pf\ \raggedright $X_2$ is connected by \stepref{lb-common} and not bipartite, so Theorem~\ref{thm:laplacian} gives a zero signless kernel. By Theorems~\ref{thm:incidence} and \ref{thm:gram}, the rank of the incidence matrix of $X_2$ is $|V(X_2)|$ and is the rank $R$ of $2I+A(L(X_2))$.~\qed
\end{proof}
\step{lb-X1notregular}{$X_1$ is not regular.}
\begin{proof}
\pf\ \raggedright If it were, of degree $k$, then $2k-2=h$ by Lemma~\ref{lem:line-connected}, so $k=s/2>0$, which is also the degree of $X_2$ by \stepref{lb-regular}; the handshake identities $Nk=2m=(N-1)k$ for $X_1$ and $X_2$, from \stepref{lb-order} and \stepref{lb-nonsingular}, would contradict $k>0$.~\qed
\end{proof}
\qedstep
\begin{proof}
\pf\ \raggedright Alternative (2) holds by \stepref{lb-regular}, \stepref{lb-nonsingular}, \stepref{lb-X1notregular}; alternative (1) fails, since it requires $X_2$ to be bipartite.~\qed
\end{proof}
\end{proof}
\qedstep
\begin{proof}
\pf\ \raggedright $X_2$ is connected by \stepref{lb-common}. The cases \stepref{lb-bipartite} and \stepref{lb-nonbipartite} are exhaustive, and in each exactly one alternative holds.~\qed
\end{proof}
\end{proof}

\begin{corollary}[The regular nonbipartite alternative]\label{cor:rival}
\lean{SimpleGraph.lineGraph_cospectral_bipartite_spectrum}
\uses{cor:positive-certificate,lem:line-connected,thm:block,thm:incidence,thm:laplacian,thm:line-bipartite,thm:mn,thm:perron}
\mbox{}\par\nobreak
\assume{\begin{pfenum}
\item $X_1$ is a finite connected simple bipartite semiregular graph with part sizes $p\le q$ and degrees $k_1>k_2>0$; $B$ is its biadjacency matrix of rank $\rho$.
\item $X_2$ is finite, simple, connected, regular and nonbipartite; $L(X_2)$ is cospectral with $L(X_1)$.
\item $s=k_1+k_2$.
\end{pfenum}}
\prove{\raggedright The positive eigenvalues of $BB^T$ include $k_1k_2$ once. Delete this
one occurrence and call the remaining list $\lambda_2,\ldots,\lambda_\rho$.
Then the spectrum of $X_2$ is
\[
 \ms{\frac s2,\quad
 \pm\sqrt{\lambda_j+\frac{(k_1-k_2)^2}{4}}\ (2\le j\le \rho),
 \quad\left(\frac{k_1-k_2}{2}\right)^{p-\rho},
 \quad\left(-\frac{k_1-k_2}{2}\right)^{q-\rho}}.
\]}
\end{corollary}
\begin{proof}
\pfsketch\ We will identify the single endpoint pair $0,s$, delete its zero member before shifting, and retain all rectangular kernel modes.
\step{rival-endpoint}{$k_1k_2$ is a positive eigenvalue of $BB^T$ of multiplicity one.}
\begin{proof}
\step{rival-rowcols}{$B\Zj_q=k_1\Zj_p$ and $B^T\Zj_p=k_2\Zj_q$.}
\begin{proof}
\pf\ \raggedright These are the constant row and column sums.~\qed
\end{proof}
\step{rival-gramtop}{$BB^T\Zj_p=k_1k_2\Zj_p$.}
\begin{proof}
\pf\ \raggedright Apply $B$ to the second equation in \stepref{rival-rowcols} and use the first.~\qed
\end{proof}
\step{rival-Cvector}{For $C=D(X_1)+A(X_1)$, $(k_1\Zj_p,k_2\Zj_q)$ is a positive $s$-eigenvector.}
\begin{proof}
\pf\ \raggedright Block multiplication with \stepref{rival-rowcols} gives scalars $k_1^2+k_1k_2=sk_1$ and $k_1k_2+k_2^2=sk_2$.~\qed
\end{proof}
\step{rival-simple}{$s$ is a simple eigenvalue of $C$.}
\begin{proof}
\pf\ \raggedright Its support is the connected graph $X_1$; use Corollary~\ref{cor:positive-certificate} and Theorem~\ref{thm:perron}.~\qed
\end{proof}
\step{rival-pair}{Each copy of $k_1k_2$ gives one $0$ and one $s$ in $C$.}
\begin{proof}
\pf\ \raggedright In Theorem~\ref{thm:block}, $(z-k_1)(z-k_2)-k_1k_2=z(z-s)$.~\qed
\end{proof}
\step{rival-one}{There is exactly one such copy.}
\begin{proof}
\pf\ \raggedright It occurs by \stepref{rival-gramtop}; two copies would contradict \stepref{rival-simple} through \stepref{rival-pair}.~\qed
\end{proof}
\qedstep
\begin{proof}
\pf\ \raggedright Use \stepref{rival-gramtop}, \stepref{rival-one}.~\qed
\end{proof}
\end{proof}
\step{rival-delete}{The signless spectrum of $X_2$ is the spectrum of $C$ with its single zero deleted.}
\begin{proof}
\step{rival-zero}{$C$ has exactly one zero eigenvalue.}
\begin{proof}
\pf\ \raggedright Theorem~\ref{thm:laplacian} for connected bipartite $X_1$.~\qed
\end{proof}
\step{rival-nonzero}{The signless matrix of $X_2$ has no zero eigenvalue and has the same nonzero spectrum as $C$.}
\begin{proof}
\pf\ \raggedright Use nonbipartiteness, Theorem~\ref{thm:laplacian}, and the incidence Gram transfer from Theorems~\ref{thm:incidence}, \ref{thm:mn}.~\qed
\end{proof}
\qedstep
\begin{proof}
\pf\ \raggedright Use \stepref{rival-zero}, \stepref{rival-nonzero}.~\qed
\end{proof}
\end{proof}
\step{rival-shift}{Subtracting $s/2$ gives exactly the stated adjacency multiset.}
\begin{proof}
\step{rival-degree}{$X_2$ has degree $s/2$.}
\begin{proof}
\pf\ \raggedright Lemma~\ref{lem:line-connected} gives degree $s-2$ for $L(X_1)$. Apply Theorem~\ref{thm:line-bipartite}; $X_1$ has an edge since $k_2>0$, and connected nonbipartite $X_2$ has no isolates.~\qed
\end{proof}
\step{rival-top}{The surviving endpoint eigenvalue $s$ becomes $s/2$.}
\begin{proof}
\pf\ \raggedright Use \stepref{rival-delete} and subtract the degree in \stepref{rival-degree}.~\qed
\end{proof}
\step{rival-paired}{Each remaining pair becomes $\pm\sqrt{\lambda_j+(k_1-k_2)^2/4}$.}
\begin{proof}
\pf\ \raggedright Subtract $s/2$ from the two roots in Theorem~\ref{thm:block} and factor $4$ inside the square root.~\qed
\end{proof}
\step{rival-left}{The $p-\rho$ left-kernel modes become $(k_1-k_2)/2$.}
\begin{proof}
\pf\ \raggedright Subtract $s/2$ from the scalar $k_1$ in Theorem~\ref{thm:block}.~\qed
\end{proof}
\step{rival-right}{The $q-\rho$ right-kernel modes become $-(k_1-k_2)/2$.}
\begin{proof}
\pf\ \raggedright Subtract $s/2$ from the scalar $k_2$ in Theorem~\ref{thm:block}.~\qed
\end{proof}
\qedstep
\begin{proof}
\pf\ \raggedright Use \stepref{rival-degree}, \stepref{rival-top}, \stepref{rival-paired}, \stepref{rival-left}, \stepref{rival-right}.~\qed
\end{proof}
\end{proof}
\qedstep
\begin{proof}
\pf\ \raggedright Use \stepref{rival-endpoint}, \stepref{rival-shift}.~\qed
\end{proof}
\end{proof}

Error 14 of \code{cgss/errata.md}~\cite{errata} records the endpoint deletion and the two kernel multiplicities in the preceding corollary.

\section{Strongly regular graphs and design spectra}

\begin{theorem}[SR1: parameter and matrix identities]\label{thm:sr1}
\lean{SimpleGraph.IsSRGWith.matrix_eq_integer, SimpleGraph.IsSRGWith.param_eq_integer}
\uses{facts:graph,facts:srg,thm:ones}
\mbox{}\par\nobreak
\assume{\begin{pfenum}
\item $X$ is a finite simple strongly regular graph with parameters $(n,k,a,c)\in\nn^4$: $n$ vertices, degree $k$, and $a,c$ common neighbours for adjacent and distinct nonadjacent pairs.
\item $A=A(X)$. The parameter equation assumes $n>0$; the final bounds assume connectedness and $0<k<n-1$.
\end{pfenum}}
\prove{\raggedright \[
 A^2=(k-c)I+(a-c)A+cJ,\qquad AJ=JA=kJ.
\]
If $n>0$, then $k(k-a-1)=c(n-k-1)$ as an equality of integers.
If $X$ is connected and $0<k<n-1$, then $0<c\le k$ and $a\le k-1$.}
\end{theorem}
\begin{proof}
\pfsketch\ We will count entries, apply the identity to the constant vector, and exhibit pairs on which the neighbour counts are attained.
\step{sr1-matrices}{$A^2=(k-c)I+(a-c)A+cJ$ and $AJ=JA=kJ$.}
\begin{proof}
\step{sr1-library}{$A^2=kI+aA+cA(\comp X)$.}
\begin{proof}
\pf\ \raggedright Apply \code{SimpleGraph.IsSRGWith.matrix_eq} in Facts~\ref{facts:srg} over $\re$; natural scalar multiplication becomes multiplication by the real cast.~\qed
\end{proof}
\step{sr1-rearrange}{$A^2=(k-c)I+(a-c)A+cJ$.}
\begin{proof}
\pf\ \raggedright Substitute the complement identity $A(\comp X)=J-I-A$ asserted in Theorem~\ref{thm:ones} into \stepref{sr1-library} and collect terms.~\qed
\end{proof}
\step{sr1-J}{$AJ=JA=kJ$.}
\begin{proof}
\pf\ \raggedright Apply the regular constant-vector declaration in Facts~\ref{facts:graph} to every column of $J$; transpose using symmetry for $JA$.~\qed
\end{proof}
\qedstep
\begin{proof}
\pf\ \raggedright Use \stepref{sr1-rearrange}, \stepref{sr1-J}.~\qed
\end{proof}
\end{proof}
\step{sr1-a}{If $k>0$, then $a\le k-1$.}
\begin{proof}
\pf\ \raggedright An edge $uv$ exists since $k>0$; its $a$ common neighbours lie in $N(u)\setminus\{v\}$, which has $k-1$ elements.~\qed
\end{proof}
\step{sr1-parameter}{For $n>0$, $k(k-a-1)=c(n-k-1)$ in $\ints$.}
\begin{proof}
\step{sr1-param-natural}{In $\nn$, $k(k-a-1)=(n-k-1)c$ with truncated subtraction.}
\begin{proof}
\pf\ \raggedright Apply \code{SimpleGraph.IsSRGWith.param_eq} in Facts~\ref{facts:srg} with $n>0$.~\qed
\end{proof}
\step{sr1-degreebound}{$k\le n-1$.}
\begin{proof}
\pf\ \raggedright A vertex exists since $n>0$; its $k$ neighbours belong to the other $n-1$ vertices.~\qed
\end{proof}
\step{sr1-param-zero}{\case{$k=0$.}}
\begin{proof}
\pf\ \raggedright The left side is zero over both rings. By \stepref{sr1-degreebound}, the right natural difference casts to $n-k-1$; cast \stepref{sr1-param-natural}.~\qed
\end{proof}
\step{sr1-param-positive}{\case{$k>0$.}}
\begin{proof}
\pf\ \raggedright By \stepref{sr1-a} and \stepref{sr1-degreebound}, both subtractions in \stepref{sr1-param-natural} cast without truncation; commute the right factors.~\qed
\end{proof}
\qedstep
\begin{proof}
\pf\ \raggedright The cases $k=0$ and $k>0$ exhaust $k\in\nn$.~\qed
\end{proof}
\end{proof}
\step{sr1-bounds}{\assume{$X$ is connected and $0<k<n-1$.}\prove{$0<c\le k$ and $a\le k-1$.}}
\begin{proof}
\step{sr1-bounds-a}{$a\le k-1$.}
\begin{proof}
\pf\ \raggedright Apply \stepref{sr1-a}, since $k>0$ by \asm{1}.~\qed
\end{proof}
\step{sr1-cupper}{$c\le k$.}
\begin{proof}
\pf\ \raggedright A distinct nonadjacent pair exists since $k<n-1$; its common-neighbour set lies in a neighbourhood of size $k$.~\qed
\end{proof}
\step{sr1-shortpath}{There is a path $u_0u_1u_2$ with $u_0,u_2$ distinct and nonadjacent.}
\begin{proof}
\pf\ \raggedright \pick\ a distinct nonadjacent pair, as $k<n-1$, and a shortest path between them, as $X$ is connected; its length is at least two, and an edge between its first and third vertices would shorten it.~\qed
\end{proof}
\step{sr1-cpositive}{$c\ge1$.}
\begin{proof}
\pf\ \raggedright The vertex $u_1$ in \stepref{sr1-shortpath} is a common neighbour of $u_0$ and $u_2$.~\qed
\end{proof}
\qedstep
\begin{proof}
\pf\ \raggedright Use \stepref{sr1-bounds-a}, \stepref{sr1-cupper}, \stepref{sr1-cpositive}.~\qed
\end{proof}
\end{proof}
\qedstep
\begin{proof}
\pf\ \raggedright Use \stepref{sr1-matrices}, \stepref{sr1-parameter}, \stepref{sr1-bounds}.~\qed
\end{proof}
\end{proof}

\begin{theorem}[SR3: spectrum and integral multiplicities]\label{thm:sr3}
\lean{SimpleGraph.IsSRGWith.spectrumMultiset_eq, SimpleGraph.IsSRGWith.spectralIdempotents, SimpleGraph.IsSRGWith.multiplicity_eq}
\uses{facts:order,thm:calculus,thm:multiplicity,thm:ones,thm:sp1,thm:sr1,thm:trace,thm:walks}
\mbox{}\par\nobreak
\assume{\begin{pfenum}
\item $X$ is a finite connected simple strongly regular graph with parameters $(n,k,a,c)\in\nn^4$ and $0<k<n-1$.
\item $A=A(X)$, $\Delta=(a-c)^2+4(k-c)$, and $\theta_1=(a-c+\sqrt\Delta)/2$, $\theta_2=(a-c-\sqrt\Delta)/2$.
\end{pfenum}}
\prove{\raggedright $\Delta>0$, $k>\theta_1>\theta_2$, and
\[
 \Spec(X)=\ms{k^1,\theta_1^{m_1},\theta_2^{m_2}},\qquad
 m_1=\frac{-k-(n-1)\theta_2}{\theta_1-\theta_2},\quad
 m_2=\frac{k+(n-1)\theta_1}{\theta_1-\theta_2}.
\]
Both displayed multiplicities are positive integers. Its spectral
idempotents are
\[
 E_0=\frac1nJ,\quad
 E_1=\frac{A-\theta_2I-\frac{k-\theta_2}{n}J}{\theta_1-\theta_2},\quad
 E_2=\frac{A-\theta_1I-\frac{k-\theta_1}{n}J}{\theta_2-\theta_1}.
\]
In particular $\rk E_i=\tr E_i=m_i$ for $i=1,2$,
$1+m_1+m_2=n$, and $k+m_1\theta_1+m_2\theta_2=0$.}
\end{theorem}
\begin{proof}
\pfsketch\ We will find the two roots on $\Zj^\perp$, force both to occur, solve for their dimensions, and check each idempotent on each eigenspace.
\step{sr3-roots}{$\Delta>0$ and every eigenvalue on $\Zj^\perp$ is $\theta_1$ or $\theta_2$, with $\theta_1>\theta_2$.}
\begin{proof}
\step{sr3-nonnegative}{$\Delta\ge0$.}
\begin{proof}
\pf\ \raggedright Theorem~\ref{thm:sr1} gives $k-c\ge0$ and the other term is a square.~\qed
\end{proof}
\step{sr3-nonzero}{$\Delta\ne0$.}
\begin{proof}
\pf\ \raggedright If the two nonnegative terms sum to zero, then $a=c=k$, contradicting $a\le k-1$ in Theorem~\ref{thm:sr1}.~\qed
\end{proof}
\step{sr3-polynomial}{On $\Zj^\perp$, $A^2-(a-c)A-(k-c)I=0$.}
\begin{proof}
\pf\ \raggedright $J$ kills this subspace by the coordinate-sum formula of Theorem~\ref{thm:ones}, applied to $A$ with $A\Zj=k\Zj$ from Theorem~\ref{thm:sr1}; substitute in the matrix identity of Theorem~\ref{thm:sr1}.~\qed
\end{proof}
\step{sr3-solve}{An eigenvalue there is one of $\theta_1,\theta_2$, which are distinct and ordered.}
\begin{proof}
\pf\ \raggedright Apply \stepref{sr3-polynomial} to its nonzero eigenvector, then complete the square and use $\sqrt\Delta>0$ (Facts~\ref{facts:order}).~\qed
\end{proof}
\qedstep
\begin{proof}
\pf\ \raggedright Use \stepref{sr3-nonnegative}, \stepref{sr3-nonzero}, \stepref{sr3-solve}.~\qed
\end{proof}
\end{proof}
\step{sr3-occurrence}{Both roots occur, and $k>\theta_1>\theta_2$ with $k$ simple.}
\begin{proof}
\step{sr3-top}{$k$ is the simple largest eigenvalue.}
\begin{proof}
\pf\ \raggedright Theorem~\ref{thm:sp1} and connectedness.~\qed
\end{proof}
\step{sr3-onevalue}{If just one value $b$ occurred on $\Zj^\perp$, then $A=bI+(k-b)J/n$.}
\begin{proof}
\pf\ \raggedright The two sides act identically on the eigenbasis from Theorem~\ref{thm:ones}.~\qed
\end{proof}
\step{sr3-twoentries}{That is impossible: there is an off-diagonal entry $1$ and another $0$.}
\begin{proof}
\pf\ \raggedright $k>0$ gives an edge, $k<n-1$ gives a distinct nonedge, while \stepref{sr3-onevalue} makes all off-diagonal entries equal.~\qed
\end{proof}
\step{sr3-strict}{Both roots are strictly below $k$.}
\begin{proof}
\pf\ \raggedright They occur in the perpendicular subspace by \stepref{sr3-roots}, \stepref{sr3-twoentries}; \stepref{sr3-top} leaves no copy of $k$ there.~\qed
\end{proof}
\qedstep
\begin{proof}
\pf\ \raggedright Use \stepref{sr3-top}, \stepref{sr3-twoentries}, \stepref{sr3-strict}.~\qed
\end{proof}
\end{proof}
\step{sr3-multiplicities}{The displayed $m_1,m_2$ are positive integers and give the stated spectrum.}
\begin{proof}
\pflet{$m_i\triangleq\dim V_A(\theta_i)$ for $i=1,2$; these are the multiplicities, by Theorem~\ref{thm:multiplicity}.}
\step{sr3-dimcounts}{$m_1,m_2>0$ and $m_1+m_2=n-1$.}
\begin{proof}
\pf\ \raggedright The eigenbasis of Theorem~\ref{thm:ones} has $\Zj/\sqrt n$ and $n-1$ perpendicular vectors, whose eigenvalues are $\theta_1$ or $\theta_2$ by \stepref{sr3-roots}; both values occur by \stepref{sr3-occurrence}. The arbitrary-basis assertion of Theorem~\ref{thm:multiplicity} makes $m_i$ the number of these vectors with eigenvalue $\theta_i$, since $k\ne\theta_i$.~\qed
\end{proof}
\step{sr3-moment}{$m_1\theta_1+m_2\theta_2=-k$.}
\begin{proof}
\pf\ \raggedright The trace is zero by Theorem~\ref{thm:walks}; expand it with Theorem~\ref{thm:trace}.~\qed
\end{proof}
\step{sr3-m1}{$m_1=(-k-(n-1)\theta_2)/(\theta_1-\theta_2)$.}
\begin{proof}
\pf\ \raggedright Substitute $m_2=n-1-m_1$ in \stepref{sr3-moment} and divide by $\theta_1-\theta_2>0$.~\qed
\end{proof}
\step{sr3-m2}{$m_2=(k+(n-1)\theta_1)/(\theta_1-\theta_2)$.}
\begin{proof}
\pf\ \raggedright Substitute $m_1=n-1-m_2$ in \stepref{sr3-moment} and divide by $\theta_1-\theta_2>0$.~\qed
\end{proof}
\qedstep
\begin{proof}
\pf\ \raggedright Use \stepref{sr3-dimcounts}, \stepref{sr3-m1}, \stepref{sr3-m2}.~\qed
\end{proof}
\end{proof}
\step{sr3-idempotents}{The three displayed matrices are the spectral idempotents, with the asserted traces and ranks.}
\begin{proof}
\step{sr3-E0}{$J/n$ acts by $1,0,0$ on the $k,\theta_1,\theta_2$ eigenspaces.}
\begin{proof}
\pf\ \raggedright Theorem~\ref{thm:ones} and simplicity of $k$ in \stepref{sr3-occurrence}.~\qed
\end{proof}
\step{sr3-E1k}{The displayed $E_1$ kills the $k$-direction.}
\begin{proof}
\pf\ \raggedright Its numerator there is $k-\theta_2-(k-\theta_2)n/n=0$.~\qed
\end{proof}
\step{sr3-E11}{$E_1$ fixes the $\theta_1$-space.}
\begin{proof}
\pf\ \raggedright Its scalar there is $(\theta_1-\theta_2)/(\theta_1-\theta_2)=1$.~\qed
\end{proof}
\step{sr3-E12}{$E_1$ kills the $\theta_2$-space.}
\begin{proof}
\pf\ \raggedright Its numerator there is $\theta_2-\theta_2=0$.~\qed
\end{proof}
\step{sr3-E2k}{$E_2$ kills the $k$-direction.}
\begin{proof}
\pf\ \raggedright Its numerator there is $k-\theta_1-(k-\theta_1)n/n=0$.~\qed
\end{proof}
\step{sr3-E21}{$E_2$ kills the $\theta_1$-space.}
\begin{proof}
\pf\ \raggedright Its numerator there is $\theta_1-\theta_1=0$.~\qed
\end{proof}
\step{sr3-E22}{$E_2$ fixes the $\theta_2$-space.}
\begin{proof}
\pf\ \raggedright Its scalar there is $(\theta_2-\theta_1)/(\theta_2-\theta_1)=1$.~\qed
\end{proof}
\step{sr3-ranks}{$\rk E_i=\tr E_i=m_i$ for $i=1,2$.}
\begin{proof}
\pf\ \raggedright The actions just checked identify the unique matrices of Theorem~\ref{thm:calculus}, which gives their ranks and traces.~\qed
\end{proof}
\qedstep
\begin{proof}
\pf\ \raggedright Use \stepref{sr3-E0}, \stepref{sr3-E1k}, \stepref{sr3-E11}, \stepref{sr3-E12}, \stepref{sr3-E2k}, \stepref{sr3-E21}, \stepref{sr3-E22}, \stepref{sr3-ranks}.~\qed
\end{proof}
\end{proof}
\qedstep
\begin{proof}
\pf\ \raggedright Use \stepref{sr3-roots}, \stepref{sr3-occurrence}, \stepref{sr3-multiplicities}, \stepref{sr3-idempotents}.~\qed
\end{proof}
\end{proof}

\begin{theorem}[SR2: complements]\label{thm:sr2}
\lean{SimpleGraph.IsSRGWith.spectrumMultiset_compl}
\uses{def:spectrum,facts:srg,thm:ones,thm:sr3}
\mbox{}\par\nobreak
\assume{\begin{pfenum}
\item $X$ is a finite connected simple strongly regular graph with parameters $(n,k,a,c)\in\nn^4$, $0<k<n-1$.
\item $\Delta=(a-c)^2+4(k-c)$, $\theta_1=(a-c+\sqrt\Delta)/2$, $\theta_2=(a-c-\sqrt\Delta)/2$, $m_1=(-k-(n-1)\theta_2)/(\theta_1-\theta_2)$, and $m_2=(k+(n-1)\theta_1)/(\theta_1-\theta_2)$.
\end{pfenum}}
\prove{\raggedright Its complement has parameters
\[
 (n,n-k-1,n-2k+c-2,n-2k+a)
\]
and spectrum
\[
 \ms{(n-k-1)^1,(-1-\theta_1)^{m_1},(-1-\theta_2)^{m_2}},
\]
where the roots and multiplicities are those of
Theorem~\ref{thm:sr3}. Coincident values in this list must be combined;
the complement need not be connected.}
\end{theorem}
\begin{proof}
\pfsketch\ We will count common nonneighbours in each of the two pair cases, then use the constant-direction decomposition.
\step{sr2-parameters}{The complement has the displayed parameters.}
\begin{proof}
\step{sr2-library}{$\comp X$ is strongly regular with Mathlib's natural complement parameters.}
\begin{proof}
\pf\ \raggedright Apply \code{SimpleGraph.IsSRGWith.compl} in Facts~\ref{facts:srg}.~\qed
\end{proof}
\step{sr2-degree}{Its degree is $n-k-1$.}
\begin{proof}
\pf\ \raggedright Among the other $n-1$ vertices, remove the $k$ original neighbours.~\qed
\end{proof}
\step{sr2-nonedge}{For a nonedge $uv$ of $X$, the complement common-neighbour count is $n-2-2k+c$.}
\begin{proof}
\pf\ \raggedright $N_X(u)\cup N_X(v)$ has size $2k-c$ and avoids $u,v$; subtract it from the other $n-2$ vertices.~\qed
\end{proof}
\step{sr2-edge}{For an edge $uv$ of $X$, the complement common-neighbour count is $n-2k+a$.}
\begin{proof}
\pf\ \raggedright The union has size $2k-a$ and contains $u,v$; remove its other $2k-a-2$ vertices from the available $n-2$.~\qed
\end{proof}
\step{sr2-casts}{These parameters equal the displayed integer expressions.}
\begin{proof}
\pf\ \raggedright The actual finite common-neighbour counts in \stepref{sr2-nonedge} and \stepref{sr2-edge} are nonnegative. They justify replacing Mathlib's truncated subtractions by the integer differences; \stepref{sr2-degree} handles the degree.~\qed
\end{proof}
\qedstep
\begin{proof}
\pf\ \raggedright Use \stepref{sr2-library}, \stepref{sr2-casts}.~\qed
\end{proof}
\end{proof}
\step{sr2-spectrum}{The displayed complement multiset is correct, even when values coincide.}
\begin{proof}
\step{sr2-input}{The spectrum of $X$ is $\ms{k^1,\theta_1^{m_1},\theta_2^{m_2}}$.}
\begin{proof}
\pf\ \raggedright Theorem~\ref{thm:sr3}.~\qed
\end{proof}
\step{sr2-map}{Its complement spectrum is $\ms{(n-k-1)^1,(-1-\theta_1)^{m_1},(-1-\theta_2)^{m_2}}$.}
\begin{proof}
\pf\ \raggedright Apply the complement formula of Theorem~\ref{thm:ones} to \stepref{sr2-input}.~\qed
\end{proof}
\step{sr2-collision}{Coincident values contribute the sum of their multiplicities.}
\begin{proof}
\pf\ \raggedright The displayed object in \stepref{sr2-map} is a multiset, by Definition~\ref{def:spectrum}.~\qed
\end{proof}
\qedstep
\begin{proof}
\pf\ \raggedright Use \stepref{sr2-map}, \stepref{sr2-collision}.~\qed
\end{proof}
\end{proof}
\qedstep
\begin{proof}
\pf\ \raggedright Use \stepref{sr2-parameters}, \stepref{sr2-spectrum}.~\qed
\end{proof}
\end{proof}

\begin{corollary}[SR4: least eigenvalue minus two]\label{cor:sr4}
\lean{SimpleGraph.IsSRGWith.parameters_of_leastEigenvalue_eq_neg_two, SimpleGraph.IsSRGWith.rank_two_add_adjMatrix}
\uses{lem:extremes,thm:calculus,thm:sr3,thm:trace}
\mbox{}\par\nobreak
\assume{\begin{pfenum}
\item $X$ is a finite connected simple strongly regular graph with parameters $(n,k,a,c)\in\nn^4$, $0<k<n-1$, and $\lambda_{\min}(X)=-2$.
\item $A=A(X)$, and $\Delta$, $\theta_1$, $\theta_2$ and $m_1$ are those of Theorem~\ref{thm:sr3}.
\end{pfenum}}
\prove{\raggedright \[
 \theta_1=a-c+2\in\ints,\quad k=c+2\theta_1,\quad
 m_1=\frac{2n-k-2}{\theta_1+2},\quad
 \rk(2I+A)=1+m_1.
\]
In particular the positive integer $\theta_1+2$ divides
$2n-k-2$ in $\ints$, and the rank is one plus that quotient.}
\end{corollary}
\begin{proof}
\pfsketch\ We will substitute the least root into its quadratic and count the positive eigenvalues after shifting by two.
\step{sr4-parameters}{$\theta_2=-2$, $\theta_1=a-c+2\in\ints$ and $k=c+2\theta_1$.}
\begin{proof}
\step{sr4-root}{$\theta_2=-2$.}
\begin{proof}
\pf\ \raggedright By Theorem~\ref{thm:sr3}, $\Spec(X)=\ms{k,\theta_1^{m_1},\theta_2^{m_2}}$ with $k>\theta_1>\theta_2$, so Lemma~\ref{lem:extremes} gives $\lambda_{\min}(X)=\theta_2$; use \asm[1]{0}.~\qed
\end{proof}
\step{sr4-sum}{$\theta_1-2=a-c$.}
\begin{proof}
\pf\ \raggedright The definitions give $\theta_1+\theta_2=a-c$; use \stepref{sr4-root}.~\qed
\end{proof}
\step{sr4-product}{$-2\theta_1=c-k$.}
\begin{proof}
\pf\ \raggedright The definitions give $\theta_1\theta_2=((a-c)^2-\Delta)/4=c-k$; use \stepref{sr4-root}.~\qed
\end{proof}
\step{sr4-integer}{$\theta_1$ is the indicated integer.}
\begin{proof}
\pf\ \raggedright Rearrange \stepref{sr4-sum}; $a,c,2$ are integer casts.~\qed
\end{proof}
\qedstep
\begin{proof}
\pf\ \raggedright Use \stepref{sr4-root}, \stepref{sr4-product}, \stepref{sr4-integer}.~\qed
\end{proof}
\end{proof}
\step{sr4-rank}{The multiplicity, rank and divisibility formulas hold.}
\begin{proof}
\step{sr4-multiplicity}{$m_1=(2n-k-2)/(\theta_1+2)$.}
\begin{proof}
\pf\ \raggedright Substitute $\theta_2=-2$ from \stepref{sr4-parameters} in the multiplicity formula of Theorem~\ref{thm:sr3}.~\qed
\end{proof}
\step{sr4-shift}{$2I+A$ has $1+m_1$ positive eigenvalues and all its other eigenvalues are zero.}
\begin{proof}
\pf\ \raggedright Theorem~\ref{thm:calculus} shifts the three values; $k+2>0$ and $\theta_1+2>0$ by Theorem~\ref{thm:sr3}.~\qed
\end{proof}
\step{sr4-rankcount}{$\rk(2I+A)=1+m_1$.}
\begin{proof}
\pf\ \raggedright Apply Theorem~\ref{thm:trace} to \stepref{sr4-shift}.~\qed
\end{proof}
\step{sr4-divides}{The positive integer $\theta_1+2$ divides $2n-k-2$.}
\begin{proof}
\pf\ \raggedright Clear the nonzero denominator in \stepref{sr4-multiplicity}; $m_1$ is a natural number by Theorem~\ref{thm:sr3}.~\qed
\end{proof}
\qedstep
\begin{proof}
\pf\ \raggedright Use \stepref{sr4-multiplicity}, \stepref{sr4-rankcount}, \stepref{sr4-divides}.~\qed
\end{proof}
\end{proof}
\qedstep
\begin{proof}
\pf\ \raggedright Use \stepref{sr4-parameters}, \stepref{sr4-rank}.~\qed
\end{proof}
\end{proof}

\begin{definition}\label{def:design}
\lean{Design, Design.incMatrix, Design.incidenceGraph, Design.flagGraph}
\uses{}
\mbox{}\par\nobreak
\assume{\begin{pfenum}
\item $D$ has finite point and block types, with the parameters and incidence axioms below.
\end{pfenum}}
A $2$-design $D$ consists of finite point and block types and an incidence
relation, with parameters $(v,k,b,r,\lambda)$: $v>1$ points, $b>0$
labelled blocks, $k>1$ points in each block, $r$ blocks at each point,
and $\lambda$ blocks at each pair of distinct points. Repeated blocks
are allowed. In Lean the fields are \code{v k b r lam}.
Its zero-one point-block matrix is $N=D.\mathrm{incMatrix}$.
The graph $I(D)=D.\mathrm{incidenceGraph}$ has the disjoint union of the
point and block types as vertices and exactly the incidences as edges.
The flag graph is $F(D)=D.\mathrm{flagGraph}=L(I(D))$.
\end{definition}

\begin{lemma}[A distinguished constant eigendirection]\label{lem:two-eigenvalues}
\lean{Matrix.IsHermitian.eq_smul_one_add_smul_allOnes_of_two_eigenvalues}
\uses{thm:ones}
\mbox{}\par\nobreak
\assume{\begin{pfenum}
\item $U$ is finite with $|U|=n>0$, $G:U\times U\to\re$ is symmetric, and $G\Zj=b\Zj$.
\item Every eigenvalue of the restriction of $G$ to $\Zj^\perp$ equals $s$; the restriction may have dimension zero.
\end{pfenum}}
\prove{\raggedright \[
 G=sI+\frac{b-s}{n}J.
\]}
\end{lemma}
\begin{proof}
\pfsketch\ We will compare the two matrices on the orthonormal eigenbasis of Theorem~\ref{thm:ones}, which starts with $\Zj/\sqrt n$.
\step{scalarid-onones}{$(sI+\frac{b-s}nJ)\Zj=b\Zj=G\Zj$.}
\begin{proof}
\pf\ \raggedright $J\Zj=n\Zj$ and $s+(b-s)n/n=b$; the last equality is \asm[1]{0}.~\qed
\end{proof}
\step{scalarid-onperp}{$(sI+\frac{b-s}nJ)\Zv=s\Zv=G\Zv$ for each vector $\Zv$ in $\Zj^\perp$ of the eigenbasis of $G$ in Theorem~\ref{thm:ones}.}
\begin{proof}
\pf\ \raggedright $J\Zv=0$ by the coordinate-sum formula of Theorem~\ref{thm:ones}; the eigenvalue of $\Zv$ is $s$ by \asm[2]{0}.~\qed
\end{proof}
\qedstep
\begin{proof}
\pf\ \raggedright By \stepref{scalarid-onones} and \stepref{scalarid-onperp}, the two matrices agree on the basis of Theorem~\ref{thm:ones}, also when $\Zj^\perp=0$; thus they are equal.~\qed
\end{proof}
\end{proof}

\begin{theorem}[Designs from two Gram eigenvalues]\label{thm:two-eigenvalues}
\lean{Matrix.gram_eq_of_two_eigenvalues}
\uses{lem:two-eigenvalues,thm:gram,thm:ones,thm:trace}
\mbox{}\par\nobreak
\assume{\begin{pfenum}
\item For the incidence assertion $N$ is a $v\times b$ zero-one matrix, $v>1$, with positive constant row sum $r$, constant column sum $k>1$, and $\lambda>0$ satisfying $r(k-1)=\lambda(v-1)$.
\item Also $\Spec(NN^T)=\ms{(rk)^1,(r-\lambda)^{v-1}}$; the last bound assumes $r>\lambda$.
\end{pfenum}}
\prove{\raggedright If
$\Spec(NN^T)=\ms{(rk)^1,(r-\lambda)^{v-1}}$, then
\[
 NN^T=(r-\lambda)I+\lambda J.
\]
Thus any two distinct rows have inner product $\lambda$, and $N$
is the incidence matrix of a $2$-design with these parameters, with
columns treated as block labels. If $r>\lambda$, then $b\ge v$.}
\end{theorem}
\begin{proof}
\pfsketch\ We will compare actions on the two orthogonal summands and apply the identity to the incidence Gram matrix.
\step{two-design}{$NN^T=(r-\lambda)I+\lambda J$, so $N$ is the stated design incidence matrix with labelled blocks.}
\begin{proof}
\step{two-constant}{$NN^T\Zj_v=rk\Zj_v$.}
\begin{proof}
\pf\ \raggedright Apply the constant column-sum equation and then the row-sum equation.~\qed
\end{proof}
\step{two-separate}{$rk-(r-\lambda)=\lambda v>0$.}
\begin{proof}
\pf\ \raggedright Use $r(k-1)=\lambda(v-1)$ and collect terms.~\qed
\end{proof}
\step{two-simple}{Every eigenvalue of the restriction of $NN^T$ to $\Zj_v^\perp$ is $r-\lambda$.}
\begin{proof}
\pf\ \raggedright By Theorem~\ref{thm:ones} for the symmetric matrix $NN^T$ and \stepref{two-constant}, $\Spec(NN^T)=\ms{rk,\mu_1,\ldots,\mu_{v-1}}$, where the $\mu_i$ are the eigenvalues on $\Zj_v^\perp$. Compare with \asm[2]{0}: since $rk\ne r-\lambda$ by \stepref{two-separate}, every $\mu_i$ is $r-\lambda$.~\qed
\end{proof}
\step{two-coefficient}{$NN^T=(r-\lambda)I+\lambda J$.}
\begin{proof}
\pf\ \raggedright Apply Lemma~\ref{lem:two-eigenvalues} to $NN^T$ with $b=rk$, $s=r-\lambda$; \stepref{two-simple} gives the hypotheses and \stepref{two-separate} gives the coefficient $\lambda$.~\qed
\end{proof}
\step{two-pairs}{Every pair of distinct rows has inner product $\lambda$.}
\begin{proof}
\pf\ \raggedright Take an off-diagonal entry of \stepref{two-coefficient}; entries count common columns because $N$ is zero-one.~\qed
\end{proof}
\qedstep
\begin{proof}
\pf\ \raggedright Use \stepref{two-coefficient}, \stepref{two-pairs}.~\qed
\end{proof}
\end{proof}
\step{two-fisher}{If $r>\lambda$, then $b\ge v$.}
\begin{proof}
\step{two-fullrank}{$NN^T$ has rank $v$.}
\begin{proof}
\pf\ \raggedright Both displayed eigenvalues are positive, so Theorem~\ref{thm:trace} counts $v$ nonzero entries.~\qed
\end{proof}
\step{two-columnbound}{$v=\rk N\le b$.}
\begin{proof}
\pf\ \raggedright Theorem~\ref{thm:gram} and the dimension of the column domain.~\qed
\end{proof}
\qedstep
\begin{proof}
\pf\ \raggedright Use \stepref{two-columnbound}.~\qed
\end{proof}
\end{proof}
\qedstep
\begin{proof}
\pf\ \raggedright Use \stepref{two-design}, \stepref{two-fisher}.~\qed
\end{proof}
\end{proof}

\begin{theorem}[Spectra of design incidence and flag graphs]\label{thm:design-spectra}
\lean{Design.spectrumMultiset_incidenceGraph, Design.spectrumMultiset_flagGraph}
\uses{cor:scalar-ones,def:design,lem:line-connected,thm:block,thm:calculus,thm:gram,thm:incidence,thm:mn,thm:trace}
\mbox{}\par\nobreak
\assume{\begin{pfenum}
\item $D$ is a design as in Definition~\ref{def:design}, and $N=D.\mathrm{incMatrix}$ is its zero-one incidence matrix with $v>1$ point rows and $b$ labelled block columns; each column has $k>1$ ones, each row $r$ ones, and each distinct row pair has inner product $\lambda>0$.
\item $r>\lambda$, $m=vr=bk$, $I(D)$ is the finite simple bipartite incidence graph, and $F(D)=L(I(D))$.
\end{pfenum}}
\prove{\raggedright Define
\[
 \gamma^{\pm}=\frac{r+k\pm\sqrt{(r-k)^2+4(r-\lambda)}}2.
\]
Then $b\ge v$, $m-v-b+1\ge0$, and
\[
 \begin{aligned}
 \Spec(I(D))&=\ms{\sqrt{rk},-\sqrt{rk},
    (\sqrt{r-\lambda})^{v-1},(-\sqrt{r-\lambda})^{v-1},0^{b-v}},\\
 \Spec(F(D))&=\ms{r+k-2,(\gamma^+-2)^{v-1},
   (\gamma^--2)^{v-1},(k-2)^{b-v},(-2)^{m-v-b+1}}.
 \end{aligned}
\]
Repeated displayed values are combined. Both graphs are connected,
and the flag graph is $(r+k-2)$-regular.}
\end{theorem}
\begin{proof}
\pfsketch\ We will compute the design Gram spectrum, then the adjacency and signless block spectra, and finally the flag spectrum.
\step{design-gram}{$r(k-1)=\lambda(v-1)$, and $NN^T=(r-\lambda)I+\lambda J$ has spectrum $\ms{(rk)^1,(r-\lambda)^{v-1}}$ and rank $v\le b$.}
\begin{proof}
\step{design-flags}{$vr=bk$.}
\begin{proof}
\pf\ \raggedright Count all incident point-block pairs by rows and by columns.~\qed
\end{proof}
\step{design-paircount}{$r(k-1)=\lambda(v-1)$.}
\begin{proof}
\pf\ \raggedright Fix a point; count pairs of a block through it and another point of that block, first by blocks and then by other points.~\qed
\end{proof}
\step{design-entries}{$NN^T=(r-\lambda)I+\lambda J$.}
\begin{proof}
\pf\ \raggedright The diagonal counts $r$ incident blocks; distinct-row entries count $\lambda$ common blocks.~\qed
\end{proof}
\step{design-eigenvalues}{The two Gram eigenvalues are $rk$ once and $r-\lambda$ with multiplicity $v-1$.}
\begin{proof}
\pf\ \raggedright Corollary~\ref{cor:scalar-ones} applied to \stepref{design-entries}, with $r-\lambda+\lambda v=rk$ from \stepref{design-paircount}.~\qed
\end{proof}
\step{design-rank}{$\rk N=v\le b$.}
\begin{proof}
\pf\ \raggedright The eigenvalues in \stepref{design-eigenvalues} are positive; use Theorems~\ref{thm:trace} and \ref{thm:gram}.~\qed
\end{proof}
\qedstep
\begin{proof}
\pf\ \raggedright Use \stepref{design-flags}, \stepref{design-paircount}, \stepref{design-entries}, \stepref{design-eigenvalues}, \stepref{design-rank}.~\qed
\end{proof}
\end{proof}
\step{design-adjacency}{The displayed spectrum of $I(D)$ holds.}
\begin{proof}
\step{design-blockA}{Its adjacency matrix is $\pmat{0&N\\N^T&0}$.}
\begin{proof}
\pf\ \raggedright Incidence joins the two disjoint vertex parts exactly where $N$ has entry one.~\qed
\end{proof}
\step{design-Avalues}{The paired values and $b-v$ zeros are precisely the stated list.}
\begin{proof}
\pf\ \raggedright Apply Theorem~\ref{thm:block} with $s=t=0$ to \stepref{design-gram}, \stepref{design-blockA}.~\qed
\end{proof}
\qedstep
\begin{proof}
\pf\ \raggedright Use \stepref{design-Avalues}.~\qed
\end{proof}
\end{proof}
\step{design-connected}{Both graphs are connected, and $F(D)$ is $(r+k-2)$-regular.}
\begin{proof}
\step{design-pointpaths}{Any two distinct point vertices are joined by a path of length two.}
\begin{proof}
\pf\ \raggedright Their $\lambda>0$ common blocks supply an intermediate vertex.~\qed
\end{proof}
\step{design-blockpaths}{Every block is joined to a point.}
\begin{proof}
\pf\ \raggedright Its size is $k>0$.~\qed
\end{proof}
\step{design-Iconnected}{$I(D)$ is connected and has an edge.}
\begin{proof}
\pf\ \raggedright The paths in \stepref{design-pointpaths}, \stepref{design-blockpaths} join all vertices; $v>0,r>0$ supplies an incidence.~\qed
\end{proof}
\step{design-Fconnected}{The flag graph is connected and has degree $r+k-2$.}
\begin{proof}
\pf\ \raggedright Apply both assertions of Lemma~\ref{lem:line-connected} to \stepref{design-Iconnected}; endpoints of an incidence have degrees $r,k$.~\qed
\end{proof}
\qedstep
\begin{proof}
\pf\ \raggedright Use \stepref{design-Iconnected}, \stepref{design-Fconnected}.~\qed
\end{proof}
\end{proof}
\step{design-flagspectrum}{The displayed spectrum of $F(D)$ holds and $m-v-b+1\ge0$.}
\begin{proof}
\step{design-C}{The signless matrix of $I(D)$ has scalar diagonal blocks $rI_v,kI_b$.}
\begin{proof}
\pf\ \raggedright These are the point and block degrees.~\qed
\end{proof}
\step{design-endpoints}{The Gram value $rk$ gives the pair $0,r+k$; each $r-\lambda$ gives $\gamma^\pm$; the remaining $b-v$ values are $k$.}
\begin{proof}
\pf\ \raggedright Apply Theorem~\ref{thm:block} to \stepref{design-C} and \stepref{design-gram}.~\qed
\end{proof}
\step{design-gammapos}{Both $\gamma^\pm$ are positive.}
\begin{proof}
\pf\ \raggedright They are real with positive sum $r+k$ and positive product $rk-r+\lambda=\lambda v$ by \stepref{design-gram}.~\qed
\end{proof}
\step{design-extra}{$2I+A(F(D))$ has the preceding nonzero list and $m-v-b+1$ zeros.}
\begin{proof}
\pf\ \raggedright Theorems~\ref{thm:incidence} and \ref{thm:mn}, using the single zero in \stepref{design-endpoints}, \stepref{design-gammapos}.~\qed
\end{proof}
\step{design-shift}{The asserted list follows, and its final exponent is nonnegative.}
\begin{proof}
\pf\ \raggedright Subtract two from \stepref{design-extra} using Theorem~\ref{thm:calculus}; a kernel dimension is nonnegative.~\qed
\end{proof}
\qedstep
\begin{proof}
\pf\ \raggedright Use \stepref{design-shift}.~\qed
\end{proof}
\end{proof}
\qedstep
\begin{proof}
\pf\ \raggedright Use \stepref{design-gram}, \stepref{design-adjacency}, \stepref{design-connected}, \stepref{design-flagspectrum}.~\qed
\end{proof}
\end{proof}

\begin{corollary}[The trivial design and complete bipartite graphs]\label{cor:complete-bipartite}
\lean{SimpleGraph.spectrumMultiset_lineGraph_completeBipartite}
\uses{cor:scalar-ones,thm:block,thm:calculus,thm:incidence,thm:mn}
\mbox{}\par\nobreak
\assume{\begin{pfenum}
\item $p,q\in\nn$ satisfy $p,q\ge1$; $K_{p,q}$ is the finite simple complete bipartite graph.
\item The design consequence concerns a zero-one incidence matrix of a $2$-design with $v>1$ points, $b\ge1$ labelled blocks, constant block size $k>1$, and replication $r=\lambda>0$.
\end{pfenum}}
\prove{\raggedright \[
 \Spec(L(K_{p,q}))=
 \ms{p+q-2,(p-2)^{q-1},(q-2)^{p-1},(-2)^{(p-1)(q-1)}}.
\]
Values that coincide are combined. This covers the design case
$r=\lambda>0$, in which every point belongs to every block.}
\end{corollary}
\begin{proof}
\pfsketch\ We will use the one-dimensional image of the all-ones biadjacency matrix.
\step{kb-gram}{The biadjacency matrix has rank one and positive Gram spectrum $\ms{pq}$.}
\begin{proof}
\step{kb-image}{The image of $B=J_{p\times q}$ is $\operatorname{span}(\Zj_p)$ of dimension one.}
\begin{proof}
\pf\ \raggedright Every column equals the nonzero vector $\Zj_p$, and a column exists since $q\ge1$.~\qed
\end{proof}
\step{kb-Gram}{$BB^T=qJ_p$ has the single positive eigenvalue $pq$.}
\begin{proof}
\pf\ \raggedright Apply Corollary~\ref{cor:scalar-ones} with $b=0,t=q$.~\qed
\end{proof}
\qedstep
\begin{proof}
\pf\ \raggedright Use \stepref{kb-image}, \stepref{kb-Gram}.~\qed
\end{proof}
\end{proof}
\step{kb-line}{The displayed line-graph multiset holds.}
\begin{proof}
\step{kb-C}{The signless spectrum is $\ms{0,p+q,q^{p-1},p^{q-1}}$.}
\begin{proof}
\pf\ \raggedright Apply Theorem~\ref{thm:block} with $s=q,t=p$ and \stepref{kb-gram}.~\qed
\end{proof}
\step{kb-connect}{$K_{p,q}$ is connected.}
\begin{proof}
\pf\ \raggedright Opposite-part vertices have an edge; same-part vertices have a two-edge path through the nonempty other part.~\qed
\end{proof}
\step{kb-zeros}{The edge incidence rank is $p+q-1$; its other Gram product has $(p-1)(q-1)$ zeros.}
\begin{proof}
\pf\ \raggedright Use Theorem~\ref{thm:incidence} and the identity $pq-(p+q-1)=(p-1)(q-1)$.~\qed
\end{proof}
\step{kb-subtract}{The required list follows, including $p=1$ or $q=1$.}
\begin{proof}
\pf\ \raggedright Delete the one zero in \stepref{kb-C}, add the zeros in \stepref{kb-zeros} by Theorem~\ref{thm:mn}, and subtract two using Theorem~\ref{thm:calculus}.~\qed
\end{proof}
\qedstep
\begin{proof}
\pf\ \raggedright Use \stepref{kb-subtract}.~\qed
\end{proof}
\end{proof}
\step{kb-trivial}{The design consequence holds.}
\begin{proof}
\step{kb-count}{$r(k-1)=r(v-1)$, so $k=v$.}
\begin{proof}
\pf\ \raggedright Fix a point and count pairs of a block through it and another point in that block: the counts are $r(k-1)$ and $\lambda(v-1)$. Substitute $\lambda=r$ and cancel $r>0$.~\qed
\end{proof}
\step{kb-allpoints}{Each block contains every point.}
\begin{proof}
\pf\ \raggedright A zero-one column on $v$ rows with sum $v$ has every entry one, by \stepref{kb-count}.~\qed
\end{proof}
\qedstep
\begin{proof}
\pf\ \raggedright Use \stepref{kb-allpoints}.~\qed
\end{proof}
\end{proof}
\qedstep
\begin{proof}
\pf\ \raggedright Use \stepref{kb-line}, \stepref{kb-trivial}.~\qed
\end{proof}
\end{proof}

\begin{theorem}[The polynomial on the constant direction]\label{thm:hoffman-poly}
\lean{SimpleGraph.IsRegularOfDegree.aeval_eq_smul_allOnes_of_connected}
\uses{thm:calculus,thm:ones,thm:sp1}
\mbox{}\par\nobreak
\assume{\begin{pfenum}
\item $X$ is a finite connected simple $k$-regular graph on $n>0$ vertices, with adjacency matrix $A$.
\item $f\in\re[t]$.
\end{pfenum}}
\prove{\raggedright If $f$ vanishes at every eigenvalue different from $k$, then
\[
 f(A)=\frac{f(k)}n J.
\]
Conversely this matrix identity implies that $f$ vanishes at every
eigenvalue on $\Zj^\perp$.}
\end{theorem}
\begin{proof}
\pfsketch\ We will evaluate on the constant eigendirection and its orthogonal complement.
\step{hoffman-forward}{\assume{$f$ vanishes at every eigenvalue different from $k$.}\prove{$f(A)=f(k)J/n$.}}
\begin{proof}
\step{hoffman-simple}{The $k$-eigenspace is exactly $\operatorname{span}(\Zj)$.}
\begin{proof}
\pf\ \raggedright Theorem~\ref{thm:sp1} makes it one-dimensional and the constant vector belongs to it.~\qed
\end{proof}
\step{hoffman-constant}{Both matrices act by $f(k)$ on $\Zj$.}
\begin{proof}
\pf\ \raggedright Theorem~\ref{thm:calculus} for $f(A)$ and $J\Zj=n\Zj$ for the other matrix.~\qed
\end{proof}
\step{hoffman-other}{Both matrices kill every eigenvector in $\Zj^\perp$.}
\begin{proof}
\pf\ \raggedright Its eigenvalue differs from $k$ by \stepref{hoffman-simple}; use the vanishing hypothesis and $J\Zx=0$.~\qed
\end{proof}
\qedstep
\begin{proof}
\pf\ \raggedright The two actions agree on the eigenbasis in Theorem~\ref{thm:ones}.~\qed
\end{proof}
\end{proof}
\step{hoffman-converse}{\assume{$f(A)=f(k)J/n$.}\prove{$f$ vanishes at every eigenvalue on $\Zj^\perp$.}}
\begin{proof}
\step{hoffman-apply}{On a nonzero such eigenvector, $f(\mu)\Zv=0$.}
\begin{proof}
\pf\ \raggedright Apply the assumed matrix identity and Theorem~\ref{thm:calculus}; $J\Zv=0$.~\qed
\end{proof}
\step{hoffman-cancel}{$f(\mu)=0$.}
\begin{proof}
\pf\ \raggedright A nonzero vector cannot be killed by a nonzero scalar; use \stepref{hoffman-apply}.~\qed
\end{proof}
\qedstep
\begin{proof}
\pf\ \raggedright Use \stepref{hoffman-cancel}.~\qed
\end{proof}
\end{proof}
\qedstep
\begin{proof}
\pf\ \raggedright Use \stepref{hoffman-forward}, \stepref{hoffman-converse}.~\qed
\end{proof}
\end{proof}

\section{Smith graphs}

\begin{definition}\label{def:smith-models}
\lean{SimpleGraph.threeArmTree, SimpleGraph.smithA, SimpleGraph.smithD, SimpleGraph.smithE}
\uses{}
\mbox{}\par\nobreak
\assume{\begin{pfenum}
\item The parameters satisfy $p,q,r>0$, $\ell\ge2$ for $A$, $\ell\ge4$ for $D$, and $h\in\{6,7,8\}$ for $E$.
\end{pfenum}}
Write $T(p,q,r)$ for the tree consisting of one vertex of degree three
and three internally disjoint arms of positive edge-lengths $p,q,r$.
Define $\widetilde{\mathrm A}_{\ell}$ to be the cycle on $\ell+1$
vertices, $\ell\ge2$; $\widetilde{\mathrm D}_4=K_{1,4}$; and, for
$\ell\ge5$, define $\widetilde{\mathrm D}_{\ell}$ by a path of length
$\ell-4$ with two new leaves at each endpoint. Define
\[
 \widetilde{\mathrm E}_6=T(2,2,2),\quad
 \widetilde{\mathrm E}_7=T(1,3,3),\quad
 \widetilde{\mathrm E}_8=T(1,2,5).
\]
These are the Smith graphs of [CGSS], Figure 1; the finite-type and
affine-type names are kept distinct.
A labelled realization uses a central vertex and arm vertices $(a,j)$,
$1\le j\le h_a$, joined consecutively and to the centre at $j=1$.
For $\widetilde{\mathrm D}_\ell$ label the path consecutively and label
each of its four added leaves separately; cycles use cyclic indices.
Write $\code{SimpleGraph.smithE}(h)$ for $h\in\{6,7,8\}$.
\end{definition}

\begin{theorem}[SM2: positive integer eigenvectors]\label{thm:smith-vectors}
\lean{SimpleGraph.largestEigenvalue_smithA, SimpleGraph.largestEigenvalue_smithD, SimpleGraph.largestEigenvalue_smithE}
\uses{cor:positive-certificate,def:smith-models,thm:calculus,thm:gram,thm:perron,thm:psd,thm:strict,thm:trace}
\mbox{}\par\nobreak
\assume{\begin{pfenum}
\item $X$ is one of the finite simple Smith graphs just defined: $\widetilde{\mathrm A}_\ell$ with $\ell\ge2$, $\widetilde{\mathrm D}_\ell$ with $\ell\ge4$, or $T(2,2,2),T(1,3,3),T(1,2,5)$.
\item $A=A(X)$ and $N=|V(X)|$. The last assertion concerns any finite family in a real inner product space with Gram matrix $2I-A$.
\end{pfenum}}
\prove{\raggedright Every Smith graph has largest eigenvalue $2$, of multiplicity one.
For a cycle, a certificate vector is $\Zj$. For
$\widetilde{\mathrm D}_{\ell}$, it has value $2$ on the central path
(or on the centre when $\ell=4$) and value $1$ on the four leaves.
For $T(p,q,r)$ in the three exceptional cases, take central value
$L=3,4,6$, respectively; at distance $j$ on an arm of length $h$ take
\[
 \Zx(j)=\frac{L(h+1-j)}{h+1}\qquad(1\le j\le h).
\]
For a Smith graph on $N$ vertices, $2I-A$ is positive semidefinite
of rank $N-1$. Deleting any vertex from a Gram representation with
this Gram matrix leaves a basis of its span.}
\end{theorem}
\begin{proof}
\pfsketch\ We will check the neighbour sum in every vertex type, then use positivity and strict monotonicity.
\step{smith-AD}{The proposed coordinates are positive integers and satisfy $A\Zx=2\Zx$ on cycles and affine $D$ graphs.}
\begin{proof}
\step{smith-cycle}{At a cycle vertex the neighbour sum is $1+1=2$.}
\begin{proof}
\pf\ \raggedright Its two neighbours each have coordinate one.~\qed
\end{proof}
\step{smith-Dleaf}{At a $D$ leaf the neighbour sum is $2=2\cdot1$.}
\begin{proof}
\pf\ \raggedright Its unique neighbour belongs to the central path or is the centre of $D_4$.~\qed
\end{proof}
\step{smith-Dinternal}{At an internal path vertex the sum is $2+2=4$.}
\begin{proof}
\pf\ \raggedright The two path neighbours have coordinate two.~\qed
\end{proof}
\step{smith-Dendpoint}{At a nontrivial path endpoint the sum is $2+1+1=4$.}
\begin{proof}
\pf\ \raggedright Its neighbours are one path vertex and the two new leaves.~\qed
\end{proof}
\step{smith-D4}{At the centre of $D_4$ the sum is $1+1+1+1=4$.}
\begin{proof}
\pf\ \raggedright It has precisely its four leaves as neighbours.~\qed
\end{proof}
\qedstep
\begin{proof}
\pf\ \raggedright The coordinates $1$ and $2$ are positive integers. Every vertex of a cycle is covered by \stepref{smith-cycle}; every vertex of an affine $D$ graph is a leaf, an internal path vertex, an endpoint of a nontrivial path, or the centre of $\widetilde{\mathrm D}_4$, covered by \stepref{smith-Dleaf}, \stepref{smith-Dinternal}, \stepref{smith-Dendpoint}, \stepref{smith-D4}.~\qed
\end{proof}
\end{proof}
\step{smith-E}{The proposed exceptional coordinates are positive integers and satisfy $A\Zx=2\Zx$.}
\begin{proof}
\step{smith-E6coords}{For $T(2,2,2)$, the centre has value $3$ and each arm has values $2,1$.}
\begin{proof}
\pf\ \raggedright Substitute $L=3,h=2$ in the stated formula.~\qed
\end{proof}
\step{smith-E7coords}{For $T(1,3,3)$, the centre has value $4$; the short arm has value $2$ and each long arm has values $3,2,1$.}
\begin{proof}
\pf\ \raggedright Substitute $L=4$ and $h=1,3$.~\qed
\end{proof}
\step{smith-E8coords}{For $T(1,2,5)$, the centre has value $6$ and the three arms have values $3$; $4,2$; and $5,4,3,2,1$.}
\begin{proof}
\pf\ \raggedright Substitute $L=6$ and $h=1,2,5$.~\qed
\end{proof}
\step{smith-arm}{At an internal arm vertex, the two neighbour coordinates sum to twice its coordinate.}
\begin{proof}
\pf\ \raggedright The values $L(h+1-j)/(h+1)$ form an arithmetic progression, including central value at $j=0$.~\qed
\end{proof}
\step{smith-leaf}{At an arm leaf the single neighbour coordinate is twice the leaf coordinate.}
\begin{proof}
\pf\ \raggedright At $j=h$, the values at $h-1,h$ are $2L/(h+1),L/(h+1)$.~\qed
\end{proof}
\step{smith-centre6}{The $E_6$ central neighbour sum is $2+2+2=6=2\cdot3$.}
\begin{proof}
\pf\ \raggedright Use \stepref{smith-E6coords}.~\qed
\end{proof}
\step{smith-centre7}{The $E_7$ central neighbour sum is $2+3+3=8=2\cdot4$.}
\begin{proof}
\pf\ \raggedright Use \stepref{smith-E7coords}.~\qed
\end{proof}
\step{smith-centre8}{The $E_8$ central neighbour sum is $3+4+5=12=2\cdot6$.}
\begin{proof}
\pf\ \raggedright Use \stepref{smith-E8coords}.~\qed
\end{proof}
\qedstep
\begin{proof}
\pf\ \raggedright The values in \stepref{smith-E6coords}, \stepref{smith-E7coords}, \stepref{smith-E8coords} are positive integers. Every vertex is the centre, covered by \stepref{smith-centre6}, \stepref{smith-centre7}, \stepref{smith-centre8}, an arm vertex other than a leaf, covered by \stepref{smith-arm}, or an arm leaf, covered by \stepref{smith-leaf}.~\qed
\end{proof}
\end{proof}
\step{smith-spectral}{$\lambda_{\max}(X)=2$, its multiplicity is one, and $2I-A\succeq0$ has rank $N-1$.}
\begin{proof}
\step{smith-positivecert}{The largest eigenvalue is two.}
\begin{proof}
\pf\ \raggedright Apply Corollary~\ref{cor:positive-certificate} to \stepref{smith-AD} or \stepref{smith-E}, according to the model.~\qed
\end{proof}
\step{smith-connected}{$X$ is connected.}
\begin{proof}
\step{smith-Aconnected}{\case{$X$ is a cycle.}}
\begin{proof}
\pf\ \raggedright Following cyclic indices reaches every vertex from any chosen vertex.~\qed
\end{proof}
\step{smith-Dconnected}{\case{$X$ is an affine $D$ graph.}}
\begin{proof}
\pf\ \raggedright Each leaf is attached to the central path or centre, and that path is connected.~\qed
\end{proof}
\step{smith-Econnected}{\case{$X$ is one of the three exceptional trees.}}
\begin{proof}
\pf\ \raggedright Every arm vertex has a path along its arm to the common centre.~\qed
\end{proof}
\qedstep
\begin{proof}
\pf\ \raggedright By \asm[1]{0} the three cases are exhaustive.~\qed
\end{proof}
\end{proof}
\step{smith-simple}{The eigenvalue two is simple.}
\begin{proof}
\pf\ \raggedright Theorem~\ref{thm:perron}, using \stepref{smith-connected}.~\qed
\end{proof}
\step{smith-psdrank}{The shifted matrix is positive semidefinite with exactly one zero eigenvalue.}
\begin{proof}
\pf\ \raggedright Theorems~\ref{thm:psd}, \ref{thm:calculus}, \ref{thm:trace} and \stepref{smith-positivecert}, \stepref{smith-simple}.~\qed
\end{proof}
\qedstep
\begin{proof}
\pf\ \raggedright Use \stepref{smith-positivecert}, \stepref{smith-simple}, \stepref{smith-psdrank}.~\qed
\end{proof}
\end{proof}
\step{smith-delete}{Deleting any representing vector leaves a basis of the original span.}
\begin{proof}
\step{smith-proper}{Deleting one vertex gives a nonempty proper subgraph of $X$.}
\begin{proof}
\pf\ \raggedright Every defined Smith graph has at least three vertices.~\qed
\end{proof}
\step{smith-strict}{That graph has largest eigenvalue less than two.}
\begin{proof}
\pf\ \raggedright Theorem~\ref{thm:strict} and \stepref{smith-spectral}.~\qed
\end{proof}
\step{smith-minor}{The remaining Gram matrix is positive definite of order $N-1$.}
\begin{proof}
\pf\ \raggedright Its eigenvalues are $2-\mu>0$ by \stepref{smith-strict} and Theorem~\ref{thm:psd}.~\qed
\end{proof}
\step{smith-basis}{The $N-1$ remaining vectors are independent and span the original space.}
\begin{proof}
\pf\ \raggedright Theorem~\ref{thm:gram} gives their span dimension $N-1$ from \stepref{smith-minor}; the original span also has dimension $N-1$ by \stepref{smith-spectral}.~\qed
\end{proof}
\qedstep
\begin{proof}
\pf\ \raggedright Use \stepref{smith-basis}.~\qed
\end{proof}
\end{proof}
\qedstep
\begin{proof}
\pf\ \raggedright Use \stepref{smith-AD}, \stepref{smith-E}, \stepref{smith-spectral}, \stepref{smith-delete}.~\qed
\end{proof}
\end{proof}

\begin{lemma}[Trees of maximum degree two]\label{lem:tree-path}
\lean{SimpleGraph.IsTree.isPath_of_degree_le_two}
\uses{}
\mbox{}\par\nobreak
\assume{\begin{pfenum}
\item $X$ is a finite tree and each vertex has degree at most two.
\end{pfenum}}
\prove{\raggedright A finite tree whose vertices have degree at most two is a nonempty path.}
\end{lemma}
\begin{proof}
\pfsketch\ We will take a longest simple path and show that no vertex lies outside it.
\step{tree-longest}{\pick\ a simple path $P$ of maximum length.}
\begin{proof}
\pf\ \raggedright Path supports are lists of distinct vertices of a finite set, so their lengths have a finite maximum; a tree is nonempty, so a singleton path exists.~\qed
\end{proof}
\step{tree-attach}{A vertex outside $P$ would have a path to $P$ meeting $P$ only at its last vertex.}
\begin{proof}
\pf\ \raggedright Use connectedness, and stop at the first vertex of $P$ on a path from the outside vertex.~\qed
\end{proof}
\step{tree-attachend}{Attachment at an endpoint of $P$ contradicts maximal length.}
\begin{proof}
\pf\ \raggedright Concatenate the nonempty outside segment with $P$; by \stepref{tree-attach} no vertex repeats.~\qed
\end{proof}
\step{tree-attachinside}{Attachment at an internal vertex of $P$ contradicts degree at most two.}
\begin{proof}
\pf\ \raggedright The attachment edge and the two path edges at that vertex have distinct other endpoints.~\qed
\end{proof}
\step{tree-edges}{$X$ has no edge outside $P$.}
\begin{proof}
\pf\ \raggedright Both endpoints lie on $P$ by \stepref{tree-attach}, \stepref{tree-attachend}, \stepref{tree-attachinside}; an edge between nonconsecutive vertices of $P$ would close a cycle.~\qed
\end{proof}
\qedstep
\begin{proof}
\pf\ \raggedright By \stepref{tree-attachend} and \stepref{tree-attachinside}, an attachment vertex is an endpoint or an internal vertex, and both are impossible; thus every vertex lies on $P$, and by \stepref{tree-edges} $X$ is the path $P$.~\qed
\end{proof}
\end{proof}

\begin{theorem}[SM3: Smith's theorem]\label{thm:smith}
\lean{SimpleGraph.smith, SimpleGraph.largestEigenvalue_lt_two_iff}
\uses{def:smith-models,facts:tree,lem:extremes,lem:tree-path,prop:cospectral,thm:smith-vectors,thm:strict}
\mbox{}\par\nobreak
\assume{\begin{pfenum}
\item $X$ is a finite connected simple graph; in particular its vertex set is nonempty.
\item The Smith graphs and $T(p,q,r)$ have the explicit meanings above; the paths in the conclusion are nonempty.
\end{pfenum}}
\prove{\raggedright A connected finite graph has largest eigenvalue $2$ iff it is
isomorphic to a Smith graph. It has largest eigenvalue less than
$2$ iff it is a path, a tree $T(1,1,r)$ for $r\ge1$, or one of
$T(1,2,2),T(1,2,3),T(1,2,4)$.}
\end{theorem}
\begin{proof}
\pfsketch\ We will first classify graphs that contain no Smith subgraph and then use strict monotonicity to distinguish equality from a strict bound.
\step{smith-containing}{\assume{$X$ contains a Smith graph $Z$ as a subgraph, through an edge-preserving vertex injection.}\prove{$\lambda_{\max}(X)=2$ if the injection is onto vertices and edges, and $\lambda_{\max}(X)>2$ otherwise.}}
\begin{proof}
\step{smith-equalmodel}{\case{The injection is onto vertices and edges.}}
\begin{proof}
\pf\ \raggedright It is then an isomorphism $Z\cong X$. By Proposition~\ref{prop:cospectral}, $X$ and $Z$ have the same spectrum, and thus the same largest eigenvalue by Lemma~\ref{lem:extremes}; Theorem~\ref{thm:smith-vectors} gives $\lambda_{\max}(Z)=2$.~\qed
\end{proof}
\step{smith-propermodel}{\case{The injection omits a vertex or an edge.}}
\begin{proof}
\pf\ \raggedright $Z$ is a nonempty proper subgraph of the connected graph $X$; Theorem~\ref{thm:strict} and Theorem~\ref{thm:smith-vectors} give $2=\lambda_{\max}(Z)<\lambda_{\max}(X)$.~\qed
\end{proof}
\qedstep
\begin{proof}
\pf\ \raggedright The two cases are exhaustive.~\qed
\end{proof}
\end{proof}
\step{smith-tree}{\assume{$X$ contains no Smith subgraph.}\prove{$X$ is a tree of maximum degree at most three, with at most one degree-three vertex.}}
\begin{proof}
\step{smith-nocycle}{$X$ has no cycle.}
\begin{proof}
\pf\ \raggedright A simple cycle has length at least three and is a Smith graph $\widetilde{\mathrm A}_{\ell}$.~\qed
\end{proof}
\step{smith-unique-paths}{There is a unique simple path between every two vertices.}
\begin{proof}
\pf\ \raggedright Connectedness and acyclicity in \stepref{smith-nocycle}; apply Facts~\ref{facts:tree}.~\qed
\end{proof}
\step{smith-degree4}{No vertex has four distinct neighbours.}
\begin{proof}
\pf\ \raggedright That vertex and four incident edges would give the Smith subgraph $K_{1,4}$.~\qed
\end{proof}
\step{smith-two-branches}{Two degree-three vertices would give a Smith subgraph.}
\begin{proof}
\step{smith-minimalpair}{\pick\ two degree-three vertices at minimum distance; internal vertices of their path have degree two.}
\begin{proof}
\pf\ \raggedright An internal degree-three vertex would form a shorter pair; internal path vertices already have two neighbours.~\qed
\end{proof}
\step{smith-offpath}{The two off-path neighbours at each endpoint are outside the path and distinct from all three other choices.}
\begin{proof}
\pf\ \raggedright A repeated vertex or return to the path creates two distinct simple paths, contradicting \stepref{smith-unique-paths}.~\qed
\end{proof}
\step{smith-Dsubgraph}{The connecting path and four extra edges form $\widetilde{\mathrm D}_{d+4}$, where $d\ge1$ is the path length.}
\begin{proof}
\pf\ \raggedright The explicit definition in \ref{def:smith-models} has that central path and four leaves; \stepref{smith-offpath} makes the vertex map injective.~\qed
\end{proof}
\qedstep
\begin{proof}
\pf\ \raggedright The subgraph in \stepref{smith-Dsubgraph} contradicts the assumption of \stepref{smith-tree}.~\qed
\end{proof}
\end{proof}
\qedstep
\begin{proof}
\pf\ \raggedright Use \stepref{smith-unique-paths}, \stepref{smith-degree4}, \stepref{smith-two-branches}.~\qed
\end{proof}
\end{proof}
\step{smith-shapes}{\assume{$X$ contains no Smith subgraph.}\prove{$X$ is a path or $T(p,q,r)$ with $1\le p\le q\le r$.}}
\begin{proof}
\step{smith-three-arms}{A tree with exactly one degree-three vertex consists of three nonempty path arms.}
\begin{proof}
\step{smith-deletecentre}{Deleting the central vertex gives exactly three components, one at each neighbour.}
\begin{proof}
\pf\ \raggedright Unique paths in \stepref{smith-tree} put distinct neighbours in distinct components, and every remaining vertex reaches the centre through one of them.~\qed
\end{proof}
\step{smith-armdegree}{Each component has maximum degree at most two, and its attachment vertex has degree at most one within the component.}
\begin{proof}
\pf\ \raggedright No other vertex has degree three by \stepref{smith-tree}; the attachment loses its edge to the centre.~\qed
\end{proof}
\step{smith-armspaths}{Each component is a path attached at an endpoint.}
\begin{proof}
\pf\ \raggedright Each component is connected and acyclic as a subtree. Apply Lemma~\ref{lem:tree-path}, using \stepref{smith-armdegree}; its attachment cannot be internal because its component degree is at most one.~\qed
\end{proof}
\qedstep
\begin{proof}
\pf\ \raggedright The three components give positive arm lengths; reorder them to obtain $p\le q\le r$.~\qed
\end{proof}
\end{proof}
\step{smith-shape-pathcase}{\case{$X$ has no vertex of degree three.}}
\begin{proof}
\pf\ \raggedright By \stepref{smith-tree}, $X$ is a tree of maximum degree at most two; apply Lemma~\ref{lem:tree-path}.~\qed
\end{proof}
\step{smith-shape-armcase}{\case{$X$ has a vertex of degree three.}}
\begin{proof}
\pf\ \raggedright It is the only one, by \stepref{smith-tree}; apply \stepref{smith-three-arms}.~\qed
\end{proof}
\qedstep
\begin{proof}
\pf\ \raggedright The two cases are exhaustive, and by \stepref{smith-tree} no vertex has degree greater than three.~\qed
\end{proof}
\end{proof}
\step{smith-armbounds}{\assume{$X=T(p,q,r)$, $1\le p\le q\le r$, and $X$ contains no Smith subgraph.}\prove{The only possibilities are $(1,1,r)$ for $r\ge1$ and $(1,2,r)$ for $2\le r\le4$.}}
\begin{proof}
\step{smith-p}{$p=1$.}
\begin{proof}
\pf\ \raggedright If $p\ge2$, truncate each arm to length two to obtain $T(2,2,2)$.~\qed
\end{proof}
\step{smith-q}{$q=1$ or $q=2$.}
\begin{proof}
\pf\ \raggedright If $q\ge3$, \stepref{smith-p} and $r\ge q$ allow truncation to $T(1,3,3)$.~\qed
\end{proof}
\step{smith-r}{If $q=2$, then $2\le r\le4$.}
\begin{proof}
\pf\ \raggedright Ordering gives $r\ge2$; $r\ge5$ permits truncation to $T(1,2,5)$.~\qed
\end{proof}
\qedstep
\begin{proof}
\pf\ \raggedright Use \stepref{smith-p}, \stepref{smith-q}, \stepref{smith-r}.~\qed
\end{proof}
\end{proof}
\step{smith-embeddings}{Every graph on the finite-type list is a proper subgraph of a Smith graph, and has largest eigenvalue less than two.}
\begin{proof}
\step{smith-path-embedding}{A path on $N\ge1$ vertices embeds properly in a cycle on $N+2$ vertices.}
\begin{proof}
\pf\ \raggedright Use $N$ consecutive cycle vertices with their $N-1$ path edges; at least two vertices are omitted.~\qed
\end{proof}
\step{smith-Dembedding}{$T(1,1,r)$ embeds properly in $\widetilde{\mathrm D}_{r+4}$.}
\begin{proof}
\pf\ \raggedright Adjoin two new leaves at the far endpoint of its length-$r$ arm; that arm becomes the central path.~\qed
\end{proof}
\step{smith-E6embedding}{$T(1,2,2)$ embeds properly in $T(1,2,5)$.}
\begin{proof}
\pf\ \raggedright Keep the first two vertices of the length-five arm and omit the last three.~\qed
\end{proof}
\step{smith-E7embedding}{$T(1,2,3)$ embeds properly in $T(1,2,5)$.}
\begin{proof}
\pf\ \raggedright Keep the first three vertices of the length-five arm and omit the last two.~\qed
\end{proof}
\step{smith-E8embedding}{$T(1,2,4)$ embeds properly in $T(1,2,5)$.}
\begin{proof}
\pf\ \raggedright Keep the first four vertices of the length-five arm and omit the last one.~\qed
\end{proof}
\step{smith-lt}{Each listed graph has largest eigenvalue less than two.}
\begin{proof}
\pf\ \raggedright Apply Theorems~\ref{thm:strict}, \ref{thm:smith-vectors} to the five explicit embeddings.~\qed
\end{proof}
\qedstep
\begin{proof}
\pf\ \raggedright Use \stepref{smith-path-embedding}, \stepref{smith-Dembedding}, \stepref{smith-E6embedding}, \stepref{smith-E7embedding}, \stepref{smith-E8embedding}, \stepref{smith-lt}.~\qed
\end{proof}
\end{proof}
\step{smith-eq-forward}{\assume{$\lambda_{\max}(X)=2$.}\prove{$X$ is isomorphic to a Smith graph.}}
\begin{proof}
\step{smith-eq-containing}{\case{$X$ contains a Smith subgraph.}}
\begin{proof}
\pf\ \raggedright By \stepref{smith-containing} and \asm{1}, the injection is onto vertices and edges, so it is an isomorphism.~\qed
\end{proof}
\step{smith-eq-avoiding}{\case{$X$ contains no Smith subgraph.}}
\begin{proof}
\pf\ \raggedright By \stepref{smith-shapes} and \stepref{smith-armbounds}, $X$ is on the finite-type list, so $\lambda_{\max}(X)<2$ by \stepref{smith-embeddings}; this contradicts \asm{1}.~\qed
\end{proof}
\qedstep
\begin{proof}
\pf\ \raggedright The two cases are exhaustive.~\qed
\end{proof}
\end{proof}
\step{smith-eq-reverse}{\assume{$X$ is isomorphic to a Smith graph.}\prove{$\lambda_{\max}(X)=2$.}}
\begin{proof}
\pf\ \raggedright Use Theorem~\ref{thm:smith-vectors}, Proposition~\ref{prop:cospectral} and Lemma~\ref{lem:extremes}.~\qed
\end{proof}
\step{smith-lt-forward}{\assume{$\lambda_{\max}(X)<2$.}\prove{$X$ is on the finite-type list.}}
\begin{proof}
\step{smith-lt-containing}{$X$ contains no Smith subgraph.}
\begin{proof}
\pf\ \raggedright Both cases of \stepref{smith-containing} give $\lambda_{\max}(X)\ge2$, contrary to \asm{1}.~\qed
\end{proof}
\step{smith-lt-avoiding}{$X$ lies on the stated list.}
\begin{proof}
\pf\ \raggedright Apply \stepref{smith-shapes} and \stepref{smith-armbounds} with \stepref{smith-lt-containing}.~\qed
\end{proof}
\qedstep
\begin{proof}
\pf\ \raggedright Use \stepref{smith-lt-avoiding}.~\qed
\end{proof}
\end{proof}
\step{smith-lt-reverse}{\assume{$X$ is on the finite-type list.}\prove{$\lambda_{\max}(X)<2$.}}
\begin{proof}
\pf\ \raggedright This is \stepref{smith-embeddings}.~\qed
\end{proof}
\qedstep
\begin{proof}
\pf\ \raggedright Use \stepref{smith-eq-forward}, \stepref{smith-eq-reverse}, \stepref{smith-lt-forward} and \stepref{smith-lt-reverse}.~\qed
\end{proof}
\end{proof}

\section{Exact certificates and rational transfer}

\begin{lemma}[Moments on an annihilating support]\label{lem:annihilating-moments}
\lean{Matrix.IsHermitian.eigenvalues_mem_of_annihilating, Matrix.IsHermitian.trace_pow_eq_multiplicities}
\uses{thm:calculus,thm:trace}
\mbox{}\par\nobreak
\assume{\begin{pfenum}
\item $A$ is real symmetric of finite order; $(\theta_i)_{i<s}$ is injective, $s\ge0$, and $\prod_{i<s}(A-\theta_iI)=0$.
\end{pfenum}}
\prove{\raggedright every eigenvalue belongs to the nodes and
\[
 \tr A^j=\sum_{i<s}\mult_A(\theta_i)\theta_i^j\qquad(j\ge0).
\]}
\end{lemma}
\begin{proof}
\pfsketch\ We will evaluate the annihilator on eigenvectors, and then group the trace-power sum of Theorem~\ref{thm:trace} by nodes.
\step{ann-evaluate}{Every eigenvalue $\mu$ satisfies $\prod_{i<s}(\mu-\theta_i)=0$.}
\begin{proof}
\pf\ \raggedright Apply Theorem~\ref{thm:calculus} to the given annihilator.~\qed
\end{proof}
\step{ann-support}{Each eigenvalue equals a node.}
\begin{proof}
\pf\ \raggedright A product of real numbers vanishes only if a factor vanishes; use \stepref{ann-evaluate}.~\qed
\end{proof}
\step{ann-unique}{Each eigenvalue occurrence has exactly one node index.}
\begin{proof}
\pf\ \raggedright Existence is \stepref{ann-support}; injectivity of the node function gives uniqueness.~\qed
\end{proof}
\step{ann-group}{$\tr A^j=\sum_{i<s}\mult_A(\theta_i)\theta_i^j$ for every $j\ge0$.}
\begin{proof}
\pf\ \raggedright Group the trace-power sum of Theorem~\ref{thm:trace} by the unique index in \stepref{ann-unique}.~\qed
\end{proof}
\step{ann-empty}{If $s=0$, the matrix has order zero.}
\begin{proof}
\pf\ \raggedright The annihilator is $I=0$; any diagonal index would imply $1=0$. Thus the spectrum and all traces are empty sums.~\qed
\end{proof}
\qedstep
\begin{proof}
\pf\ \raggedright Use \stepref{ann-support}, \stepref{ann-group}, \stepref{ann-empty}.~\qed
\end{proof}
\end{proof}

\begin{theorem}[Moments determine multiplicities on a known support]\label{thm:certificate}
\lean{Matrix.IsHermitian.spectrumMultiset_eq_of_annihilatingPolynomial_trace}
\uses{facts:roots,facts:vandermonde,lem:annihilating-moments,thm:multiplicity}
\mbox{}\par\nobreak
\assume{\begin{pfenum}
\item $A$ is real symmetric of order $n$; the nodes $(\theta_i)_{i<s}$ are distinct with $s\ge0$, and $m_i\in\nn$.
\item $\prod_{i<s}(A-\theta_i I)=0$ and $\tr A^j=\sum_{i<s}m_i\theta_i^j$ for $0\le j<s$.
\end{pfenum}}
\prove{\raggedright $\Spec(A)=\ms{\theta_0^{m_0},\ldots,\theta_{s-1}^{m_{s-1}}}$.
In particular $\sum_i m_i=n$, including the empty node list.}
\end{theorem}
\begin{proof}
\pfsketch\ We will compare the assumed moments with the true moments of Lemma~\ref{lem:annihilating-moments}, and conclude by Vandermonde uniqueness as in \code{LineGraphs/Spectral/Basic.lean}.
\step{cert-support}{The spectrum is supported on the nodes, and its moments have coefficients $d_i=\mult_A(\theta_i)$.}
\begin{proof}
\pf\ \raggedright Apply Lemma~\ref{lem:annihilating-moments} to $A$.~\qed
\end{proof}
\step{cert-difference}{$\sum_{i<s}(d_i-m_i)\theta_i^j=0$ for every $j<s$.}
\begin{proof}
\pf\ \raggedright Subtract the assumed moments from those in \stepref{cert-support}.~\qed
\end{proof}
\step{cert-vandermonde}{The real vector $(d_i-m_i)_{i<s}$ is zero.}
\begin{proof}
\pf\ \raggedright Apply \code{Matrix.eq_zero_of_forall_pow_sum_mul_pow_eq_zero} in Facts~\ref{facts:vandermonde} to the injective nodes and \stepref{cert-difference}; its finite index size may be zero.~\qed
\end{proof}
\step{cert-counts}{$d_i=m_i$ as natural numbers for every $i<s$.}
\begin{proof}
\pf\ \raggedright Evaluate \stepref{cert-vandermonde} at $i$ and use injectivity of the natural-number cast into $\re$.~\qed
\end{proof}
\step{cert-list}{$\Spec(A)=\ms{\theta_i^{m_i}:i<s}$.}
\begin{proof}
\pf\ \raggedright Use \stepref{cert-support}, \stepref{cert-counts} and multiset extensionality in Facts~\ref{facts:roots}; outside the nodes both counts are zero.~\qed
\end{proof}
\step{cert-order}{$\sum_i m_i=n$, even when $s=0$.}
\begin{proof}
\pf\ \raggedright Take cardinalities in \stepref{cert-list} and apply Theorem~\ref{thm:multiplicity}.~\qed
\end{proof}
\qedstep
\begin{proof}
\pf\ \raggedright Use \stepref{cert-list}, \stepref{cert-order}.~\qed
\end{proof}
\end{proof}

\begin{corollary}[Two matrices on a common annihilating support]\label{cor:certificate-pair}
\lean{Matrix.IsHermitian.charpoly_eq_of_annihilatingPolynomial_trace}
\uses{lem:annihilating-moments,prop:cospectral,thm:certificate}
\mbox{}\par\nobreak
\assume{\begin{pfenum}
\item $A,B$ are real symmetric on finite index sets, annihilate the product on the same distinct nodes $(\theta_i)_{i<s}$ with $s\ge0$, and have equal traces for $j<s$.
\end{pfenum}}
\prove{\raggedright $\Spec(A)=\Spec(B)$ and $\chi_A=\chi_B$.}
\end{corollary}
\begin{proof}
\pfsketch\ We will read the moments of $B$ from Lemma~\ref{lem:annihilating-moments} and certify $A$ with them by Theorem~\ref{thm:certificate}.
\step{pair-Bmoments}{$\tr B^j=\sum_i\mult_B(\theta_i)\theta_i^j$.}
\begin{proof}
\pf\ \raggedright Apply Lemma~\ref{lem:annihilating-moments} to $B$.~\qed
\end{proof}
\step{pair-Amoments}{$A$ has these moments for $j<s$.}
\begin{proof}
\pf\ \raggedright Compose the assumed trace equalities with \stepref{pair-Bmoments}.~\qed
\end{proof}
\step{pair-Alist}{$\Spec(A)=\ms{\theta_i^{\mult_B(\theta_i)}:i<s}$.}
\begin{proof}
\pf\ \raggedright Apply Theorem~\ref{thm:certificate} to $A$, choosing the natural multiplicities of $B$ and using \stepref{pair-Amoments}.~\qed
\end{proof}
\step{pair-Blist}{$\Spec(B)$ is the same multiset.}
\begin{proof}
\pf\ \raggedright The support assertion of Lemma~\ref{lem:annihilating-moments} for $B$ identifies all counts, both on and off the nodes.~\qed
\end{proof}
\step{pair-charpoly}{$\chi_A=\chi_B$.}
\begin{proof}
\pf\ \raggedright Apply Proposition~\ref{prop:cospectral} to \stepref{pair-Alist} and \stepref{pair-Blist}.~\qed
\end{proof}
\qedstep
\begin{proof}
\pf\ \raggedright Use \stepref{pair-Alist}, \stepref{pair-Blist}, \stepref{pair-charpoly}.~\qed
\end{proof}
\end{proof}

\begin{theorem}[Similarity and intertwining certificates]\label{thm:similarity}
\lean{Matrix.charpoly_eq_of_invertible_intertwiner, Matrix.charpoly_eq_of_scaled_inverse_intertwiner, SimpleGraph.cospectral_of_intertwiner}
\uses{def:graph,facts:charpoly,facts:det}
\mbox{}\par\nobreak
\assume{\begin{pfenum}
\item $U$ is finite; $A,B,Q:U\times U\to F$ are matrices over a field $F$, $\det Q\ne0$, and $QA=BQ$.
\item The involution variant assumes $Q^2=I$ and $QAQ=B$ over that field.
\item The integer variant assumes $A,B,Q,R$ are integer matrices of the same finite order, $d\in\ints\setminus\{0\}$, $QA=BQ$ and $QR=RQ=dI$.
\end{pfenum}}
\prove{\raggedright $\chi_A=\chi_B$.
In particular $Q^2=I$ and $QAQ=B$ suffice.
For integer matrices $A,B,Q,R$ and an integer $d\ne0$, it is enough
to check
\[
 QA=BQ,\qquad QR=RQ=dI.
\]
These identities imply equality of the real characteristic polynomials
of $A$ and $B$. For adjacency matrices, they prove graph cospectrality.}
\end{theorem}
\begin{proof}
\pfsketch\ We will conjugate by an invertible matrix and then construct the rational inverse explicitly for the integer certificate.
\step{similarity-field}{$\chi_A=\chi_B$ over $F$.}
\begin{proof}
\step{similarity-inverse}{$Q$ has an inverse and $B=QAQ^{-1}$.}
\begin{proof}
\pf\ \raggedright The determinant-unit criterion in Facts~\ref{facts:det} and right multiplication of $QA=BQ$ by $Q^{-1}$.~\qed
\end{proof}
\step{similarity-det}{$\det(tI-B)=\det(tI-A)$.}
\begin{proof}
\pf\ \raggedright Apply \code{Matrix.charpoly_units_conj} in Facts~\ref{facts:charpoly} to the unit with value $Q$ and inverse $Q^{-1}$ from \stepref{similarity-inverse}.~\qed
\end{proof}
\qedstep
\begin{proof}
\pf\ \raggedright Use \stepref{similarity-det}.~\qed
\end{proof}
\end{proof}
\step{similarity-involution}{$Q^2=I$ and $QAQ=B$ suffice.}
\begin{proof}
\step{similarity-Qinverse}{$Q^{-1}=Q$.}
\begin{proof}
\pf\ \raggedright The equation $Q^2=I$ gives a two-sided inverse.~\qed
\end{proof}
\step{similarity-Qconj}{$\chi_A=\chi_B$.}
\begin{proof}
\pf\ \raggedright Apply \code{Matrix.charpoly_units_conj} in Facts~\ref{facts:charpoly} directly to the unit with value and inverse $Q$, using $QAQ=B$.~\qed
\end{proof}
\qedstep
\begin{proof}
\pf\ \raggedright Use \stepref{similarity-Qinverse}, \stepref{similarity-Qconj}.~\qed
\end{proof}
\end{proof}
\step{similarity-integer}{The integer certificate implies equality of real characteristic polynomials.}
\begin{proof}
\step{similarity-castQ}{Over $\rats$, $Q((1/d)R)=((1/d)R)Q=I$.}
\begin{proof}
\pf\ \raggedright Cast the two identities into $\rats$, distribute $1/d$, and cancel $d\ne0$.~\qed
\end{proof}
\step{similarity-rational}{The rational characteristic polynomials are equal.}
\begin{proof}
\pf\ \raggedright Apply \code{Matrix.charpoly_units_conj} in Facts~\ref{facts:charpoly} over $\rats$ to the unit from \stepref{similarity-castQ}; $QA=BQ$ gives its conjugation equation.~\qed
\end{proof}
\step{similarity-real}{The equality persists over $\re$.}
\begin{proof}
\pf\ \raggedright Apply \code{Matrix.charpoly_map} in Facts~\ref{facts:charpoly} to the inclusion $\rats\hookrightarrow\re$.~\qed
\end{proof}
\step{similarity-graphs}{For adjacency matrices this is graph cospectrality.}
\begin{proof}
\pf\ \raggedright Use Definition~\ref{def:graph}.~\qed
\end{proof}
\qedstep
\begin{proof}
\pf\ \raggedright Use \stepref{similarity-real}, \stepref{similarity-graphs}.~\qed
\end{proof}
\end{proof}
\qedstep
\begin{proof}
\pf\ \raggedright Use \stepref{similarity-field}, \stepref{similarity-involution}, \stepref{similarity-integer}.~\qed
\end{proof}
\end{proof}

\begin{theorem}[Seidel switching]\label{thm:seidel}
\lean{SimpleGraph.seidelMatrix_switch, SimpleGraph.spectrumMultiset_seidelMatrix_switch}
\uses{def:graph,prop:cospectral,thm:calculus,thm:similarity}
\mbox{}\par\nobreak
\assume{\begin{pfenum}
\item $X$ is a finite simple graph on a finite vertex set $V$, and $U\subseteq V$.
\item $X^U$ complements adjacency across $U,V\setminus U$ and retains adjacency within each part; $T$ is diagonal with value $-1$ on $U$ and $1$ outside.
\end{pfenum}}
\prove{\raggedright \[
 S(X^U)=T S(X)T,\qquad T^2=I.
\]
Thus Seidel switching preserves the Seidel spectrum with multiplicity.
The matrix used in [CGSS] Corollary 5.10 is $-S=2A-J+I$, and its
spectrum is obtained from that of $S$ by negating every entry.}
\end{theorem}
\begin{proof}
\pfsketch\ We will check the conjugation entry by entry and apply the characteristic-polynomial certificate.
\step{seidel-conj}{$S(X^U)=TS(X)T$ and $T^2=I$.}
\begin{proof}
\step{seidel-entries}{\assume{$u,v\in V$.}\prove{$S(X^U)(u,v)=(TS(X)T)(u,v)=T(u,u)S(X)(u,v)T(v,v)$.}}
\begin{proof}
\step{seidel-diag}{\case{$u=v$.}}
\begin{proof}
\pf\ \raggedright Both entries are zero, since Seidel diagonal entries are zero by Definition~\ref{def:graph}.~\qed
\end{proof}
\step{seidel-inside}{\case{$u\ne v$ lie in the same part.}}
\begin{proof}
\pf\ \raggedright The sign product is $1$ and adjacency is unchanged by switching.~\qed
\end{proof}
\step{seidel-crossedge}{\case{$u,v$ lie in different parts and $uv$ is an edge of $X$.}}
\begin{proof}
\pf\ \raggedright Switching removes the edge, so the entry $-1$ becomes $1$; the sign product in the conjugation is $-1$.~\qed
\end{proof}
\step{seidel-crossnonedge}{\case{$u,v$ lie in different parts and $uv$ is not an edge of $X$.}}
\begin{proof}
\pf\ \raggedright Switching inserts the edge, so the entry $1$ becomes $-1$; the sign product is again $-1$.~\qed
\end{proof}
\qedstep
\begin{proof}
\pf\ \raggedright The four cases are exhaustive.~\qed
\end{proof}
\end{proof}
\step{seidel-square}{$T^2=I$.}
\begin{proof}
\pf\ \raggedright The off-diagonal entries vanish and each diagonal sign squares to one.~\qed
\end{proof}
\qedstep
\begin{proof}
\pf\ \raggedright Use \stepref{seidel-entries}, \stepref{seidel-square}.~\qed
\end{proof}
\end{proof}
\step{seidel-spectrum}{Seidel spectra with multiplicities are preserved, and $-S$ has the negated spectrum.}
\begin{proof}
\step{seidel-charpoly}{The two Seidel characteristic polynomials agree.}
\begin{proof}
\pf\ \raggedright Apply Theorem~\ref{thm:similarity} to \stepref{seidel-conj}.~\qed
\end{proof}
\step{seidel-multiset}{The two Seidel multisets agree.}
\begin{proof}
\pf\ \raggedright The matrices are real symmetric; use Proposition~\ref{prop:cospectral}.~\qed
\end{proof}
\step{seidel-minus}{$-S=2A-J+I$ and its eigenvalues are those of $S$ negated, with repetitions.}
\begin{proof}
\pf\ \raggedright Negate Definition~\ref{def:graph} and apply Theorem~\ref{thm:calculus} to $-t$.~\qed
\end{proof}
\qedstep
\begin{proof}
\pf\ \raggedright Use \stepref{seidel-charpoly}, \stepref{seidel-multiset}, \stepref{seidel-minus}.~\qed
\end{proof}
\end{proof}
\qedstep
\begin{proof}
\pf\ \raggedright Use \stepref{seidel-conj}, \stepref{seidel-spectrum}.~\qed
\end{proof}
\end{proof}

\begin{theorem}[Rational eigenvalues of integer matrices]\label{thm:rational}
\lean{Matrix.eigenvalue_isInteger_of_rat}
\uses{facts:charpoly,facts:det,facts:rational}
\mbox{}\par\nobreak
\assume{\begin{pfenum}
\item $U$ is finite, $A:U\times U\to\ints$, $b\in\rats$, and a nonzero real vector satisfies $A\Zv=b\Zv$ after casting $A,b$ to $\re$.
\end{pfenum}}
\prove{\raggedright If an integer square matrix has a rational eigenvalue $b$, viewed as
an eigenvalue of its real matrix, then $b$ is an integer.}
\end{theorem}
\begin{proof}
\pfsketch\ We will place the rational root in the monic integer characteristic polynomial and clear its reduced denominator.
\step{rat-root}{The integer characteristic polynomial $f$ is monic of degree $N>0$ and $f(b)=0$ over $\rats$.}
\begin{proof}
\step{rat-realzero}{The real determinant of $bI-A$ is zero.}
\begin{proof}
\pf\ \raggedright The given nonzero eigenvector lies in its kernel; the determinant criterion in Facts~\ref{facts:det} excludes invertibility.~\qed
\end{proof}
\step{rat-cast}{$f(b)=0$ in $\rats$.}
\begin{proof}
\pf\ \raggedright Use \code{Matrix.charpoly_map} in Facts~\ref{facts:charpoly} and injectivity of $\rats\hookrightarrow\re$ on \stepref{rat-realzero}.~\qed
\end{proof}
\step{rat-monic}{$f$ is monic and its degree $N$ is positive.}
\begin{proof}
\pf\ \raggedright Characteristic polynomials are monic (Facts~\ref{facts:charpoly}); a degree-zero monic polynomial is $1$ and cannot vanish at $b$.~\qed
\end{proof}
\qedstep
\begin{proof}
\pf\ \raggedright Use \stepref{rat-cast}, \stepref{rat-monic}.~\qed
\end{proof}
\end{proof}
\step{rat-integrality}{The rational root $b$ is an integer.}
\begin{proof}
\pf\ \raggedright Apply \code{isInteger_of_is_root_of_monic} in Facts~\ref{facts:rational} to the integer characteristic polynomial and the root equation in \stepref{rat-root}, with UFD $\ints$ and fraction field $\rats$.~\qed
\end{proof}
\qedstep
\begin{proof}
\pf\ \raggedright Use \stepref{rat-integrality}.~\qed
\end{proof}
\end{proof}

\begin{lemma}[Adjacency matrix guards]\label{lem:adjacency-guards}
\lean{Matrix.IsAdjMatrix.of_symmetric_zero_one}
\uses{facts:graph}
\mbox{}\par\nobreak
\assume{\begin{pfenum}
\item $U$ is finite and $A:U\times U\to\ints$ is symmetric, with zero diagonal and entries in $\{0,1\}$.
\end{pfenum}}
\prove{\raggedright The relation $A(u,v)=1$ defines a simple
graph whose integer adjacency matrix is $A$.}
\end{lemma}
\begin{proof}
\pfsketch\ We will recognize the three guards as the fields of Mathlib's \code{Matrix.IsAdjMatrix}, whose graph has adjacency $A(u,v)=1$.
\step{check-isadj}{$A$ satisfies \code{Matrix.IsAdjMatrix}.}
\begin{proof}
\pf\ \raggedright Its three fields, zero-one entries, symmetry and zero diagonal, are the three guards of \asm[1]{0} (Facts~\ref{facts:graph}).~\qed
\end{proof}
\step{check-adjacency}{The relation $A(u,v)=1$ is a simple graph $X_A$.}
\begin{proof}
\pf\ \raggedright Use \code{Matrix.IsAdjMatrix.toGraph} of \stepref{check-isadj} in Facts~\ref{facts:graph}: symmetry of $A$ makes the relation symmetric, and the zero diagonal makes it irreflexive.~\qed
\end{proof}
\step{check-adjmatrix}{$A(X_A)=A$.}
\begin{proof}
\pf\ \raggedright Apply \code{Matrix.IsAdjMatrix.adjMatrix_toGraph_eq} in Facts~\ref{facts:graph} to \stepref{check-isadj}; on each off-diagonal pair, the zero-one guard makes the test for equality with $1$ return the entry itself.~\qed
\end{proof}
\qedstep
\begin{proof}
\pf\ \raggedright Use \stepref{check-adjacency}, \stepref{check-adjmatrix}.~\qed
\end{proof}
\end{proof}

\begin{theorem}[A precise Boolean certificate interface]\label{thm:checker}
\lean{Matrix.spectrumCertificate_sound, Matrix.positiveEigenvectorCertificate_sound, Matrix.similarityCertificate_sound}
\uses{cor:positive-certificate,def:graph,lem:adjacency-guards,thm:certificate,thm:similarity,thm:smith-vectors}
\mbox{}\par\nobreak
\assume{\begin{pfenum}
\item Inputs are finite lists of integer rows, integer nodes, natural multiplicities, integer vectors or integer scalars as appropriate; all list lengths and index bounds are checked before arithmetic access.
\item The three Boolean procedures and recursive operations are specified below; every guard is combined by Boolean conjunction.
\end{pfenum}}
\prove{\raggedright Encode an integer $n\times n$ matrix by a list of $n$ rows of length
$n$. Consider the following three checkers, using exact integers and
natural numbers, with no floating-point operations.
\begin{enumerate}
\item A spectral checker first verifies the square shape and symmetry,
then a possibly empty list of distinct integer nodes $\theta_i$ and an
equally long list of natural multiplicities $m_i$. It verifies the
matrix product $\prod_i(A-\theta_iI)=0$ and the trace identities
of Theorem~\ref{thm:certificate} for $0\le j<s$.
\item A largest-eigenvalue checker verifies square shape, $n>0$,
symmetry, nonnegative entries, a length-$n$ integer vector with
strictly positive entries, and $A\Zx=b\Zx$ for the proposed integer $b$.
\item A similarity checker verifies that $A,B,Q,R$ have the same
square shape, that $d\ne0$, and the three identities in
Theorem~\ref{thm:similarity}.
\end{enumerate}
Returning \code{true} implies, respectively, the stated real multiset
spectrum, the stated largest eigenvalue, or equality of real
characteristic polynomials. Graph certificates additionally require zero
diagonal, zero-one entries, and symmetry for every proposed adjacency
matrix. In the similarity checker these guards apply to both $A$ and $B$.}
\end{theorem}
\begin{proof}
\pfsketch\ We will prove the list operations agree with matrix arithmetic, maintain the product and power invariants, cast to real matrices, and apply the three soundness theorems.
\step{check-dot}{The recursive list dot product computes the finite coordinate sum.}
\begin{proof}
\step{check-dotdef}{Define $\operatorname{dot}([],[])=0$ and $\operatorname{dot}(h::v,k::w)=hk+\operatorname{dot}(v,w)$ for equally long lists.}
\begin{proof}
\pf\ \raggedright These equations are the specification of the exact integer operation.~\qed
\end{proof}
\step{check-dotbase}{For length zero, $\operatorname{dot}([],[])=\sum_{i\in\operatorname{Fin}(0)}v(i)w(i)$.}
\begin{proof}
\pf\ \raggedright Both sides are zero by \stepref{check-dotdef} and the empty-sum definition.~\qed
\end{proof}
\step{check-dotstep}{If the equality holds on the tails, it holds on the full lists.}
\begin{proof}
\pf\ \raggedright Both sides split into the same first term $hk$ and the tail sum in \stepref{check-dotdef}.~\qed
\end{proof}
\qedstep
\begin{proof}
\pf\ \raggedright Induction on the common list length using \stepref{check-dotbase}, \stepref{check-dotstep}.~\qed
\end{proof}
\end{proof}
\step{check-operations}{The guarded operations compute the exact matrices and Boolean predicates.}
\begin{proof}
\step{check-shape}{A true shape guard defines a matrix on $\operatorname{Fin}(n)$ by row and column lookup.}
\begin{proof}
\pf\ \raggedright Each of the $n$ rows has length $n$, so every bounded pair of indices has an integer entry.~\qed
\end{proof}
\step{check-mul}{Computing each entry as a row-column dot product gives matrix multiplication.}
\begin{proof}
\pf\ \raggedright Apply \stepref{check-dot} to each row and column; its sum is the defining matrix product.~\qed
\end{proof}
\step{check-entryops}{Entrywise addition and scalar multiplication, the Kronecker-delta identity, and the diagonal sum give the named matrix operations.}
\begin{proof}
\pf\ \raggedright These are their coordinate definitions.~\qed
\end{proof}
\step{check-booleans}{List equality and a universal list test reflect equality and universal quantification.}
\begin{proof}
\step{check-boolbase}{The empty conjunction is true and the empty equal-length lists agree.}
\begin{proof}
\pf\ \raggedright These are the base equations for the Boolean tests.~\qed
\end{proof}
\step{check-boolstep}{A nonempty conjunction is true iff its first test and tail conjunction are true.}
\begin{proof}
\pf\ \raggedright The Boolean truth table for conjunction.~\qed
\end{proof}
\step{check-equality}{An equality test is true iff lengths agree and all corresponding integers agree.}
\begin{proof}
\pf\ \raggedright After the length guard, induct using \stepref{check-boolbase}, \stepref{check-boolstep} and exact integer equality.~\qed
\end{proof}
\qedstep
\begin{proof}
\pf\ \raggedright List induction proves the universal test; \stepref{check-equality} proves the equality test.~\qed
\end{proof}
\end{proof}
\qedstep
\begin{proof}
\pf\ \raggedright Use \stepref{check-shape}, \stepref{check-mul}, \stepref{check-entryops}, \stepref{check-booleans}.~\qed
\end{proof}
\end{proof}
\step{check-invariants}{The product accumulator and power accumulator compute exactly the annihilator and all required powers.}
\begin{proof}
\step{check-product0}{The initial product accumulator $I$ is the empty product.}
\begin{proof}
\pf\ \raggedright This is the definition of a product over an empty prefix.~\qed
\end{proof}
\step{check-productS}{Multiplying by the next $A-\theta_iI$ appends that factor to the prefix product.}
\begin{proof}
\pf\ \raggedright Use the product recursion and exact multiplication from \stepref{check-operations}.~\qed
\end{proof}
\step{check-power0}{$P_0=I=A^0$.}
\begin{proof}
\pf\ \raggedright Definition of the initial accumulator and of zeroth power.~\qed
\end{proof}
\step{check-powerS}{If $P_j=A^j$, then $P_{j+1}=P_jA=A^{j+1}$.}
\begin{proof}
\pf\ \raggedright Substitute the induction hypothesis in the specified recurrence.~\qed
\end{proof}
\step{check-traces}{The trace tests are precisely the required moment equations, including degree zero.}
\begin{proof}
\pf\ \raggedright Induct with \stepref{check-power0}, \stepref{check-powerS} and use the diagonal sum in \stepref{check-operations}.~\qed
\end{proof}
\qedstep
\begin{proof}
\pf\ \raggedright Induction on the processed nodes with \stepref{check-product0} and \stepref{check-productS} gives the annihilator; \stepref{check-traces} gives the moments.~\qed
\end{proof}
\end{proof}
\step{check-real}{True checks give all real hypotheses of the three certificate theorems.}
\begin{proof}
\step{check-cast-algebra}{Integer matrix equations and moments remain valid over $\re$.}
\begin{proof}
\pf\ \raggedright The inclusion $\ints\hookrightarrow\re$ preserves each finite sum and product defining these operations.~\qed
\end{proof}
\step{check-cast-order}{Nonnegative entries, positive coordinates, distinct nodes and a nonzero scalar remain so after casting.}
\begin{proof}
\pf\ \raggedright The integer inclusion is order preserving and injective.~\qed
\end{proof}
\step{check-spectrum}{The spectral checker returning true gives the stated multiset.}
\begin{proof}
\pf\ \raggedright Apply Theorem~\ref{thm:certificate} to \stepref{check-operations}, \stepref{check-invariants}, \stepref{check-cast-algebra} and \stepref{check-cast-order}.~\qed
\end{proof}
\step{check-largest}{The positive-vector checker returning true gives the stated largest eigenvalue.}
\begin{proof}
\pf\ \raggedright Apply Corollary~\ref{cor:positive-certificate} to the reflected shape, symmetry, nonnegativity, positivity and eigenvector checks.~\qed
\end{proof}
\step{check-similarity}{The similarity checker returning true gives equal real characteristic polynomials.}
\begin{proof}
\pf\ \raggedright Apply Theorem~\ref{thm:similarity} to the reflected three identities and $d\ne0$.~\qed
\end{proof}
\qedstep
\begin{proof}
\pf\ \raggedright Use \stepref{check-spectrum}, \stepref{check-largest}, \stepref{check-similarity}.~\qed
\end{proof}
\end{proof}
\step{check-graph}{With the symmetry, zero-diagonal and zero-one guards on every proposed adjacency matrix, the checkers prove the corresponding graph statements; without the symmetry guard they need not.}
\begin{proof}
\step{check-missingguard}{Zero diagonal and zero-one entries alone do not ensure an undirected graph in the similarity checker.}
\begin{proof}
\pf\ \raggedright Take $A=B=\pmat{0&1\\0&0}$, $Q=R=I$, $d=1$: all similarity identities and both entry guards hold, but $A(1,2)\ne A(2,1)$.~\qed
\end{proof}
\step{check-onegraph}{For the spectral and positive-vector checkers, the guards on $A$ give a graph with the stated spectrum or largest eigenvalue.}
\begin{proof}
\pf\ \raggedright Apply Lemma~\ref{lem:adjacency-guards} to $A$; the real adjacency matrix of the resulting graph is the cast of $A$, so \stepref{check-real} applies to it; use Definition~\ref{def:graph}.~\qed
\end{proof}
\step{check-both}{For the similarity checker, the guards on both $A$ and $B$ give cospectral graphs.}
\begin{proof}
\pf\ \raggedright Apply Lemma~\ref{lem:adjacency-guards} separately to $A$ and $B$; the real adjacency matrices of the two graphs are the casts of $A$ and $B$, so \stepref{check-real} gives equal characteristic polynomials, which is cospectrality by Definition~\ref{def:graph}.~\qed
\end{proof}
\qedstep
\begin{proof}
\pf\ \raggedright Use \stepref{check-missingguard}, \stepref{check-onegraph}, \stepref{check-both}.~\qed
\end{proof}
\end{proof}
\qedstep
\begin{proof}
\pf\ \raggedright Use \stepref{check-real}, \stepref{check-graph}.~\qed
\end{proof}
\end{proof}

The concrete C13, C15 and C17 data are those of the ordinary version.
The sign convention is $S=J-I-2A$, as in SW2 of \code{cgss/prerequisites.md}. For C13 the adjacency nodes are $4,2,0,-2$, with multiplicities $1,3,3,5$
and moments $12,0,48,48$. Its Seidel nodes are $-5,-1,3$, with
multiplicities $3,3,6$ and moments $12,0,132$; negating the Seidel
matrix negates these nodes. For C15 the nodes are $6,4,1,0,-2$,
with multiplicities $1,5,4,5,15$ and moments $30,0,180,420,2820$.
C17 uses the positive integer vectors in Theorem~\ref{thm:smith-vectors}.
The checks in \path{checks/check_spectraltools.py} verify these finite
identities using exact integers. These data are not themselves a Lean
proof.

\section{Further spectral applications}

\begin{theorem}[Scaled idempotents and singular principal minors]\label{thm:scaled}
\lean{Matrix.IsHermitian.posSemidef_of_sq_eq_pos_smul, Matrix.IsHermitian.rank_eq_trace_div_of_sq_eq_smul, Matrix.det_principal_submatrix_eq_zero_of_rank_lt}
\uses{facts:det,facts:dim,thm:calculus,thm:psd,thm:trace}
\mbox{}\par\nobreak
\assume{\begin{pfenum}
\item $U$ is finite of order $n$, $C:U\times U\to\re$ is symmetric, $b>0$ and $C^2=bC$.
\item For the rank bound $Q$ is finite, $B:U\times Q\to\re$, and $CB=0$. A principal submatrix is indexed by a subset of $U$.
\end{pfenum}}
\prove{\raggedright $C\succeq0$, its eigenvalues lie in $\{0,b\}$,
and $\rk C=\tr C/b$. If $B$ is an $n\times q$ matrix with $CB=0$,
then $\rk C\le n-\rk B$. Every principal submatrix of order
greater than $\rk C$ is singular.}
\end{theorem}
\begin{proof}
\pfsketch\ We will read the two possible eigenvalues, compare image with kernel, and bound the rank of each coordinate restriction.
\step{scaled-values}{$C\succeq0$ and $\rk C=\tr C/b$.}
\begin{proof}
\step{scaled-roots}{Every eigenvalue $\mu$ satisfies $\mu(\mu-b)=0$.}
\begin{proof}
\pf\ \raggedright Apply Theorem~\ref{thm:calculus} to $C^2-bC=0$.~\qed
\end{proof}
\step{scaled-positive}{Every eigenvalue is $0$ or $b>0$, so $C\succeq0$.}
\begin{proof}
\pf\ \raggedright Factor cancellation over $\re$ in \stepref{scaled-roots} and Theorem~\ref{thm:psd}.~\qed
\end{proof}
\step{scaled-trace}{$\tr C=b\rk C$.}
\begin{proof}
\pf\ \raggedright Theorem~\ref{thm:trace} sums the $b$ entries and counts precisely those entries as nonzero.~\qed
\end{proof}
\step{scaled-divide}{$\rk C=\tr C/b$.}
\begin{proof}
\pf\ \raggedright Divide \stepref{scaled-trace} by $b>0$.~\qed
\end{proof}
\qedstep
\begin{proof}
\pf\ \raggedright Use \stepref{scaled-positive}, \stepref{scaled-divide}.~\qed
\end{proof}
\end{proof}
\step{scaled-null}{$\rk C\le n-\rk B$.}
\begin{proof}
\step{scaled-image}{$\im B\subseteq\ker C$.}
\begin{proof}
\pf\ \raggedright For every $\Zx$, $C(B\Zx)=(CB)\Zx=0$.~\qed
\end{proof}
\step{scaled-dim}{$\rk B\le\dim\ker C=n-\rk C$.}
\begin{proof}
\pf\ \raggedright Dimension monotonicity and rank-nullity in Facts~\ref{facts:dim}.~\qed
\end{proof}
\qedstep
\begin{proof}
\pf\ \raggedright Rearrange the dimension inequality in \stepref{scaled-dim} over the integers.~\qed
\end{proof}
\end{proof}
\step{scaled-minor}{Every principal submatrix of order greater than $\rk C$ is singular.}
\begin{proof}
\step{scaled-factor}{A principal submatrix factors as coordinate inclusion, $C$, and coordinate projection.}
\begin{proof}
\pf\ \raggedright Its entry at retained indices is the corresponding entry of $C$.~\qed
\end{proof}
\step{scaled-rank}{Its rank is at most $\rk C$.}
\begin{proof}
\pf\ \raggedright The image in \stepref{scaled-factor} lies in the projection of $\im C$, whose dimension cannot exceed $\dim\im C$ (Facts~\ref{facts:dim}).~\qed
\end{proof}
\step{scaled-determinant}{If its order exceeds that rank, its determinant is zero.}
\begin{proof}
\pf\ \raggedright Rank-nullity gives a nonzero kernel; use the determinant criterion in Facts~\ref{facts:det}.~\qed
\end{proof}
\qedstep
\begin{proof}
\pf\ \raggedright Use \stepref{scaled-determinant}.~\qed
\end{proof}
\end{proof}
\qedstep
\begin{proof}
\pf\ \raggedright Use \stepref{scaled-values}, \stepref{scaled-null}, \stepref{scaled-minor}.~\qed
\end{proof}
\end{proof}

\begin{theorem}[Symmetric matrices and rank-one Gram matrices]\label{thm:symmetric-space}
\lean{Matrix.finrank_symmetric, Matrix.finrank_traceless_symmetric, Matrix.gram_rankOne_eq_inner_sq, SimpleGraph.card_le_symm_dimension_of_representation}
\uses{facts:counting,facts:dim,facts:order,thm:gram,thm:psd,thm:trace}
\mbox{}\par\nobreak
\assume{\begin{pfenum}
\item $\ell\in\nn$ and the matrix space consists of all real symmetric $\ell\times\ell$ matrices.
\item The traceless dimension assertion assumes $\ell>0$. For the Gram identity $(\Zv_u)_{u\in U}$ is a finite family of arbitrary real vectors in $\re^\ell$, with $|U|=n$.
\item The last assertion assumes unit vectors and $\ip{\Zv_u}{\Zv_v}\in\{0,1/2\}$ for all distinct $u,v\in U$.
\end{pfenum}}
\prove{\raggedright The real symmetric $\ell\times\ell$ matrices form an inner product
space for $\ip{P}{Q}=\tr(PQ)$, of dimension $\ell(\ell+1)/2$.
For $\ell>0$ the traceless subspace has dimension
$\ell(\ell+1)/2-1$. For arbitrary real vectors $\Zv_u$, their rank-one
matrices $P_u=\Zv_u\Zv_u^T$ have Gram matrix
$\tr(P_uP_v)=\ip{\Zv_u}{\Zv_v}^2$.
Thus, if $n$ unit vectors in $\re^\ell$ have distinct-pair
inner products in $\{0,1/2\}$, then $n\le\ell(\ell+1)/2$.}
\end{theorem}
\begin{proof}
\pfsketch\ We will use coordinates on unordered pairs for the dimension and a positive definite Gram matrix for the rank-one family.
\step{symm-inner}{$\ip{P}{Q}=\tr(PQ)$ is an inner product on the symmetric matrices.}
\begin{proof}
\step{symm-coordinate}{$\tr(PQ)=\sum_{i,j}P(i,j)Q(i,j)$.}
\begin{proof}
\pf\ \raggedright Expand trace and multiplication, using $Q(j,i)=Q(i,j)$.~\qed
\end{proof}
\step{symm-bilinear}{This form is symmetric and bilinear.}
\begin{proof}
\pf\ \raggedright Each coordinate product in \stepref{symm-coordinate} is symmetric bilinear; finite sums preserve these equalities.~\qed
\end{proof}
\step{symm-definite}{$\tr(P^2)=\sum_{i,j}P(i,j)^2\ge0$, with equality iff $P=0$.}
\begin{proof}
\pf\ \raggedright Facts~\ref{facts:order} applied to the coordinate squares in \stepref{symm-coordinate}.~\qed
\end{proof}
\qedstep
\begin{proof}
\pf\ \raggedright Use \stepref{symm-bilinear}, \stepref{symm-definite}.~\qed
\end{proof}
\end{proof}
\step{symm-dimension}{The dimension is $\ell(\ell+1)/2$, and for $\ell>0$ the traceless dimension is one less.}
\begin{proof}
\step{symm-basis}{Use one diagonal unit matrix for each index and one symmetric off-diagonal unit pair for each unordered distinct pair.}
\begin{proof}
\pf\ \raggedright These are symmetric matrices by their entry definitions.~\qed
\end{proof}
\step{symm-independent}{This family is independent.}
\begin{proof}
\pf\ \raggedright A coefficient in a zero combination is recovered by reading its unique diagonal or unordered-pair entry.~\qed
\end{proof}
\step{symm-span}{It spans every symmetric matrix.}
\begin{proof}
\pf\ \raggedright Take the coefficients to be its diagonal and one entry of each off-diagonal pair.~\qed
\end{proof}
\step{symm-card}{The number of these matrices is $\ell+\binom\ell2=\ell(\ell+1)/2$.}
\begin{proof}
\pf\ \raggedright Apply \code{Sym2.card} to their index type of unordered pairs with repetitions in Facts~\ref{facts:counting}.~\qed
\end{proof}
\step{symm-tracesurj}{For $\ell>0$, the trace map onto $\re$ is surjective.}
\begin{proof}
\pf\ \raggedright A diagonal unit matrix has trace one; multiplying it by any real number gives that trace.~\qed
\end{proof}
\step{symm-traceless}{The trace kernel has dimension $\ell(\ell+1)/2-1$.}
\begin{proof}
\pf\ \raggedright Rank-nullity in Facts~\ref{facts:dim}, since \stepref{symm-tracesurj} gives image dimension one.~\qed
\end{proof}
\qedstep
\begin{proof}
\pf\ \raggedright Use \stepref{symm-independent}, \stepref{symm-span}, \stepref{symm-card}, \stepref{symm-traceless}.~\qed
\end{proof}
\end{proof}
\step{symm-rankone}{For $P_u=\Zv_u\Zv_u^T$, $\tr(P_uP_v)=\ip{\Zv_u}{\Zv_v}^2$.}
\begin{proof}
\step{symm-expand}{$\tr(P_uP_v)=\sum_{i,j}\Zv_u(i)\Zv_u(j)\Zv_v(j)\Zv_v(i)$.}
\begin{proof}
\pf\ \raggedright Expand the two rank-one matrices, multiplication and trace.~\qed
\end{proof}
\step{symm-factor}{This sum equals $(\sum_i\Zv_u(i)\Zv_v(i))^2$.}
\begin{proof}
\pf\ \raggedright Reorder real factors and distribute the square of the finite sum.~\qed
\end{proof}
\qedstep
\begin{proof}
\pf\ \raggedright Use \stepref{symm-expand}, \stepref{symm-factor}.~\qed
\end{proof}
\end{proof}
\step{symm-bound}{$n\le\ell(\ell+1)/2$ under the final hypothesis.}
\begin{proof}
\pflet{$G\triangleq(\ip{\Zv_u}{\Zv_v})_{u,v\in U}$ is the Gram matrix of the vectors, and $H\triangleq(\tr(P_uP_v))_{u,v\in U}$ that of the rank-one matrices.}
\step{symm-original}{$H=(I+G)/2$.}
\begin{proof}
\pf\ \raggedright On the diagonal both are $1$; off the diagonal, squaring $0$ or $1/2$ gives half its value, by \stepref{symm-rankone}.~\qed
\end{proof}
\step{symm-positive}{For $\Zx\ne0$, $\Zx^T(I+G)\Zx/2>0$.}
\begin{proof}
\pf\ \raggedright $G\succeq0$ by Theorem~\ref{thm:gram}, and $\norm{\Zx}^2>0$.~\qed
\end{proof}
\step{symm-fullrank}{$H$ has rank $n$.}
\begin{proof}
\pf\ \raggedright By \stepref{symm-original} and \stepref{symm-positive}, $H$ is positive definite; Theorem~\ref{thm:psd} gives $n$ positive eigenvalues, and Theorem~\ref{thm:trace} counts them.~\qed
\end{proof}
\step{symm-spanbound}{The span of the $P_u$ has dimension $n$, at most the symmetric-space dimension.}
\begin{proof}
\pf\ \raggedright Use Theorem~\ref{thm:gram}, \stepref{symm-fullrank}, and \stepref{symm-dimension}.~\qed
\end{proof}
\step{symm-empty}{The argument includes $n=0$; if $\ell=0$, a unit-vector family must be empty.}
\begin{proof}
\pf\ \raggedright The zero space has no unit vector, and its empty family has rank and cardinality zero.~\qed
\end{proof}
\qedstep
\begin{proof}
\pf\ \raggedright Use \stepref{symm-spanbound}, \stepref{symm-empty}.~\qed
\end{proof}
\end{proof}
\qedstep
\begin{proof}
\pf\ \raggedright Use \stepref{symm-inner}, \stepref{symm-dimension}, \stepref{symm-rankone}, \stepref{symm-bound}.~\qed
\end{proof}
\end{proof}

\begin{definition}\label{def:hadamard-pairs}
\lean{Matrix.hadamardPairs, Matrix.hadamardSignedAdj}
\uses{}
\mbox{}\par\nobreak
\assume{\begin{pfenum}
\item $H$ is a real matrix on finite row and column types; the signed matrix uses square $H$ of order $n$ and $x>0$.
\end{pfenum}}
For a real matrix $H$ on a finite row type $U$ and column type $V$,
let $P=\code{Matrix.hadamardPairs}\ H$ have rows indexed by the
unordered two-element subsets $\{i,j\}\subseteq U$, with
$P(\{i,j\},u)=H(i,u)H(j,u)$. This does not depend on the ordering
of $i,j$. For square $H$ of order $n$ and real $x>0$, define
$\code{Matrix.hadamardSignedAdj}\ H\ x=(PP^T-nI)/x$ and write it $C_H$.
The name does not by itself assert that the entries are $0,1,-1$;
that requires an additional hypothesis.
\end{definition}

\begin{theorem}[The Hadamard pair matrix]\label{thm:hadamard}
\lean{Matrix.IsHadamard.hadamardPairs_transpose_mul, Matrix.IsHadamard.spectrumMultiset_hadamardPairs_gram, Matrix.IsHadamard.dvd_order_of_integral_hadamardSignedAdj}
\uses{cor:scalar-ones,facts:dim,thm:calculus,thm:gram,thm:mn,thm:rational,thm:trace}
\mbox{}\par\nobreak
\assume{\begin{pfenum}
\item $n\ge2$ and $H$ is a real $n\times n$ matrix with entries $\pm1$ and $HH^T=nI$.
\item $P$ has one row for each unordered pair of distinct row indices, $P(\{i,j\},u)=H(i,u)H(j,u)$, and $m=\binom n2$.
\item For the scaled assertions $x>0$ and $C_H=(PP^T-nI)/x$. For divisibility $n>2$, $x$ is a positive integer, and $C_H$ has integer entries.
\end{pfenum}}
\prove{\raggedright \[
 P^TP=\frac{n^2}{2}I-\frac n2J,\qquad
 \Spec(PP^T)=\ms{(n^2/2)^{n-1},0^{m-n+1}},\qquad
 \rk P=n-1.
\]
For a positive real $x$, put $C_H=(PP^T-nI)/x$. Then
\[
 (xC_H+nI)\left(xC_H-\frac{n(n-2)}2I\right)=0,
\]
and its spectrum consists of $n(n-2)/(2x)$ with multiplicity
$n-1$ and $-n/x$ with multiplicity $m-n+1$.
If $n>2$, $x$ is a positive integer, and $C_H$ has integer
entries, then $x\mid n$.}
\end{theorem}
\begin{proof}
\pfsketch\ We will derive column orthogonality, expand a sum over row pairs, transfer the nonzero Gram spectrum, and use an occurring rational eigenvalue.
\step{had-columns}{$H^TH=nI$.}
\begin{proof}
\step{had-rightinverse}{$H(H^T/n)=I$.}
\begin{proof}
\pf\ \raggedright Divide the assumed row-orthogonality equation by $n>0$.~\qed
\end{proof}
\step{had-surjective}{$H$ is surjective.}
\begin{proof}
\pf\ \raggedright For any $\Zy$, the vector $(H^T/n)\Zy$ maps to $\Zy$ by \stepref{had-rightinverse}.~\qed
\end{proof}
\step{had-invertible}{$H$ is invertible with inverse $H^T/n$.}
\begin{proof}
\pf\ \raggedright A surjective endomorphism in finite equal dimension is injective (Facts~\ref{facts:dim}); its right inverse is its inverse.~\qed
\end{proof}
\step{had-leftinverse}{$H^TH=nI$.}
\begin{proof}
\pf\ \raggedright Multiply the left-inverse identity from \stepref{had-invertible} by $n$.~\qed
\end{proof}
\qedstep
\begin{proof}
\pf\ \raggedright Use \stepref{had-leftinverse}.~\qed
\end{proof}
\end{proof}
\step{had-pairproduct}{$P^TP=(n^2/2)I-(n/2)J$.}
\begin{proof}
\step{had-pairsum}{For fixed $u,v$, put $b_i=H(i,u)H(i,v)$; then $(P^TP)(u,v)=\sum_{i<j}b_ib_j$.}
\begin{proof}
\pf\ \raggedright Expand the product and use the unordered distinct-pair indexing of $P$.~\qed
\end{proof}
\step{had-squareexpand}{$2\sum_{i<j}b_ib_j=(\sum_i b_i)^2-\sum_i b_i^2$.}
\begin{proof}
\pf\ \raggedright Split the double sum into its diagonal and the two orientations of each distinct pair.~\qed
\end{proof}
\step{had-bsum}{$\sum_i b_i=n\delta_{uv}$ and $\sum_i b_i^2=n$.}
\begin{proof}
\pf\ \raggedright The first is \stepref{had-columns}; each squared product of two signs equals one for the second.~\qed
\end{proof}
\step{had-entries}{$(P^TP)(u,v)=(n^2\delta_{uv}-n)/2$.}
\begin{proof}
\pf\ \raggedright Substitute \stepref{had-bsum} in \stepref{had-squareexpand} and divide by two; $\delta_{uv}^2=\delta_{uv}$.~\qed
\end{proof}
\qedstep
\begin{proof}
\pf\ \raggedright Use \stepref{had-entries}.~\qed
\end{proof}
\end{proof}
\step{had-spectrum}{$\rk P=n-1$ and $\Spec(PP^T)=\ms{(n^2/2)^{n-1},0^{m-n+1}}$.}
\begin{proof}
\step{had-list}{$P^TP$ has eigenvalue zero once and $n^2/2>0$ with multiplicity $n-1$.}
\begin{proof}
\pf\ \raggedright Apply Corollary~\ref{cor:scalar-ones} to \stepref{had-pairproduct}.~\qed
\end{proof}
\step{had-rank}{$\rk P=n-1$.}
\begin{proof}
\pf\ \raggedright Theorems~\ref{thm:gram} and \ref{thm:trace} applied to \stepref{had-list}.~\qed
\end{proof}
\step{had-transfer}{$PP^T$ has the stated positive list and $m-(n-1)$ zeros.}
\begin{proof}
\pf\ \raggedright Apply Theorem~\ref{thm:mn} to \stepref{had-list}, \stepref{had-rank}.~\qed
\end{proof}
\step{had-zerocount}{$m-n+1=(n-1)(n-2)/2\ge0$, with strict positivity for $n>2$.}
\begin{proof}
\pf\ \raggedright Substitute $m=n(n-1)/2$ and factor.~\qed
\end{proof}
\qedstep
\begin{proof}
\pf\ \raggedright Use \stepref{had-rank}, \stepref{had-transfer}, \stepref{had-zerocount}.~\qed
\end{proof}
\end{proof}
\step{had-scaled}{The polynomial identity and spectrum of $C_H$ hold.}
\begin{proof}
\step{had-annihilator}{$PP^T(PP^T-(n^2/2)I)=0$.}
\begin{proof}
\pf\ \raggedright The polynomial vanishes at both values in \stepref{had-spectrum}; use Theorem~\ref{thm:calculus}.~\qed
\end{proof}
\step{had-substitute}{$(xC_H+nI)(xC_H-n(n-2)I/2)=0$.}
\begin{proof}
\pf\ \raggedright Substitute $PP^T=xC_H+nI$ in \stepref{had-annihilator}.~\qed
\end{proof}
\step{had-scalelist}{The two values become $n(n-2)/(2x)$ and $-n/x$, with unchanged multiplicities.}
\begin{proof}
\pf\ \raggedright Apply Theorem~\ref{thm:calculus} to $(z-n)/x$, using $x>0$.~\qed
\end{proof}
\qedstep
\begin{proof}
\pf\ \raggedright Use \stepref{had-substitute}, \stepref{had-scalelist}.~\qed
\end{proof}
\end{proof}
\step{had-divisibility}{Under the integer hypotheses, $x\mid n$.}
\begin{proof}
\step{had-actualroot}{$-n/x$ actually occurs as a rational eigenvalue of $C_H$.}
\begin{proof}
\pf\ \raggedright Its multiplicity is positive for $n>2$ because $m-n+1=(n-1)(n-2)/2>0$ in \stepref{had-spectrum}, \stepref{had-scaled}; $n,x$ are integers and $x\ne0$.~\qed
\end{proof}
\step{had-integralroot}{$-n/x=z$ for some integer $z$.}
\begin{proof}
\pf\ \raggedright Apply Theorem~\ref{thm:rational} to the integer matrix $C_H$ and \stepref{had-actualroot}.~\qed
\end{proof}
\step{had-divide}{$n=x(-z)$ in $\ints$.}
\begin{proof}
\pf\ \raggedright Clear the nonzero denominator in \stepref{had-integralroot} and use injectivity of the integer cast; this is the divisibility witness.~\qed
\end{proof}
\qedstep
\begin{proof}
\pf\ \raggedright Use \stepref{had-divide}.~\qed
\end{proof}
\end{proof}
\qedstep
\begin{proof}
\pf\ \raggedright Use \stepref{had-pairproduct}, \stepref{had-spectrum}, \stepref{had-scaled}, \stepref{had-divisibility}.~\qed
\end{proof}
\end{proof}

\begin{proposition}[A row-sum assertion in Section 6]\label{rem:hadamard-row}
\lean{Matrix.IsHadamard.hadamardPairs_row_sum_counterexample}
\uses{thm:gram,thm:hadamard}
\mbox{}\par\nobreak
\assume{\begin{pfenum}
\item $H$ is the real matrix displayed below, $\Zh_i$ is its row $i$, and $n=x=4$, $P$ is its matrix of unordered distinct row products, $m=6$, and $C_H=(PP^T-4I)/4$.
\end{pfenum}}
\prove{\raggedright The matrix
\[
 H=\pmat{1&1&1&1\\1&-1&1&-1\\1&1&-1&-1\\1&-1&-1&1}
\]
is Hadamard and its unique four-row set agrees in $(n+x)/8=1$ columns,
but $xC_HJ=-nJ$ is false. The valid identity $P\Zj_n=0$ does not
imply $P^T\Zj_m=0$. In particular the row-sum assertion in Section 6
of CGSS is not a consequence of these hypotheses.
More generally, for every real $m\times n$ matrix $P$, $x>0$ and
$C=(PP^T-nI_m)/x$, we have
\[
 xCJ=-nJ\quad\Longleftrightarrow\quad PP^T\Zj_m=0
 \quad\Longleftrightarrow\quad P^T\Zj_m=0.
\]
If the rows of $P$ are the distinct-row products of a real square matrix
$H$ with $HH^T=nI_n$, then $P\Zj_n=0$.
}
\end{proposition}
\begin{proof}
\pfsketch\ We will list all six pair rows and read the three two-by-two blocks. We will also check the four-row agreement count.
\step{row-Hadamard}{The displayed $H$ has entries $\pm1$ and $HH^T=4I$.}
\begin{proof}
\step{row-norms}{Each row has squared norm four.}
\begin{proof}
\pf\ \raggedright Every row has four sign entries.~\qed
\end{proof}
\step{row-orthogonality}{The six distinct-row products are the balanced sign rows $\Zh_2,\Zh_3,\Zh_4,\Zh_4,\Zh_3,\Zh_2$, each of coordinate sum zero.}
\begin{proof}
\pf\ \raggedright Multiply the four displayed entries for each of the six unordered row pairs.~\qed
\end{proof}
\step{row-GramH}{$HH^T=4I$.}
\begin{proof}
\pf\ \raggedright The diagonal values in \stepref{row-norms} and off-diagonal sums in \stepref{row-orthogonality} give every entry.~\qed
\end{proof}
\qedstep
\begin{proof}
\pf\ \raggedright Use \stepref{row-GramH}.~\qed
\end{proof}
\end{proof}
\step{row-pairs}{The rows of $P$ consist of each of the three nonconstant rows of $H$ twice.}
\begin{proof}
\step{row-firstpairs}{Writing $\Zh_i$ for row $i$, the pairs $(1,2),(1,3),(1,4)$ give $\Zh_2,\Zh_3,\Zh_4$.}
\begin{proof}
\pf\ \raggedright Row one consists of ones.~\qed
\end{proof}
\step{row-23}{Pair $(2,3)$ gives $(1,-1,-1,1)=\Zh_4$.}
\begin{proof}
\pf\ \raggedright Coordinate multiplication of the displayed rows.~\qed
\end{proof}
\step{row-24}{Pair $(2,4)$ gives $(1,1,-1,-1)=\Zh_3$.}
\begin{proof}
\pf\ \raggedright Coordinate multiplication of the displayed rows.~\qed
\end{proof}
\step{row-34}{Pair $(3,4)$ gives $(1,-1,1,-1)=\Zh_2$.}
\begin{proof}
\pf\ \raggedright Coordinate multiplication of the displayed rows.~\qed
\end{proof}
\qedstep
\begin{proof}
\pf\ \raggedright Use \stepref{row-firstpairs}, \stepref{row-23}, \stepref{row-24}, \stepref{row-34}.~\qed
\end{proof}
\end{proof}
\step{row-counterexample}{$C_H$ is, after reordering, three blocks $\pmat{0&1\\1&0}$; thus $xC_HJ=4J\ne-4J$.}
\begin{proof}
\step{row-GramP}{$PP^T$ is three blocks $4J_2$.}
\begin{proof}
\pf\ \raggedright The repeated equal rows in \stepref{row-pairs} have inner product four, and different nonconstant rows are orthogonal by \stepref{row-Hadamard}.~\qed
\end{proof}
\step{row-C}{$C_H$ has the three stated blocks.}
\begin{proof}
\pf\ \raggedright Subtract $4I$ and divide by four in \stepref{row-GramP}.~\qed
\end{proof}
\step{row-sum}{Each row sum of $C_H$ is $1$, so $xC_HJ=4J$.}
\begin{proof}
\pf\ \raggedright Each block in \stepref{row-C} has one entry one in each row.~\qed
\end{proof}
\step{row-agreement}{The only four-row set agrees in exactly one column, equal to $(n+x)/8$.}
\begin{proof}
\pf\ \raggedright Column one is all ones; each other column has both signs. Here $(4+4)/8=1$.~\qed
\end{proof}
\qedstep
\begin{proof}
\pf\ \raggedright Use \stepref{row-GramP}, \stepref{row-C}, \stepref{row-sum}, \stepref{row-agreement}.~\qed
\end{proof}
\end{proof}
\step{row-dimensions}{\assume{$P$ is any real $m\times n$ matrix, $x>0$, and $C=(PP^T-nI_m)/x$.}\prove{$xCJ=-nJ$ iff $PP^T\Zj_m=0$, iff $P^T\Zj_m=0$. If the rows of $P$ are distinct-row products of a real square matrix $H$ with $HH^T=nI_n$, then $P\Zj_n=0$.}}
\begin{proof}
\step{row-right}{\assume{$P$ consists of the distinct-row products of a real matrix $H$ with $HH^T=nI_n$.}\prove{$P\Zj_n=0$.}}
\begin{proof}
\pf\ \raggedright Each row of $P$ is a product of two distinct orthogonal rows of $H$; its coordinate sum is their inner product.~\qed
\end{proof}
\step{row-left}{$xCJ=-nJ$ is equivalent to $PP^T\Zj_m=0$, equivalently $P^T\Zj_m=0$.}
\begin{proof}
\pf\ \raggedright Substitute the definition of $C$ from \asm{1} and use the kernel equality in Theorem~\ref{thm:gram} with matrix $P^T$.~\qed
\end{proof}
\qedstep
\begin{proof}
\pf\ \raggedright Use \stepref{row-right}, \stepref{row-left}.~\qed
\end{proof}
\end{proof}
\qedstep
\begin{proof}
\pf\ \raggedright Use \stepref{row-counterexample}, \stepref{row-dimensions}.~\qed
\end{proof}
\end{proof}

Error 9 of \code{cgss/errata.md}~\cite{errata} records the row-sum error. The shorter divisibility repair is Theorem~\ref{thm:hadamard}: $-n/x$ occurs with positive multiplicity $(n-1)(n-2)/2$ when $n>2$, so rational-root integrality gives $x\mid n$.
\section{Further directions}
Not every proposed declaration is in Lean. The Lean formalization is in
\code{LineGraphs/}, and \code{dependency-diagram.md}, at the top of the
repository, maps the numbered results of [CGSS] to its declarations.
The spectral certificates, positive-vector bound and similarity lemmas
already have proofs in \code{LineGraphs/Spectral/}; that directory also
defines switching equivalence. The Boolean soundness lemmas and their
\code{Matrix} aliases are in \code{LineGraphs/Computations/}. Each
instantiated C13, C15 or C17 certificate requires a kernel-checked proof
of its true-checker hypothesis; for C13 and C15 these are in
\code{LineGraphs/Computations/}. No structured proof obligation in this
document is deferred to an unproved project declaration.

\begin{thebibliography}{CGSS}
\bibitem[Err]{errata} \emph{Errata to CGSS}, \code{cgss/errata.md}, Errors 9 and 14.
\bibitem[CGSS]{CGSS} P. J. Cameron, J. M. Goethals, J. J. Seidel and E. E. Shult,
Line graphs, root systems, and elliptic geometry, \emph{J. Algebra}
\textbf{43} (1976), 305--327.
\bibitem[L95]{L95} L. Lamport, How to write a proof,
\emph{Amer. Math. Monthly} \textbf{102} (1995), 600--608.
\end{thebibliography}
\end{document}
